% Môn: Vật lý
% Lớp: 12
% Chương: Chương III. Từ trường
% Bài: L12C3 Bài 20. Bài tập về từ trường
% Số câu: 324
% App đọc file này trực tiếp — sửa rồi Commit trên GitHub, bấm Đồng bộ GitHub .tex

% L12C3 Bài 20. Bài tập về từ trường

% ===== Câu 1 =====
% ID: L12C3B20-01-DS
% Mức: TH
\dangbt{Bài toán vận dụng cao kết hợp nhiều nội dung của Chương III}
\begin{ex}
% ID: L12C3B20-01-DS
% Mức: TH
Xét bài toán thuộc dạng bài toán vận dụng cao kết hợp nhiều nội dung của chương iii. Hãy xác định tính đúng, sai của bốn phát biểu sau.
\choiceTF
{Có thể cộng đại số mọi vectơ mà không xét phương chiều.}
{Có thể dùng trực tiếp số đo góc theo độ trong mọi biểu thức máy tính mà không kiểm tra chế độ.}
{\True Định luật Lenz quyết định chiều, còn định luật Faraday xác định độ lớn suất điện động cảm ứng.}
{\True Kết quả phải được kiểm tra về đơn vị và ý nghĩa vật lí.}
\loigiai{
\begin{itemchoice}
\item (Sai) Có thể cộng đại số mọi vectơ mà không xét phương chiều. (Vì): Có thể cộng đại số mọi vectơ mà không xét phương chiều. b) Sai: Có thể dùng trực tiếp số đo góc theo độ trong mọi biểu thức máy tính mà không kiểm tra chế độ. c) Đúng: Định luật Lenz quyết định chiều, còn định luật Faraday xác định độ lớn suất điện động cảm ứng. d) Đúng: Kết quả phải được kiểm tra về đơn vị và ý nghĩa vật lí.. Bài toán tổng hợp phải tách các giai đoạn, xác định đúng phương chiều, chọn định luật phù hợp, đổi đơn vị thống nhất và kiểm tra ý nghĩa vật lí của kết quả
\item (Sai) Có thể dùng trực tiếp số đo góc theo độ trong mọi biểu thức máy tính mà không kiểm tra chế độ.
\item (Đúng) Định luật Lenz quyết định chiều, còn định luật Faraday xác định độ lớn suất điện động cảm ứng.
\item (Đúng) Kết quả phải được kiểm tra về đơn vị và ý nghĩa vật lí.
\end{itemchoice}
}
\end{ex}

% ===== Câu 2 =====
% ID: L12C3B20-01-TL
% Mức: TH
\begin{bt}
% ID: L12C3B20-01-TL
% Mức: TH
Trình bày kiến thức cốt lõi và các bước giải bài thuộc dạng bài toán vận dụng cao kết hợp nhiều nội dung của chương iii.
\loigiai{
Bài toán tổng hợp phải tách các giai đoạn, xác định đúng phương chiều, chọn định luật phù hợp, đổi đơn vị thống nhất và kiểm tra ý nghĩa vật lí của kết quả. Các bước: tóm tắt dữ kiện; xác định phương, chiều và điều kiện; chọn hệ thức; đổi đơn vị; tính toán; kiểm tra kết quả.
}
\end{bt}

% ===== Câu 3 =====
% ID: L12C3B20-01-TLN
% Mức: TH
\begin{ex}
% ID: L12C3B20-01-TLN
% Mức: TH
Trong bốn phát biểu sau về bài toán vận dụng cao kết hợp nhiều nội dung của chương iii, có bao nhiêu phát biểu đúng? Chỉ ghi một số nguyên. (1) Khi giải bài tổng hợp cần xác định rõ chiều của $B$, dòng điện, lực từ và dòng điện cảm ứng. (2) Có thể cộng đại số mọi vectơ mà không xét phương chiều. (3) Các đại lượng phải được đổi về cùng hệ đơn vị trước khi thay số. (4) Có thể dùng trực tiếp số đo góc theo độ trong mọi biểu thức máy tính mà không kiểm tra chế độ.
\shortans{2}
\loigiai{
Đối chiếu từng phát biểu với kiến thức: Bài toán tổng hợp phải tách các giai đoạn, xác định đúng phương chiều, chọn định luật phù hợp, đổi đơn vị thống nhất và kiểm tra ý nghĩa vật lí của kết quả. Có 2 phát biểu đúng.
}
\end{ex}

% ===== Câu 4 =====
% ID: L12C3B20-01-TN
% Mức: NB
\begin{ex}
% ID: L12C3B20-01-TN
% Mức: NB
Phát biểu nào sau đây đúng về bài toán vận dụng cao kết hợp nhiều nội dung của chương iii?
\choice
{\True Khi giải bài tổng hợp cần xác định rõ chiều của $B$, dòng điện, lực từ và dòng điện cảm ứng.}
{Có thể cộng đại số mọi vectơ mà không xét phương chiều.}
{Có thể dùng trực tiếp số đo góc theo độ trong mọi biểu thức máy tính mà không kiểm tra chế độ.}
{Hiệu suất của mọi thiết bị điện luôn bằng 100 phần trăm.}
\loigiai{
Khi giải bài tổng hợp cần xác định rõ chiều của $B$, dòng điện, lực từ và dòng điện cảm ứng. Do đó chọn A. Bài toán tổng hợp phải tách các giai đoạn, xác định đúng phương chiều, chọn định luật phù hợp, đổi đơn vị thống nhất và kiểm tra ý nghĩa vật lí của kết quả.
}
\end{ex}

% ===== Câu 5 =====
% ID: L12C3B20-02-DS
% Mức: TH
\begin{ex}
% ID: L12C3B20-02-DS
% Mức: TH
Xét bài toán thuộc dạng bài toán vận dụng cao kết hợp nhiều nội dung của chương iii. Hãy xác định tính đúng, sai của bốn phát biểu sau.
\choiceTF
{\True Các đại lượng phải được đổi về cùng hệ đơn vị trước khi thay số.}
{Hiệu suất của mọi thiết bị điện luôn bằng 100 phần trăm.}
{\True Kết quả phải được kiểm tra về đơn vị và ý nghĩa vật lí.}
{Có thể cộng đại số mọi vectơ mà không xét phương chiều.}
\loigiai{
\begin{itemchoice}
\item (Đúng) Các đại lượng phải được đổi về cùng hệ đơn vị trước khi thay số. (Vì): Các đại lượng phải được đổi về cùng hệ đơn vị trước khi thay số. b) Sai: Hiệu suất của mọi thiết bị điện luôn bằng 100 phần trăm. c) Đúng: Kết quả phải được kiểm tra về đơn vị và ý nghĩa vật lí. d) Sai: Có thể cộng đại số mọi vectơ mà không xét phương chiều.. Bài toán tổng hợp phải tách các giai đoạn, xác định đúng phương chiều, chọn định luật phù hợp, đổi đơn vị thống nhất và kiểm tra ý nghĩa vật lí của kết quả
\item (Sai) Hiệu suất của mọi thiết bị điện luôn bằng 100 phần trăm.
\item (Đúng) Kết quả phải được kiểm tra về đơn vị và ý nghĩa vật lí.
\item (Sai) Có thể cộng đại số mọi vectơ mà không xét phương chiều.
\end{itemchoice}
}
\end{ex}

% ===== Câu 6 =====
% ID: L12C3B20-02-TL
% Mức: TH
\begin{bt}
% ID: L12C3B20-02-TL
% Mức: TH
Phân tích các đại lượng, phương chiều và điều kiện áp dụng trong dạng bài toán vận dụng cao kết hợp nhiều nội dung của chương iii.
\loigiai{
Bài toán tổng hợp phải tách các giai đoạn, xác định đúng phương chiều, chọn định luật phù hợp, đổi đơn vị thống nhất và kiểm tra ý nghĩa vật lí của kết quả. Cần xác định rõ đại lượng vô hướng hay vectơ, quy ước chiều và điều kiện để công thức được áp dụng.
}
\end{bt}

% ===== Câu 7 =====
% ID: L12C3B20-02-TLN
% Mức: TH
\begin{ex}
% ID: L12C3B20-02-TLN
% Mức: TH
Trong bốn phát biểu sau về bài toán vận dụng cao kết hợp nhiều nội dung của chương iii, có bao nhiêu phát biểu đúng? Chỉ ghi một số nguyên. (1) Định luật Lenz quyết định chiều, còn định luật Faraday xác định độ lớn suất điện động cảm ứng. (2) Kết quả phải được kiểm tra về đơn vị và ý nghĩa vật lí. (3) Hiệu suất của mọi thiết bị điện luôn bằng 100 phần trăm. (4) Trong bài tổng hợp có thể bỏ qua điều kiện áp dụng của từng định luật.
\shortans{2}
\loigiai{
Đối chiếu từng phát biểu với kiến thức: Bài toán tổng hợp phải tách các giai đoạn, xác định đúng phương chiều, chọn định luật phù hợp, đổi đơn vị thống nhất và kiểm tra ý nghĩa vật lí của kết quả. Có 2 phát biểu đúng.
}
\end{ex}

% ===== Câu 8 =====
% ID: L12C3B20-02-TN
% Mức: NB
\begin{ex}
% ID: L12C3B20-02-TN
% Mức: NB
Trong Chương III, kết luận nào đúng khi giải bài thuộc dạng bài toán vận dụng cao kết hợp nhiều nội dung của chương iii?
\choice
{Có thể dùng trực tiếp số đo góc theo độ trong mọi biểu thức máy tính mà không kiểm tra chế độ.}
{\True Các đại lượng phải được đổi về cùng hệ đơn vị trước khi thay số.}
{Hiệu suất của mọi thiết bị điện luôn bằng 100 phần trăm.}
{Trong bài tổng hợp có thể bỏ qua điều kiện áp dụng của từng định luật.}
\loigiai{
Các đại lượng phải được đổi về cùng hệ đơn vị trước khi thay số. Do đó chọn B. Bài toán tổng hợp phải tách các giai đoạn, xác định đúng phương chiều, chọn định luật phù hợp, đổi đơn vị thống nhất và kiểm tra ý nghĩa vật lí của kết quả.
}
\end{ex}

% ===== Câu 9 =====
% ID: L12C3B20-03-DS
% Mức: VD
\begin{ex}
% ID: L12C3B20-03-DS
% Mức: VD
Xét bài toán thuộc dạng bài toán vận dụng cao kết hợp nhiều nội dung của chương iii. Hãy xác định tính đúng, sai của bốn phát biểu sau.
\choiceTF
{Hiệu suất của mọi thiết bị điện luôn bằng 100 phần trăm.}
{\True Kết quả phải được kiểm tra về đơn vị và ý nghĩa vật lí.}
{Có thể cộng đại số mọi vectơ mà không xét phương chiều.}
{\True Các đại lượng phải được đổi về cùng hệ đơn vị trước khi thay số.}
\loigiai{
\begin{itemchoice}
\item (Sai) Hiệu suất của mọi thiết bị điện luôn bằng 100 phần trăm. (Vì): Hiệu suất của mọi thiết bị điện luôn bằng 100 phần trăm. b) Đúng: Kết quả phải được kiểm tra về đơn vị và ý nghĩa vật lí. c) Sai: Có thể cộng đại số mọi vectơ mà không xét phương chiều. d) Đúng: Các đại lượng phải được đổi về cùng hệ đơn vị trước khi thay số.. Bài toán tổng hợp phải tách các giai đoạn, xác định đúng phương chiều, chọn định luật phù hợp, đổi đơn vị thống nhất và kiểm tra ý nghĩa vật lí của kết quả
\item (Đúng) Kết quả phải được kiểm tra về đơn vị và ý nghĩa vật lí.
\item (Sai) Có thể cộng đại số mọi vectơ mà không xét phương chiều.
\item (Đúng) Các đại lượng phải được đổi về cùng hệ đơn vị trước khi thay số.
\end{itemchoice}
}
\end{ex}

% ===== Câu 10 =====
% ID: L12C3B20-03-TL
% Mức: VD
\begin{bt}
% ID: L12C3B20-03-TL
% Mức: VD
Nêu một sai lầm thường gặp khi giải bài toán vận dụng cao kết hợp nhiều nội dung của chương iii và cách khắc phục.
\loigiai{
Bài toán tổng hợp phải tách các giai đoạn, xác định đúng phương chiều, chọn định luật phù hợp, đổi đơn vị thống nhất và kiểm tra ý nghĩa vật lí của kết quả. Sai lầm thường gặp là bỏ qua phương chiều, góc hoặc đơn vị. Khắc phục bằng hình vẽ, quy ước dấu và đổi đơn vị trước khi tính.
}
\end{bt}

% ===== Câu 11 =====
% ID: L12C3B20-03-TLN
% Mức: VD
\begin{ex}
% ID: L12C3B20-03-TLN
% Mức: VD
Trong bốn phát biểu sau về bài toán vận dụng cao kết hợp nhiều nội dung của chương iii, có bao nhiêu phát biểu đúng? Chỉ ghi một số nguyên. (1) Khi giải bài tổng hợp cần xác định rõ chiều của $B$, dòng điện, lực từ và dòng điện cảm ứng. (2) Định luật Lenz quyết định chiều, còn định luật Faraday xác định độ lớn suất điện động cảm ứng. (3) Có thể dùng trực tiếp số đo góc theo độ trong mọi biểu thức máy tính mà không kiểm tra chế độ. (4) Trong bài tổng hợp có thể bỏ qua điều kiện áp dụng của từng định luật.
\shortans{2}
\loigiai{
Đối chiếu từng phát biểu với kiến thức: Bài toán tổng hợp phải tách các giai đoạn, xác định đúng phương chiều, chọn định luật phù hợp, đổi đơn vị thống nhất và kiểm tra ý nghĩa vật lí của kết quả. Có 2 phát biểu đúng.
}
\end{ex}

% ===== Câu 12 =====
% ID: L12C3B20-03-TN
% Mức: NB
\begin{ex}
% ID: L12C3B20-03-TN
% Mức: NB
Chọn nhận định phù hợp với kiến thức Vật lí về bài toán vận dụng cao kết hợp nhiều nội dung của chương iii.
\choice
{Hiệu suất của mọi thiết bị điện luôn bằng 100 phần trăm.}
{Trong bài tổng hợp có thể bỏ qua điều kiện áp dụng của từng định luật.}
{\True Định luật Lenz quyết định chiều, còn định luật Faraday xác định độ lớn suất điện động cảm ứng.}
{Có thể cộng đại số mọi vectơ mà không xét phương chiều.}
\loigiai{
Định luật Lenz quyết định chiều, còn định luật Faraday xác định độ lớn suất điện động cảm ứng. Do đó chọn C. Bài toán tổng hợp phải tách các giai đoạn, xác định đúng phương chiều, chọn định luật phù hợp, đổi đơn vị thống nhất và kiểm tra ý nghĩa vật lí của kết quả.
}
\end{ex}

% ===== Câu 13 =====
% ID: L12C3B20-04-DS
% Mức: VD
\begin{ex}
% ID: L12C3B20-04-DS
% Mức: VD
Xét bài toán thuộc dạng bài toán vận dụng cao kết hợp nhiều nội dung của chương iii. Hãy xác định tính đúng, sai của bốn phát biểu sau.
\choiceTF
{\True Kết quả phải được kiểm tra về đơn vị và ý nghĩa vật lí.}
{\True Khi giải bài tổng hợp cần xác định rõ chiều của $B$, dòng điện, lực từ và dòng điện cảm ứng.}
{Có thể dùng trực tiếp số đo góc theo độ trong mọi biểu thức máy tính mà không kiểm tra chế độ.}
{Hiệu suất của mọi thiết bị điện luôn bằng 100 phần trăm.}
\loigiai{
\begin{itemchoice}
\item (Đúng) Kết quả phải được kiểm tra về đơn vị và ý nghĩa vật lí. (Vì): Kết quả phải được kiểm tra về đơn vị và ý nghĩa vật lí. b) Đúng: Khi giải bài tổng hợp cần xác định rõ chiều của $B$, dòng điện, lực từ và dòng điện cảm ứng. c) Sai: Có thể dùng trực tiếp số đo góc theo độ trong mọi biểu thức máy tính mà không kiểm tra chế độ. d) Sai: Hiệu suất của mọi thiết bị điện luôn bằng 100 phần trăm.. Bài toán tổng hợp phải tách các giai đoạn, xác định đúng phương chiều, chọn định luật phù hợp, đổi đơn vị thống nhất và kiểm tra ý nghĩa vật lí của kết quả
\item (Đúng) Khi giải bài tổng hợp cần xác định rõ chiều của $B$, dòng điện, lực từ và dòng điện cảm ứng.
\item (Sai) Có thể dùng trực tiếp số đo góc theo độ trong mọi biểu thức máy tính mà không kiểm tra chế độ.
\item (Sai) Hiệu suất của mọi thiết bị điện luôn bằng 100 phần trăm.
\end{itemchoice}
}
\end{ex}

% ===== Câu 14 =====
% ID: L12C3B20-04-TL
% Mức: VD
\begin{bt}
% ID: L12C3B20-04-TL
% Mức: VD
Đề xuất quy trình kiểm tra kết quả của một bài toán bài toán vận dụng cao kết hợp nhiều nội dung của chương iii.
\loigiai{
Bài toán tổng hợp phải tách các giai đoạn, xác định đúng phương chiều, chọn định luật phù hợp, đổi đơn vị thống nhất và kiểm tra ý nghĩa vật lí của kết quả. Kiểm tra thứ nguyên, dấu, giới hạn đặc biệt và sự phù hợp với ý nghĩa vật lí.
}
\end{bt}

% ===== Câu 15 =====
% ID: L12C3B20-04-TLN
% Mức: VD
\begin{ex}
% ID: L12C3B20-04-TLN
% Mức: VD
Trong bốn phát biểu sau về bài toán vận dụng cao kết hợp nhiều nội dung của chương iii, có bao nhiêu phát biểu đúng? Chỉ ghi một số nguyên. (1) Khi giải bài tổng hợp cần xác định rõ chiều của $B$, dòng điện, lực từ và dòng điện cảm ứng. (2) Có thể cộng đại số mọi vectơ mà không xét phương chiều. (3) Các đại lượng phải được đổi về cùng hệ đơn vị trước khi thay số. (4) Có thể dùng trực tiếp số đo góc theo độ trong mọi biểu thức máy tính mà không kiểm tra chế độ.
\shortans{2}
\loigiai{
Đối chiếu từng phát biểu với kiến thức: Bài toán tổng hợp phải tách các giai đoạn, xác định đúng phương chiều, chọn định luật phù hợp, đổi đơn vị thống nhất và kiểm tra ý nghĩa vật lí của kết quả. Có 2 phát biểu đúng.
}
\end{ex}

% ===== Câu 16 =====
% ID: L12C3B20-04-TN
% Mức: TH
\begin{ex}
% ID: L12C3B20-04-TN
% Mức: TH
Khi phân tích bài toán vận dụng cao kết hợp nhiều nội dung của chương iii, phát biểu nào đúng?
\choice
{Trong bài tổng hợp có thể bỏ qua điều kiện áp dụng của từng định luật.}
{Có thể cộng đại số mọi vectơ mà không xét phương chiều.}
{Có thể dùng trực tiếp số đo góc theo độ trong mọi biểu thức máy tính mà không kiểm tra chế độ.}
{\True Kết quả phải được kiểm tra về đơn vị và ý nghĩa vật lí.}
\loigiai{
Kết quả phải được kiểm tra về đơn vị và ý nghĩa vật lí. Do đó chọn D. Bài toán tổng hợp phải tách các giai đoạn, xác định đúng phương chiều, chọn định luật phù hợp, đổi đơn vị thống nhất và kiểm tra ý nghĩa vật lí của kết quả.
}
\end{ex}

% ===== Câu 17 =====
% ID: L12C3B20-05-DS
% Mức: VD
\begin{ex}
% ID: L12C3B20-05-DS
% Mức: VD
Xét bài toán thuộc dạng bài toán vận dụng cao kết hợp nhiều nội dung của chương iii. Hãy xác định tính đúng, sai của bốn phát biểu sau.
\choiceTF
{Có thể cộng đại số mọi vectơ mà không xét phương chiều.}
{\True Các đại lượng phải được đổi về cùng hệ đơn vị trước khi thay số.}
{\True Định luật Lenz quyết định chiều, còn định luật Faraday xác định độ lớn suất điện động cảm ứng.}
{Trong bài tổng hợp có thể bỏ qua điều kiện áp dụng của từng định luật.}
\loigiai{
\begin{itemchoice}
\item (Sai) Có thể cộng đại số mọi vectơ mà không xét phương chiều. (Vì): Có thể cộng đại số mọi vectơ mà không xét phương chiều. b) Đúng: Các đại lượng phải được đổi về cùng hệ đơn vị trước khi thay số. c) Đúng: Định luật Lenz quyết định chiều, còn định luật Faraday xác định độ lớn suất điện động cảm ứng. d) Sai: Trong bài tổng hợp có thể bỏ qua điều kiện áp dụng của từng định luật.. Bài toán tổng hợp phải tách các giai đoạn, xác định đúng phương chiều, chọn định luật phù hợp, đổi đơn vị thống nhất và kiểm tra ý nghĩa vật lí của kết quả
\item (Đúng) Các đại lượng phải được đổi về cùng hệ đơn vị trước khi thay số.
\item (Đúng) Định luật Lenz quyết định chiều, còn định luật Faraday xác định độ lớn suất điện động cảm ứng.
\item (Sai) Trong bài tổng hợp có thể bỏ qua điều kiện áp dụng của từng định luật.
\end{itemchoice}
}
\end{ex}

% ===== Câu 18 =====
% ID: L12C3B20-05-TL
% Mức: VD
\begin{bt}
% ID: L12C3B20-05-TL
% Mức: VD
Xây dựng một tình huống thực tiễn thuộc dạng bài toán vận dụng cao kết hợp nhiều nội dung của chương iii và chỉ ra định luật cần dùng.
\loigiai{
Bài toán tổng hợp phải tách các giai đoạn, xác định đúng phương chiều, chọn định luật phù hợp, đổi đơn vị thống nhất và kiểm tra ý nghĩa vật lí của kết quả. Có thể chọn thiết bị hoặc hiện tượng thực tế tương ứng, nêu dữ kiện và chỉ rõ định luật dùng để giải thích hoặc tính toán.
}
\end{bt}

% ===== Câu 19 =====
% ID: L12C3B20-05-TLN
% Mức: VD
\begin{ex}
% ID: L12C3B20-05-TLN
% Mức: VD
Trong bốn phát biểu sau về bài toán vận dụng cao kết hợp nhiều nội dung của chương iii, có bao nhiêu phát biểu đúng? Chỉ ghi một số nguyên. (1) Định luật Lenz quyết định chiều, còn định luật Faraday xác định độ lớn suất điện động cảm ứng. (2) Kết quả phải được kiểm tra về đơn vị và ý nghĩa vật lí. (3) Hiệu suất của mọi thiết bị điện luôn bằng 100 phần trăm. (4) Trong bài tổng hợp có thể bỏ qua điều kiện áp dụng của từng định luật.
\shortans{2}
\loigiai{
Đối chiếu từng phát biểu với kiến thức: Bài toán tổng hợp phải tách các giai đoạn, xác định đúng phương chiều, chọn định luật phù hợp, đổi đơn vị thống nhất và kiểm tra ý nghĩa vật lí của kết quả. Có 2 phát biểu đúng.
}
\end{ex}

% ===== Câu 20 =====
% ID: L12C3B20-05-TN
% Mức: TH
\begin{ex}
% ID: L12C3B20-05-TN
% Mức: TH
Cơ sở Vật lí nào sau đây cần được sử dụng trong dạng bài toán vận dụng cao kết hợp nhiều nội dung của chương iii?
\choice
{\True Khi giải bài tổng hợp cần xác định rõ chiều của $B$, dòng điện, lực từ và dòng điện cảm ứng.}
{Có thể cộng đại số mọi vectơ mà không xét phương chiều.}
{Có thể dùng trực tiếp số đo góc theo độ trong mọi biểu thức máy tính mà không kiểm tra chế độ.}
{Hiệu suất của mọi thiết bị điện luôn bằng 100 phần trăm.}
\loigiai{
Khi giải bài tổng hợp cần xác định rõ chiều của $B$, dòng điện, lực từ và dòng điện cảm ứng. Do đó chọn A. Bài toán tổng hợp phải tách các giai đoạn, xác định đúng phương chiều, chọn định luật phù hợp, đổi đơn vị thống nhất và kiểm tra ý nghĩa vật lí của kết quả.
}
\end{ex}

% ===== Câu 21 =====
% ID: L12C3B20-06-DS
% Mức: TH
\dangbt{Tổng hợp lực từ tác dụng lên dây dẫn mang dòng điện}
\begin{ex}
% ID: L12C3B20-06-DS
% Mức: TH
Xét bài toán thuộc dạng tổng hợp lực từ tác dụng lên dây dẫn mang dòng điện. Hãy xác định tính đúng, sai của bốn phát biểu sau.
\choiceTF
{Lực từ luôn cùng chiều với dòng điện.}
{Khi dây dẫn vuông góc với từ trường thì lực từ bằng 0.}
{\True Khi dây dẫn song song với đường sức từ thì lực từ bằng 0.}
{\True Đổi chiều dòng điện thì chiều lực từ đổi ngược.}
\loigiai{
\begin{itemchoice}
\item (Sai) Lực từ luôn cùng chiều với dòng điện. (Vì): lpha); chiều lực xác định bằng quy tắc bàn tay trái
\item (Sai) Khi dây dẫn vuông góc với từ trường thì lực từ bằng 0.
\item (Đúng) Khi dây dẫn song song với đường sức từ thì lực từ bằng 0.
\item (Đúng) Đổi chiều dòng điện thì chiều lực từ đổi ngược.
\end{itemchoice}
}
\end{ex}

% ===== Câu 22 =====
% ID: L12C3B20-06-TL
% Mức: VDC
\dangbt{Bài toán vận dụng cao kết hợp nhiều nội dung của Chương III}
\begin{bt}
% ID: L12C3B20-06-TL
% Mức: VDC
So sánh một nhận định đúng và một nhận định sai trong dạng bài toán vận dụng cao kết hợp nhiều nội dung của chương iii; giải thích căn cứ Vật lí.
\loigiai{
Nhận định đúng: Khi giải bài tổng hợp cần xác định rõ chiều của $B$, dòng điện, lực từ và dòng điện cảm ứng. Nhận định sai: Có thể cộng đại số mọi vectơ mà không xét phương chiều. Căn cứ: Bài toán tổng hợp phải tách các giai đoạn, xác định đúng phương chiều, chọn định luật phù hợp, đổi đơn vị thống nhất và kiểm tra ý nghĩa vật lí của kết quả.
}
\end{bt}

% ===== Câu 23 =====
% ID: L12C3B20-06-TLN
% Mức: VDC
\begin{ex}
% ID: L12C3B20-06-TLN
% Mức: VDC
Trong bốn phát biểu sau về bài toán vận dụng cao kết hợp nhiều nội dung của chương iii, có bao nhiêu phát biểu đúng? Chỉ ghi một số nguyên. (1) Khi giải bài tổng hợp cần xác định rõ chiều của $B$, dòng điện, lực từ và dòng điện cảm ứng. (2) Định luật Lenz quyết định chiều, còn định luật Faraday xác định độ lớn suất điện động cảm ứng. (3) Có thể dùng trực tiếp số đo góc theo độ trong mọi biểu thức máy tính mà không kiểm tra chế độ. (4) Trong bài tổng hợp có thể bỏ qua điều kiện áp dụng của từng định luật.
\shortans{2}
\loigiai{
Đối chiếu từng phát biểu với kiến thức: Bài toán tổng hợp phải tách các giai đoạn, xác định đúng phương chiều, chọn định luật phù hợp, đổi đơn vị thống nhất và kiểm tra ý nghĩa vật lí của kết quả. Có 2 phát biểu đúng.
}
\end{ex}

% ===== Câu 24 =====
% ID: L12C3B20-06-TN
% Mức: TH
\begin{ex}
% ID: L12C3B20-06-TN
% Mức: TH
Phát biểu nào sau đây đúng về bài toán vận dụng cao kết hợp nhiều nội dung của chương iii?
\choice
{Có thể dùng trực tiếp số đo góc theo độ trong mọi biểu thức máy tính mà không kiểm tra chế độ.}
{\True Các đại lượng phải được đổi về cùng hệ đơn vị trước khi thay số.}
{Hiệu suất của mọi thiết bị điện luôn bằng 100 phần trăm.}
{Trong bài tổng hợp có thể bỏ qua điều kiện áp dụng của từng định luật.}
\loigiai{
Các đại lượng phải được đổi về cùng hệ đơn vị trước khi thay số. Do đó chọn B. Bài toán tổng hợp phải tách các giai đoạn, xác định đúng phương chiều, chọn định luật phù hợp, đổi đơn vị thống nhất và kiểm tra ý nghĩa vật lí của kết quả.
}
\end{ex}

% ===== Câu 25 =====
% ID: L12C3B20-07-DS
% Mức: TH
\dangbt{Tổng hợp lực từ tác dụng lên dây dẫn mang dòng điện}
\begin{ex}
% ID: L12C3B20-07-DS
% Mức: TH
Xét bài toán thuộc dạng tổng hợp lực từ tác dụng lên dây dẫn mang dòng điện. Hãy xác định tính đúng, sai của bốn phát biểu sau.
\choiceTF
{\True Lực từ vuông góc với cả dòng điện và vectơ cảm ứng từ.}
{Độ lớn lực từ không phụ thuộc chiều dài đoạn dây.}
{\True Đổi chiều dòng điện thì chiều lực từ đổi ngược.}
{Lực từ luôn cùng chiều với dòng điện.}
\loigiai{
\begin{itemchoice}
\item (Đúng) Lực từ vuông góc với cả dòng điện và vectơ cảm ứng từ. (Vì): lpha); chiều lực xác định bằng quy tắc bàn tay trái
\item (Sai) Độ lớn lực từ không phụ thuộc chiều dài đoạn dây.
\item (Đúng) Đổi chiều dòng điện thì chiều lực từ đổi ngược.
\item (Sai) Lực từ luôn cùng chiều với dòng điện.
\end{itemchoice}
}
\end{ex}

% ===== Câu 26 =====
% ID: L12C3B20-07-TL
% Mức: TH
\begin{bt}
% ID: L12C3B20-07-TL
% Mức: TH
Trình bày kiến thức cốt lõi và các bước giải bài thuộc dạng tổng hợp lực từ tác dụng lên dây dẫn mang dòng điện.
\loigiai{
Với đoạn dây thẳng dài l mang dòng $I$ trong từ trường đều $B$, $F$ = BIl sin(alpha); chiều lực xác định bằng quy tắc bàn tay trái. Các bước: tóm tắt dữ kiện; xác định phương, chiều và điều kiện; chọn hệ thức; đổi đơn vị; tính toán; kiểm tra kết quả.
}
\end{bt}

% ===== Câu 27 =====
% ID: L12C3B20-07-TLN
% Mức: TH
\begin{ex}
% ID: L12C3B20-07-TLN
% Mức: TH
Trong bốn phát biểu sau về tổng hợp lực từ tác dụng lên dây dẫn mang dòng điện, có bao nhiêu phát biểu đúng? Chỉ ghi một số nguyên. (1) Độ lớn lực từ tác dụng lên đoạn dây thẳng là $F$ = BIl sin(alpha). (2) Lực từ luôn cùng chiều với dòng điện. (3) Lực từ vuông góc với cả dòng điện và vectơ cảm ứng từ. (4) Khi dây dẫn vuông góc với từ trường thì lực từ bằng 0.
\shortans{2}
\loigiai{
Đối chiếu từng phát biểu với kiến thức: Với đoạn dây thẳng dài l mang dòng $I$ trong từ trường đều $B$, $F$ = BIl sin(alpha); chiều lực xác định bằng quy tắc bàn tay trái. Có 2 phát biểu đúng.
}
\end{ex}

% ===== Câu 28 =====
% ID: L12C3B20-07-TN
% Mức: TH
\dangbt{Bài toán vận dụng cao kết hợp nhiều nội dung của Chương III}
\begin{ex}
% ID: L12C3B20-07-TN
% Mức: TH
Trong Chương III, kết luận nào đúng khi giải bài thuộc dạng bài toán vận dụng cao kết hợp nhiều nội dung của chương iii?
\choice
{Hiệu suất của mọi thiết bị điện luôn bằng 100 phần trăm.}
{Trong bài tổng hợp có thể bỏ qua điều kiện áp dụng của từng định luật.}
{\True Định luật Lenz quyết định chiều, còn định luật Faraday xác định độ lớn suất điện động cảm ứng.}
{Có thể cộng đại số mọi vectơ mà không xét phương chiều.}
\loigiai{
Định luật Lenz quyết định chiều, còn định luật Faraday xác định độ lớn suất điện động cảm ứng. Do đó chọn C. Bài toán tổng hợp phải tách các giai đoạn, xác định đúng phương chiều, chọn định luật phù hợp, đổi đơn vị thống nhất và kiểm tra ý nghĩa vật lí của kết quả.
}
\end{ex}

% ===== Câu 29 =====
% ID: L12C3B20-08-DS
% Mức: VD
\dangbt{Tổng hợp lực từ tác dụng lên dây dẫn mang dòng điện}
\begin{ex}
% ID: L12C3B20-08-DS
% Mức: VD
Xét bài toán thuộc dạng tổng hợp lực từ tác dụng lên dây dẫn mang dòng điện. Hãy xác định tính đúng, sai của bốn phát biểu sau.
\choiceTF
{Độ lớn lực từ không phụ thuộc chiều dài đoạn dây.}
{\True Đổi chiều dòng điện thì chiều lực từ đổi ngược.}
{Lực từ luôn cùng chiều với dòng điện.}
{\True Lực từ vuông góc với cả dòng điện và vectơ cảm ứng từ.}
\loigiai{
\begin{itemchoice}
\item (Sai) Độ lớn lực từ không phụ thuộc chiều dài đoạn dây. (Vì): lpha); chiều lực xác định bằng quy tắc bàn tay trái
\item (Đúng) Đổi chiều dòng điện thì chiều lực từ đổi ngược.
\item (Sai) Lực từ luôn cùng chiều với dòng điện.
\item (Đúng) Lực từ vuông góc với cả dòng điện và vectơ cảm ứng từ.
\end{itemchoice}
}
\end{ex}

% ===== Câu 30 =====
% ID: L12C3B20-08-TL
% Mức: TH
\begin{bt}
% ID: L12C3B20-08-TL
% Mức: TH
Phân tích các đại lượng, phương chiều và điều kiện áp dụng trong dạng tổng hợp lực từ tác dụng lên dây dẫn mang dòng điện.
\loigiai{
Với đoạn dây thẳng dài l mang dòng $I$ trong từ trường đều $B$, $F$ = BIl sin(alpha); chiều lực xác định bằng quy tắc bàn tay trái. Cần xác định rõ đại lượng vô hướng hay vectơ, quy ước chiều và điều kiện để công thức được áp dụng.
}
\end{bt}

% ===== Câu 31 =====
% ID: L12C3B20-08-TLN
% Mức: TH
\begin{ex}
% ID: L12C3B20-08-TLN
% Mức: TH
Trong bốn phát biểu sau về tổng hợp lực từ tác dụng lên dây dẫn mang dòng điện, có bao nhiêu phát biểu đúng? Chỉ ghi một số nguyên. (1) Khi dây dẫn song song với đường sức từ thì lực từ bằng 0. (2) Đổi chiều dòng điện thì chiều lực từ đổi ngược. (3) Độ lớn lực từ không phụ thuộc chiều dài đoạn dây. (4) Đổi chiều cảm ứng từ nhưng giữ dòng điện thì chiều lực từ không đổi.
\shortans{2}
\loigiai{
Đối chiếu từng phát biểu với kiến thức: Với đoạn dây thẳng dài l mang dòng $I$ trong từ trường đều $B$, $F$ = BIl sin(alpha); chiều lực xác định bằng quy tắc bàn tay trái. Có 2 phát biểu đúng.
}
\end{ex}

% ===== Câu 32 =====
% ID: L12C3B20-08-TN
% Mức: VD
\dangbt{Bài toán vận dụng cao kết hợp nhiều nội dung của Chương III}
\begin{ex}
% ID: L12C3B20-08-TN
% Mức: VD
Chọn nhận định phù hợp với kiến thức Vật lí về bài toán vận dụng cao kết hợp nhiều nội dung của chương iii.
\choice
{Trong bài tổng hợp có thể bỏ qua điều kiện áp dụng của từng định luật.}
{Có thể cộng đại số mọi vectơ mà không xét phương chiều.}
{Có thể dùng trực tiếp số đo góc theo độ trong mọi biểu thức máy tính mà không kiểm tra chế độ.}
{\True Kết quả phải được kiểm tra về đơn vị và ý nghĩa vật lí.}
\loigiai{
Kết quả phải được kiểm tra về đơn vị và ý nghĩa vật lí. Do đó chọn D. Bài toán tổng hợp phải tách các giai đoạn, xác định đúng phương chiều, chọn định luật phù hợp, đổi đơn vị thống nhất và kiểm tra ý nghĩa vật lí của kết quả.
}
\end{ex}

% ===== Câu 33 =====
% ID: L12C3B20-09-DS
% Mức: VD
\dangbt{Tổng hợp lực từ tác dụng lên dây dẫn mang dòng điện}
\begin{ex}
% ID: L12C3B20-09-DS
% Mức: VD
Xét bài toán thuộc dạng tổng hợp lực từ tác dụng lên dây dẫn mang dòng điện. Hãy xác định tính đúng, sai của bốn phát biểu sau.
\choiceTF
{\True Đổi chiều dòng điện thì chiều lực từ đổi ngược.}
{\True Độ lớn lực từ tác dụng lên đoạn dây thẳng là $F$ = BIl sin(alpha).}
{Khi dây dẫn vuông góc với từ trường thì lực từ bằng 0.}
{Độ lớn lực từ không phụ thuộc chiều dài đoạn dây.}
\loigiai{
\begin{itemchoice}
\item (Đúng) Đổi chiều dòng điện thì chiều lực từ đổi ngược. (Vì): lpha); chiều lực xác định bằng quy tắc bàn tay trái
\item (Đúng) Độ lớn lực từ tác dụng lên đoạn dây thẳng là $F$ = BIl sin(alpha).
\item (Sai) Khi dây dẫn vuông góc với từ trường thì lực từ bằng 0.
\item (Sai) Độ lớn lực từ không phụ thuộc chiều dài đoạn dây.
\end{itemchoice}
}
\end{ex}

% ===== Câu 34 =====
% ID: L12C3B20-09-TL
% Mức: VD
\begin{bt}
% ID: L12C3B20-09-TL
% Mức: VD
Nêu một sai lầm thường gặp khi giải tổng hợp lực từ tác dụng lên dây dẫn mang dòng điện và cách khắc phục.
\loigiai{
Với đoạn dây thẳng dài l mang dòng $I$ trong từ trường đều $B$, $F$ = BIl sin(alpha); chiều lực xác định bằng quy tắc bàn tay trái. Sai lầm thường gặp là bỏ qua phương chiều, góc hoặc đơn vị. Khắc phục bằng hình vẽ, quy ước dấu và đổi đơn vị trước khi tính.
}
\end{bt}

% ===== Câu 35 =====
% ID: L12C3B20-09-TLN
% Mức: VD
\begin{ex}
% ID: L12C3B20-09-TLN
% Mức: VD
Trong bốn phát biểu sau về tổng hợp lực từ tác dụng lên dây dẫn mang dòng điện, có bao nhiêu phát biểu đúng? Chỉ ghi một số nguyên. (1) Độ lớn lực từ tác dụng lên đoạn dây thẳng là $F$ = BIl sin(alpha). (2) Khi dây dẫn song song với đường sức từ thì lực từ bằng 0. (3) Khi dây dẫn vuông góc với từ trường thì lực từ bằng 0. (4) Đổi chiều cảm ứng từ nhưng giữ dòng điện thì chiều lực từ không đổi.
\shortans{2}
\loigiai{
Đối chiếu từng phát biểu với kiến thức: Với đoạn dây thẳng dài l mang dòng $I$ trong từ trường đều $B$, $F$ = BIl sin(alpha); chiều lực xác định bằng quy tắc bàn tay trái. Có 2 phát biểu đúng.
}
\end{ex}

% ===== Câu 36 =====
% ID: L12C3B20-09-TN
% Mức: VD
\dangbt{Bài toán vận dụng cao kết hợp nhiều nội dung của Chương III}
\begin{ex}
% ID: L12C3B20-09-TN
% Mức: VD
Khi phân tích bài toán vận dụng cao kết hợp nhiều nội dung của chương iii, phát biểu nào đúng?
\choice
{\True Khi giải bài tổng hợp cần xác định rõ chiều của $B$, dòng điện, lực từ và dòng điện cảm ứng.}
{Có thể cộng đại số mọi vectơ mà không xét phương chiều.}
{Có thể dùng trực tiếp số đo góc theo độ trong mọi biểu thức máy tính mà không kiểm tra chế độ.}
{Hiệu suất của mọi thiết bị điện luôn bằng 100 phần trăm.}
\loigiai{
Khi giải bài tổng hợp cần xác định rõ chiều của $B$, dòng điện, lực từ và dòng điện cảm ứng. Do đó chọn A. Bài toán tổng hợp phải tách các giai đoạn, xác định đúng phương chiều, chọn định luật phù hợp, đổi đơn vị thống nhất và kiểm tra ý nghĩa vật lí của kết quả.
}
\end{ex}

% ===== Câu 37 =====
% ID: L12C3B20-10-DS
% Mức: VD
\dangbt{Tổng hợp lực từ tác dụng lên dây dẫn mang dòng điện}
\begin{ex}
% ID: L12C3B20-10-DS
% Mức: VD
Xét bài toán thuộc dạng tổng hợp lực từ tác dụng lên dây dẫn mang dòng điện. Hãy xác định tính đúng, sai của bốn phát biểu sau.
\choiceTF
{Lực từ luôn cùng chiều với dòng điện.}
{\True Lực từ vuông góc với cả dòng điện và vectơ cảm ứng từ.}
{\True Khi dây dẫn song song với đường sức từ thì lực từ bằng 0.}
{Đổi chiều cảm ứng từ nhưng giữ dòng điện thì chiều lực từ không đổi.}
\loigiai{
\begin{itemchoice}
\item (Sai) Lực từ luôn cùng chiều với dòng điện. (Vì): lpha); chiều lực xác định bằng quy tắc bàn tay trái
\item (Đúng) Lực từ vuông góc với cả dòng điện và vectơ cảm ứng từ.
\item (Đúng) Khi dây dẫn song song với đường sức từ thì lực từ bằng 0.
\item (Sai) Đổi chiều cảm ứng từ nhưng giữ dòng điện thì chiều lực từ không đổi.
\end{itemchoice}
}
\end{ex}

% ===== Câu 38 =====
% ID: L12C3B20-10-TL
% Mức: VD
\begin{bt}
% ID: L12C3B20-10-TL
% Mức: VD
Đề xuất quy trình kiểm tra kết quả của một bài toán tổng hợp lực từ tác dụng lên dây dẫn mang dòng điện.
\loigiai{
Với đoạn dây thẳng dài l mang dòng $I$ trong từ trường đều $B$, $F$ = BIl sin(alpha); chiều lực xác định bằng quy tắc bàn tay trái. Kiểm tra thứ nguyên, dấu, giới hạn đặc biệt và sự phù hợp với ý nghĩa vật lí.
}
\end{bt}

% ===== Câu 39 =====
% ID: L12C3B20-10-TLN
% Mức: VD
\begin{ex}
% ID: L12C3B20-10-TLN
% Mức: VD
Trong bốn phát biểu sau về tổng hợp lực từ tác dụng lên dây dẫn mang dòng điện, có bao nhiêu phát biểu đúng? Chỉ ghi một số nguyên. (1) Độ lớn lực từ tác dụng lên đoạn dây thẳng là $F$ = BIl sin(alpha). (2) Lực từ luôn cùng chiều với dòng điện. (3) Lực từ vuông góc với cả dòng điện và vectơ cảm ứng từ. (4) Khi dây dẫn vuông góc với từ trường thì lực từ bằng 0.
\shortans{2}
\loigiai{
Đối chiếu từng phát biểu với kiến thức: Với đoạn dây thẳng dài l mang dòng $I$ trong từ trường đều $B$, $F$ = BIl sin(alpha); chiều lực xác định bằng quy tắc bàn tay trái. Có 2 phát biểu đúng.
}
\end{ex}

% ===== Câu 40 =====
% ID: L12C3B20-10-TN
% Mức: VD
\dangbt{Bài toán vận dụng cao kết hợp nhiều nội dung của Chương III}
\begin{ex}
% ID: L12C3B20-10-TN
% Mức: VD
Cơ sở Vật lí nào sau đây cần được sử dụng trong dạng bài toán vận dụng cao kết hợp nhiều nội dung của chương iii?
\choice
{Có thể dùng trực tiếp số đo góc theo độ trong mọi biểu thức máy tính mà không kiểm tra chế độ.}
{\True Các đại lượng phải được đổi về cùng hệ đơn vị trước khi thay số.}
{Hiệu suất của mọi thiết bị điện luôn bằng 100 phần trăm.}
{Trong bài tổng hợp có thể bỏ qua điều kiện áp dụng của từng định luật.}
\loigiai{
Các đại lượng phải được đổi về cùng hệ đơn vị trước khi thay số. Do đó chọn B. Bài toán tổng hợp phải tách các giai đoạn, xác định đúng phương chiều, chọn định luật phù hợp, đổi đơn vị thống nhất và kiểm tra ý nghĩa vật lí của kết quả.
}
\end{ex}

% ===== Câu 41 =====
% ID: L12C3B20-100-TN
% Mức: VD
\dangbt{Tổng hợp máy phát điện xoay chiều và đồ thị điện từ}
\begin{ex}
% ID: L12C3B20-100-TN
% Mức: VD
Cơ sở Vật lí nào sau đây cần được sử dụng trong dạng tổng hợp máy phát điện xoay chiều và đồ thị điện từ?
\choice
{Tăng tốc độ quay làm giảm biên độ suất điện động.}
{\True Khung $N$ vòng quay đều có suất điện động cực đại E0 = NBS omega.}
{Từ thông và suất điện động luôn cùng pha.}
{Suất điện động cực đại không phụ thuộc số vòng dây.}
\loigiai{
Khung $N$ vòng quay đều có suất điện động cực đại E0 = NBS omega. Do đó chọn B. Máy phát điện biến cơ năng thành điện năng nhờ cảm ứng điện từ; nếu Phi = Phi0 cos(omega t + phi) thì e = E0 sin(omega t + phi), E0 = NBS omega.
}
\end{ex}

% ===== Câu 42 =====
% ID: L12C3B20-101-TN
% Mức: NB
\dangbt{Tổng hợp máy biến áp, truyền tải điện năng và ứng dụng cảm ứng điện từ}
\begin{ex}
% ID: L12C3B20-101-TN
% Mức: NB
Phát biểu nào sau đây đúng về tổng hợp máy biến áp, truyền tải điện năng và ứng dụng cảm ứng điện từ?
\choice
{\True Máy biến áp hoạt động với dòng điện xoay chiều dựa trên cảm ứng điện từ.}
{Máy biến áp làm việc bình thường với nguồn một chiều không đổi.}
{Máy tăng áp luôn có số vòng cuộn thứ cấp nhỏ hơn cuộn sơ cấp.}
{Hao phí truyền tải tỉ lệ nghịch với bình phương điện trở dây.}
\loigiai{
Máy biến áp hoạt động với dòng điện xoay chiều dựa trên cảm ứng điện từ. Do đó chọn A. Máy biến áp lí tưởng có U2/U1 = N2/N1 và U1I1 = U2I2. Truyền tải cùng công suất ở điện áp cao làm giảm hao phí I^2R.
}
\end{ex}

% ===== Câu 43 =====
% ID: L12C3B20-102-TN
% Mức: NB
\begin{ex}
% ID: L12C3B20-102-TN
% Mức: NB
Trong Chương III, kết luận nào đúng khi giải bài thuộc dạng tổng hợp máy biến áp, truyền tải điện năng và ứng dụng cảm ứng điện từ?
\choice
{Máy tăng áp luôn có số vòng cuộn thứ cấp nhỏ hơn cuộn sơ cấp.}
{\True Máy biến áp lí tưởng thỏa U2/U1 = N2/N1.}
{Hao phí truyền tải tỉ lệ nghịch với bình phương điện trở dây.}
{Lõi máy biến áp ghép từ một khối thép đặc để tăng dòng điện Foucault.}
\loigiai{
Máy biến áp lí tưởng thỏa U2/U1 = N2/N1. Do đó chọn B. Máy biến áp lí tưởng có U2/U1 = N2/N1 và U1I1 = U2I2. Truyền tải cùng công suất ở điện áp cao làm giảm hao phí I^2R.
}
\end{ex}

% ===== Câu 44 =====
% ID: L12C3B20-103-TN
% Mức: NB
\begin{ex}
% ID: L12C3B20-103-TN
% Mức: NB
Chọn nhận định phù hợp với kiến thức Vật lí về tổng hợp máy biến áp, truyền tải điện năng và ứng dụng cảm ứng điện từ.
\choice
{Hao phí truyền tải tỉ lệ nghịch với bình phương điện trở dây.}
{Lõi máy biến áp ghép từ một khối thép đặc để tăng dòng điện Foucault.}
{\True Tăng điện áp truyền tải với cùng công suất làm giảm hao phí trên đường dây.}
{Máy biến áp làm việc bình thường với nguồn một chiều không đổi.}
\loigiai{
Tăng điện áp truyền tải với cùng công suất làm giảm hao phí trên đường dây. Do đó chọn C. Máy biến áp lí tưởng có U2/U1 = N2/N1 và U1I1 = U2I2. Truyền tải cùng công suất ở điện áp cao làm giảm hao phí I^2R.
}
\end{ex}

% ===== Câu 45 =====
% ID: L12C3B20-104-TN
% Mức: TH
\begin{ex}
% ID: L12C3B20-104-TN
% Mức: TH
Khi phân tích tổng hợp máy biến áp, truyền tải điện năng và ứng dụng cảm ứng điện từ, phát biểu nào đúng?
\choice
{Lõi máy biến áp ghép từ một khối thép đặc để tăng dòng điện Foucault.}
{Máy biến áp làm việc bình thường với nguồn một chiều không đổi.}
{Máy tăng áp luôn có số vòng cuộn thứ cấp nhỏ hơn cuộn sơ cấp.}
{\True Dòng điện Foucault được ứng dụng trong bếp từ và phanh điện từ.}
\loigiai{
Dòng điện Foucault được ứng dụng trong bếp từ và phanh điện từ. Do đó chọn D. Máy biến áp lí tưởng có U2/U1 = N2/N1 và U1I1 = U2I2. Truyền tải cùng công suất ở điện áp cao làm giảm hao phí I^2R.
}
\end{ex}

% ===== Câu 46 =====
% ID: L12C3B20-105-TN
% Mức: TH
\begin{ex}
% ID: L12C3B20-105-TN
% Mức: TH
Cơ sở Vật lí nào sau đây cần được sử dụng trong dạng tổng hợp máy biến áp, truyền tải điện năng và ứng dụng cảm ứng điện từ?
\choice
{\True Máy biến áp hoạt động với dòng điện xoay chiều dựa trên cảm ứng điện từ.}
{Máy biến áp làm việc bình thường với nguồn một chiều không đổi.}
{Máy tăng áp luôn có số vòng cuộn thứ cấp nhỏ hơn cuộn sơ cấp.}
{Hao phí truyền tải tỉ lệ nghịch với bình phương điện trở dây.}
\loigiai{
Máy biến áp hoạt động với dòng điện xoay chiều dựa trên cảm ứng điện từ. Do đó chọn A. Máy biến áp lí tưởng có U2/U1 = N2/N1 và U1I1 = U2I2. Truyền tải cùng công suất ở điện áp cao làm giảm hao phí I^2R.
}
\end{ex}

% ===== Câu 47 =====
% ID: L12C3B20-106-TN
% Mức: TH
\begin{ex}
% ID: L12C3B20-106-TN
% Mức: TH
Phát biểu nào sau đây đúng về tổng hợp máy biến áp, truyền tải điện năng và ứng dụng cảm ứng điện từ?
\choice
{Máy tăng áp luôn có số vòng cuộn thứ cấp nhỏ hơn cuộn sơ cấp.}
{\True Máy biến áp lí tưởng thỏa U2/U1 = N2/N1.}
{Hao phí truyền tải tỉ lệ nghịch với bình phương điện trở dây.}
{Lõi máy biến áp ghép từ một khối thép đặc để tăng dòng điện Foucault.}
\loigiai{
Máy biến áp lí tưởng thỏa U2/U1 = N2/N1. Do đó chọn B. Máy biến áp lí tưởng có U2/U1 = N2/N1 và U1I1 = U2I2. Truyền tải cùng công suất ở điện áp cao làm giảm hao phí I^2R.
}
\end{ex}

% ===== Câu 48 =====
% ID: L12C3B20-107-TN
% Mức: TH
\begin{ex}
% ID: L12C3B20-107-TN
% Mức: TH
Trong Chương III, kết luận nào đúng khi giải bài thuộc dạng tổng hợp máy biến áp, truyền tải điện năng và ứng dụng cảm ứng điện từ?
\choice
{Hao phí truyền tải tỉ lệ nghịch với bình phương điện trở dây.}
{Lõi máy biến áp ghép từ một khối thép đặc để tăng dòng điện Foucault.}
{\True Tăng điện áp truyền tải với cùng công suất làm giảm hao phí trên đường dây.}
{Máy biến áp làm việc bình thường với nguồn một chiều không đổi.}
\loigiai{
Tăng điện áp truyền tải với cùng công suất làm giảm hao phí trên đường dây. Do đó chọn C. Máy biến áp lí tưởng có U2/U1 = N2/N1 và U1I1 = U2I2. Truyền tải cùng công suất ở điện áp cao làm giảm hao phí I^2R.
}
\end{ex}

% ===== Câu 49 =====
% ID: L12C3B20-108-TN
% Mức: VD
\begin{ex}
% ID: L12C3B20-108-TN
% Mức: VD
Chọn nhận định phù hợp với kiến thức Vật lí về tổng hợp máy biến áp, truyền tải điện năng và ứng dụng cảm ứng điện từ.
\choice
{Lõi máy biến áp ghép từ một khối thép đặc để tăng dòng điện Foucault.}
{Máy biến áp làm việc bình thường với nguồn một chiều không đổi.}
{Máy tăng áp luôn có số vòng cuộn thứ cấp nhỏ hơn cuộn sơ cấp.}
{\True Dòng điện Foucault được ứng dụng trong bếp từ và phanh điện từ.}
\loigiai{
Dòng điện Foucault được ứng dụng trong bếp từ và phanh điện từ. Do đó chọn D. Máy biến áp lí tưởng có U2/U1 = N2/N1 và U1I1 = U2I2. Truyền tải cùng công suất ở điện áp cao làm giảm hao phí I^2R.
}
\end{ex}

% ===== Câu 50 =====
% ID: L12C3B20-109-TN
% Mức: VD
\begin{ex}
% ID: L12C3B20-109-TN
% Mức: VD
Khi phân tích tổng hợp máy biến áp, truyền tải điện năng và ứng dụng cảm ứng điện từ, phát biểu nào đúng?
\choice
{\True Máy biến áp hoạt động với dòng điện xoay chiều dựa trên cảm ứng điện từ.}
{Máy biến áp làm việc bình thường với nguồn một chiều không đổi.}
{Máy tăng áp luôn có số vòng cuộn thứ cấp nhỏ hơn cuộn sơ cấp.}
{Hao phí truyền tải tỉ lệ nghịch với bình phương điện trở dây.}
\loigiai{
Máy biến áp hoạt động với dòng điện xoay chiều dựa trên cảm ứng điện từ. Do đó chọn A. Máy biến áp lí tưởng có U2/U1 = N2/N1 và U1I1 = U2I2. Truyền tải cùng công suất ở điện áp cao làm giảm hao phí I^2R.
}
\end{ex}

% ===== Câu 51 =====
% ID: L12C3B20-11-DS
% Mức: TH
\begin{ex}
% ID: L12C3B20-11-DS
% Mức: TH
Xét bài toán thuộc dạng tổng hợp máy biến áp, truyền tải điện năng và ứng dụng cảm ứng điện từ. Hãy xác định tính đúng, sai của bốn phát biểu sau.
\choiceTF
{Máy biến áp làm việc bình thường với nguồn một chiều không đổi.}
{Máy tăng áp luôn có số vòng cuộn thứ cấp nhỏ hơn cuộn sơ cấp.}
{\True Tăng điện áp truyền tải với cùng công suất làm giảm hao phí trên đường dây.}
{\True Dòng điện Foucault được ứng dụng trong bếp từ và phanh điện từ.}
\loigiai{
\begin{itemchoice}
\item (Sai) Máy biến áp làm việc bình thường với nguồn một chiều không đổi. (Vì): Máy biến áp làm việc bình thường với nguồn một chiều không đổi. b) Sai: Máy tăng áp luôn có số vòng cuộn thứ cấp nhỏ hơn cuộn sơ cấp. c) Đúng: Tăng điện áp truyền tải với cùng công suất làm giảm hao phí trên đường dây. d) Đúng: Dòng điện Foucault được ứng dụng trong bếp từ và phanh điện từ.. Máy biến áp lí tưởng có U2/U1 = N2/N1 và U1I1 = U2I2. Truyền tải cùng công suất ở điện áp cao làm giảm hao phí I^2R
\item (Sai) Máy tăng áp luôn có số vòng cuộn thứ cấp nhỏ hơn cuộn sơ cấp.
\item (Đúng) Tăng điện áp truyền tải với cùng công suất làm giảm hao phí trên đường dây.
\item (Đúng) Dòng điện Foucault được ứng dụng trong bếp từ và phanh điện từ.
\end{itemchoice}
}
\end{ex}

% ===== Câu 52 =====
% ID: L12C3B20-11-TL
% Mức: VD
\dangbt{Tổng hợp lực từ tác dụng lên dây dẫn mang dòng điện}
\begin{bt}
% ID: L12C3B20-11-TL
% Mức: VD
Xây dựng một tình huống thực tiễn thuộc dạng tổng hợp lực từ tác dụng lên dây dẫn mang dòng điện và chỉ ra định luật cần dùng.
\loigiai{
Với đoạn dây thẳng dài l mang dòng $I$ trong từ trường đều $B$, $F$ = BIl sin(alpha); chiều lực xác định bằng quy tắc bàn tay trái. Có thể chọn thiết bị hoặc hiện tượng thực tế tương ứng, nêu dữ kiện và chỉ rõ định luật dùng để giải thích hoặc tính toán.
}
\end{bt}

% ===== Câu 53 =====
% ID: L12C3B20-11-TLN
% Mức: VD
\begin{ex}
% ID: L12C3B20-11-TLN
% Mức: VD
Trong bốn phát biểu sau về tổng hợp lực từ tác dụng lên dây dẫn mang dòng điện, có bao nhiêu phát biểu đúng? Chỉ ghi một số nguyên. (1) Khi dây dẫn song song với đường sức từ thì lực từ bằng 0. (2) Đổi chiều dòng điện thì chiều lực từ đổi ngược. (3) Độ lớn lực từ không phụ thuộc chiều dài đoạn dây. (4) Đổi chiều cảm ứng từ nhưng giữ dòng điện thì chiều lực từ không đổi.
\shortans{2}
\loigiai{
Đối chiếu từng phát biểu với kiến thức: Với đoạn dây thẳng dài l mang dòng $I$ trong từ trường đều $B$, $F$ = BIl sin(alpha); chiều lực xác định bằng quy tắc bàn tay trái. Có 2 phát biểu đúng.
}
\end{ex}

% ===== Câu 54 =====
% ID: L12C3B20-11-TN
% Mức: NB
\begin{ex}
% ID: L12C3B20-11-TN
% Mức: NB
Phát biểu nào sau đây đúng về tổng hợp lực từ tác dụng lên dây dẫn mang dòng điện?
\choice
{\True Độ lớn lực từ tác dụng lên đoạn dây thẳng là $F$ = BIl sin(alpha).}
{Lực từ luôn cùng chiều với dòng điện.}
{Khi dây dẫn vuông góc với từ trường thì lực từ bằng 0.}
{Độ lớn lực từ không phụ thuộc chiều dài đoạn dây.}
\loigiai{
Độ lớn lực từ tác dụng lên đoạn dây thẳng là $F$ = BIl sin(alpha). Do đó chọn A. Với đoạn dây thẳng dài l mang dòng $I$ trong từ trường đều $B$, $F$ = BIl sin(alpha); chiều lực xác định bằng quy tắc bàn tay trái.
}
\end{ex}

% ===== Câu 55 =====
% ID: L12C3B20-110-TN
% Mức: VD
\dangbt{Tổng hợp máy biến áp, truyền tải điện năng và ứng dụng cảm ứng điện từ}
\begin{ex}
% ID: L12C3B20-110-TN
% Mức: VD
Cơ sở Vật lí nào sau đây cần được sử dụng trong dạng tổng hợp máy biến áp, truyền tải điện năng và ứng dụng cảm ứng điện từ?
\choice
{Máy tăng áp luôn có số vòng cuộn thứ cấp nhỏ hơn cuộn sơ cấp.}
{\True Máy biến áp lí tưởng thỏa U2/U1 = N2/N1.}
{Hao phí truyền tải tỉ lệ nghịch với bình phương điện trở dây.}
{Lõi máy biến áp ghép từ một khối thép đặc để tăng dòng điện Foucault.}
\loigiai{
Máy biến áp lí tưởng thỏa U2/U1 = N2/N1. Do đó chọn B. Máy biến áp lí tưởng có U2/U1 = N2/N1 và U1I1 = U2I2. Truyền tải cùng công suất ở điện áp cao làm giảm hao phí I^2R.
}
\end{ex}

% ===== Câu 56 =====
% ID: L12C3B20-111-TN
% Mức: NB
\dangbt{Bài toán vận dụng cao kết hợp nhiều nội dung của Chương III}
\begin{ex}
% ID: L12C3B20-111-TN
% Mức: NB
Phát biểu nào sau đây đúng về bài toán vận dụng cao kết hợp nhiều nội dung của chương iii?
\choice
{\True Khi giải bài tổng hợp cần xác định rõ chiều của $B$, dòng điện, lực từ và dòng điện cảm ứng.}
{Có thể cộng đại số mọi vectơ mà không xét phương chiều.}
{Có thể dùng trực tiếp số đo góc theo độ trong mọi biểu thức máy tính mà không kiểm tra chế độ.}
{Hiệu suất của mọi thiết bị điện luôn bằng 100 phần trăm.}
\loigiai{
Khi giải bài tổng hợp cần xác định rõ chiều của $B$, dòng điện, lực từ và dòng điện cảm ứng. Do đó chọn A. Bài toán tổng hợp phải tách các giai đoạn, xác định đúng phương chiều, chọn định luật phù hợp, đổi đơn vị thống nhất và kiểm tra ý nghĩa vật lí của kết quả.
}
\end{ex}

% ===== Câu 57 =====
% ID: L12C3B20-112-TN
% Mức: NB
\begin{ex}
% ID: L12C3B20-112-TN
% Mức: NB
Trong Chương III, kết luận nào đúng khi giải bài thuộc dạng bài toán vận dụng cao kết hợp nhiều nội dung của chương iii?
\choice
{Có thể dùng trực tiếp số đo góc theo độ trong mọi biểu thức máy tính mà không kiểm tra chế độ.}
{\True Các đại lượng phải được đổi về cùng hệ đơn vị trước khi thay số.}
{Hiệu suất của mọi thiết bị điện luôn bằng 100 phần trăm.}
{Trong bài tổng hợp có thể bỏ qua điều kiện áp dụng của từng định luật.}
\loigiai{
Các đại lượng phải được đổi về cùng hệ đơn vị trước khi thay số. Do đó chọn B. Bài toán tổng hợp phải tách các giai đoạn, xác định đúng phương chiều, chọn định luật phù hợp, đổi đơn vị thống nhất và kiểm tra ý nghĩa vật lí của kết quả.
}
\end{ex}

% ===== Câu 58 =====
% ID: L12C3B20-113-TN
% Mức: NB
\begin{ex}
% ID: L12C3B20-113-TN
% Mức: NB
Chọn nhận định phù hợp với kiến thức Vật lí về bài toán vận dụng cao kết hợp nhiều nội dung của chương iii.
\choice
{Hiệu suất của mọi thiết bị điện luôn bằng 100 phần trăm.}
{Trong bài tổng hợp có thể bỏ qua điều kiện áp dụng của từng định luật.}
{\True Định luật Lenz quyết định chiều, còn định luật Faraday xác định độ lớn suất điện động cảm ứng.}
{Có thể cộng đại số mọi vectơ mà không xét phương chiều.}
\loigiai{
Định luật Lenz quyết định chiều, còn định luật Faraday xác định độ lớn suất điện động cảm ứng. Do đó chọn C. Bài toán tổng hợp phải tách các giai đoạn, xác định đúng phương chiều, chọn định luật phù hợp, đổi đơn vị thống nhất và kiểm tra ý nghĩa vật lí của kết quả.
}
\end{ex}

% ===== Câu 59 =====
% ID: L12C3B20-114-TN
% Mức: TH
\begin{ex}
% ID: L12C3B20-114-TN
% Mức: TH
Khi phân tích bài toán vận dụng cao kết hợp nhiều nội dung của chương iii, phát biểu nào đúng?
\choice
{Trong bài tổng hợp có thể bỏ qua điều kiện áp dụng của từng định luật.}
{Có thể cộng đại số mọi vectơ mà không xét phương chiều.}
{Có thể dùng trực tiếp số đo góc theo độ trong mọi biểu thức máy tính mà không kiểm tra chế độ.}
{\True Kết quả phải được kiểm tra về đơn vị và ý nghĩa vật lí.}
\loigiai{
Kết quả phải được kiểm tra về đơn vị và ý nghĩa vật lí. Do đó chọn D. Bài toán tổng hợp phải tách các giai đoạn, xác định đúng phương chiều, chọn định luật phù hợp, đổi đơn vị thống nhất và kiểm tra ý nghĩa vật lí của kết quả.
}
\end{ex}

% ===== Câu 60 =====
% ID: L12C3B20-115-TN
% Mức: TH
\begin{ex}
% ID: L12C3B20-115-TN
% Mức: TH
Cơ sở Vật lí nào sau đây cần được sử dụng trong dạng bài toán vận dụng cao kết hợp nhiều nội dung của chương iii?
\choice
{\True Khi giải bài tổng hợp cần xác định rõ chiều của $B$, dòng điện, lực từ và dòng điện cảm ứng.}
{Có thể cộng đại số mọi vectơ mà không xét phương chiều.}
{Có thể dùng trực tiếp số đo góc theo độ trong mọi biểu thức máy tính mà không kiểm tra chế độ.}
{Hiệu suất của mọi thiết bị điện luôn bằng 100 phần trăm.}
\loigiai{
Khi giải bài tổng hợp cần xác định rõ chiều của $B$, dòng điện, lực từ và dòng điện cảm ứng. Do đó chọn A. Bài toán tổng hợp phải tách các giai đoạn, xác định đúng phương chiều, chọn định luật phù hợp, đổi đơn vị thống nhất và kiểm tra ý nghĩa vật lí của kết quả.
}
\end{ex}

% ===== Câu 61 =====
% ID: L12C3B20-116-TN
% Mức: TH
\begin{ex}
% ID: L12C3B20-116-TN
% Mức: TH
Phát biểu nào sau đây đúng về bài toán vận dụng cao kết hợp nhiều nội dung của chương iii?
\choice
{Có thể dùng trực tiếp số đo góc theo độ trong mọi biểu thức máy tính mà không kiểm tra chế độ.}
{\True Các đại lượng phải được đổi về cùng hệ đơn vị trước khi thay số.}
{Hiệu suất của mọi thiết bị điện luôn bằng 100 phần trăm.}
{Trong bài tổng hợp có thể bỏ qua điều kiện áp dụng của từng định luật.}
\loigiai{
Các đại lượng phải được đổi về cùng hệ đơn vị trước khi thay số. Do đó chọn B. Bài toán tổng hợp phải tách các giai đoạn, xác định đúng phương chiều, chọn định luật phù hợp, đổi đơn vị thống nhất và kiểm tra ý nghĩa vật lí của kết quả.
}
\end{ex}

% ===== Câu 62 =====
% ID: L12C3B20-117-TN
% Mức: TH
\begin{ex}
% ID: L12C3B20-117-TN
% Mức: TH
Trong Chương III, kết luận nào đúng khi giải bài thuộc dạng bài toán vận dụng cao kết hợp nhiều nội dung của chương iii?
\choice
{Hiệu suất của mọi thiết bị điện luôn bằng 100 phần trăm.}
{Trong bài tổng hợp có thể bỏ qua điều kiện áp dụng của từng định luật.}
{\True Định luật Lenz quyết định chiều, còn định luật Faraday xác định độ lớn suất điện động cảm ứng.}
{Có thể cộng đại số mọi vectơ mà không xét phương chiều.}
\loigiai{
Định luật Lenz quyết định chiều, còn định luật Faraday xác định độ lớn suất điện động cảm ứng. Do đó chọn C. Bài toán tổng hợp phải tách các giai đoạn, xác định đúng phương chiều, chọn định luật phù hợp, đổi đơn vị thống nhất và kiểm tra ý nghĩa vật lí của kết quả.
}
\end{ex}

% ===== Câu 63 =====
% ID: L12C3B20-118-TN
% Mức: VD
\begin{ex}
% ID: L12C3B20-118-TN
% Mức: VD
Chọn nhận định phù hợp với kiến thức Vật lí về bài toán vận dụng cao kết hợp nhiều nội dung của chương iii.
\choice
{Trong bài tổng hợp có thể bỏ qua điều kiện áp dụng của từng định luật.}
{Có thể cộng đại số mọi vectơ mà không xét phương chiều.}
{Có thể dùng trực tiếp số đo góc theo độ trong mọi biểu thức máy tính mà không kiểm tra chế độ.}
{\True Kết quả phải được kiểm tra về đơn vị và ý nghĩa vật lí.}
\loigiai{
Kết quả phải được kiểm tra về đơn vị và ý nghĩa vật lí. Do đó chọn D. Bài toán tổng hợp phải tách các giai đoạn, xác định đúng phương chiều, chọn định luật phù hợp, đổi đơn vị thống nhất và kiểm tra ý nghĩa vật lí của kết quả.
}
\end{ex}

% ===== Câu 64 =====
% ID: L12C3B20-119-TN
% Mức: VD
\begin{ex}
% ID: L12C3B20-119-TN
% Mức: VD
Khi phân tích bài toán vận dụng cao kết hợp nhiều nội dung của chương iii, phát biểu nào đúng?
\choice
{\True Khi giải bài tổng hợp cần xác định rõ chiều của $B$, dòng điện, lực từ và dòng điện cảm ứng.}
{Có thể cộng đại số mọi vectơ mà không xét phương chiều.}
{Có thể dùng trực tiếp số đo góc theo độ trong mọi biểu thức máy tính mà không kiểm tra chế độ.}
{Hiệu suất của mọi thiết bị điện luôn bằng 100 phần trăm.}
\loigiai{
Khi giải bài tổng hợp cần xác định rõ chiều của $B$, dòng điện, lực từ và dòng điện cảm ứng. Do đó chọn A. Bài toán tổng hợp phải tách các giai đoạn, xác định đúng phương chiều, chọn định luật phù hợp, đổi đơn vị thống nhất và kiểm tra ý nghĩa vật lí của kết quả.
}
\end{ex}

% ===== Câu 65 =====
% ID: L12C3B20-12-DS
% Mức: TH
\dangbt{Tổng hợp máy biến áp, truyền tải điện năng và ứng dụng cảm ứng điện từ}
\begin{ex}
% ID: L12C3B20-12-DS
% Mức: TH
Xét bài toán thuộc dạng tổng hợp máy biến áp, truyền tải điện năng và ứng dụng cảm ứng điện từ. Hãy xác định tính đúng, sai của bốn phát biểu sau.
\choiceTF
{\True Máy biến áp lí tưởng thỏa U2/U1 = N2/N1.}
{Hao phí truyền tải tỉ lệ nghịch với bình phương điện trở dây.}
{\True Dòng điện Foucault được ứng dụng trong bếp từ và phanh điện từ.}
{Máy biến áp làm việc bình thường với nguồn một chiều không đổi.}
\loigiai{
\begin{itemchoice}
\item (Đúng) Máy biến áp lí tưởng thỏa U2/U1 = N2/N1. (Vì): Máy biến áp lí tưởng thỏa U2/U1 = N2/N1. b) Sai: Hao phí truyền tải tỉ lệ nghịch với bình phương điện trở dây. c) Đúng: Dòng điện Foucault được ứng dụng trong bếp từ và phanh điện từ. d) Sai: Máy biến áp làm việc bình thường với nguồn một chiều không đổi.. Máy biến áp lí tưởng có U2/U1 = N2/N1 và U1I1 = U2I2. Truyền tải cùng công suất ở điện áp cao làm giảm hao phí I^2R
\item (Sai) Hao phí truyền tải tỉ lệ nghịch với bình phương điện trở dây.
\item (Đúng) Dòng điện Foucault được ứng dụng trong bếp từ và phanh điện từ.
\item (Sai) Máy biến áp làm việc bình thường với nguồn một chiều không đổi.
\end{itemchoice}
}
\end{ex}

% ===== Câu 66 =====
% ID: L12C3B20-12-TL
% Mức: VDC
\dangbt{Tổng hợp lực từ tác dụng lên dây dẫn mang dòng điện}
\begin{bt}
% ID: L12C3B20-12-TL
% Mức: VDC
So sánh một nhận định đúng và một nhận định sai trong dạng tổng hợp lực từ tác dụng lên dây dẫn mang dòng điện; giải thích căn cứ Vật lí.
\loigiai{
Nhận định đúng: Độ lớn lực từ tác dụng lên đoạn dây thẳng là $F$ = BIl sin(alpha). Nhận định sai: Lực từ luôn cùng chiều với dòng điện. Căn cứ: Với đoạn dây thẳng dài l mang dòng $I$ trong từ trường đều $B$, $F$ = BIl sin(alpha); chiều lực xác định bằng quy tắc bàn tay trái.
}
\end{bt}

% ===== Câu 67 =====
% ID: L12C3B20-12-TLN
% Mức: VDC
\begin{ex}
% ID: L12C3B20-12-TLN
% Mức: VDC
Trong bốn phát biểu sau về tổng hợp lực từ tác dụng lên dây dẫn mang dòng điện, có bao nhiêu phát biểu đúng? Chỉ ghi một số nguyên. (1) Độ lớn lực từ tác dụng lên đoạn dây thẳng là $F$ = BIl sin(alpha). (2) Khi dây dẫn song song với đường sức từ thì lực từ bằng 0. (3) Khi dây dẫn vuông góc với từ trường thì lực từ bằng 0. (4) Đổi chiều cảm ứng từ nhưng giữ dòng điện thì chiều lực từ không đổi.
\shortans{2}
\loigiai{
Đối chiếu từng phát biểu với kiến thức: Với đoạn dây thẳng dài l mang dòng $I$ trong từ trường đều $B$, $F$ = BIl sin(alpha); chiều lực xác định bằng quy tắc bàn tay trái. Có 2 phát biểu đúng.
}
\end{ex}

% ===== Câu 68 =====
% ID: L12C3B20-12-TN
% Mức: NB
\begin{ex}
% ID: L12C3B20-12-TN
% Mức: NB
Trong Chương III, kết luận nào đúng khi giải bài thuộc dạng tổng hợp lực từ tác dụng lên dây dẫn mang dòng điện?
\choice
{Khi dây dẫn vuông góc với từ trường thì lực từ bằng 0.}
{\True Lực từ vuông góc với cả dòng điện và vectơ cảm ứng từ.}
{Độ lớn lực từ không phụ thuộc chiều dài đoạn dây.}
{Đổi chiều cảm ứng từ nhưng giữ dòng điện thì chiều lực từ không đổi.}
\loigiai{
Lực từ vuông góc với cả dòng điện và vectơ cảm ứng từ. Do đó chọn B. Với đoạn dây thẳng dài l mang dòng $I$ trong từ trường đều $B$, $F$ = BIl sin(alpha); chiều lực xác định bằng quy tắc bàn tay trái.
}
\end{ex}

% ===== Câu 69 =====
% ID: L12C3B20-120-TN
% Mức: VD
\dangbt{Bài toán vận dụng cao kết hợp nhiều nội dung của Chương III}
\begin{ex}
% ID: L12C3B20-120-TN
% Mức: VD
Cơ sở Vật lí nào sau đây cần được sử dụng trong dạng bài toán vận dụng cao kết hợp nhiều nội dung của chương iii?
\choice
{Có thể dùng trực tiếp số đo góc theo độ trong mọi biểu thức máy tính mà không kiểm tra chế độ.}
{\True Các đại lượng phải được đổi về cùng hệ đơn vị trước khi thay số.}
{Hiệu suất của mọi thiết bị điện luôn bằng 100 phần trăm.}
{Trong bài tổng hợp có thể bỏ qua điều kiện áp dụng của từng định luật.}
\loigiai{
Các đại lượng phải được đổi về cùng hệ đơn vị trước khi thay số. Do đó chọn B. Bài toán tổng hợp phải tách các giai đoạn, xác định đúng phương chiều, chọn định luật phù hợp, đổi đơn vị thống nhất và kiểm tra ý nghĩa vật lí của kết quả.
}
\end{ex}

% ===== Câu 70 =====
% ID: L12C3B20-13-DS
% Mức: VD
\dangbt{Tổng hợp máy biến áp, truyền tải điện năng và ứng dụng cảm ứng điện từ}
\begin{ex}
% ID: L12C3B20-13-DS
% Mức: VD
Xét bài toán thuộc dạng tổng hợp máy biến áp, truyền tải điện năng và ứng dụng cảm ứng điện từ. Hãy xác định tính đúng, sai của bốn phát biểu sau.
\choiceTF
{Hao phí truyền tải tỉ lệ nghịch với bình phương điện trở dây.}
{\True Dòng điện Foucault được ứng dụng trong bếp từ và phanh điện từ.}
{Máy biến áp làm việc bình thường với nguồn một chiều không đổi.}
{\True Máy biến áp lí tưởng thỏa U2/U1 = N2/N1.}
\loigiai{
\begin{itemchoice}
\item (Sai) Hao phí truyền tải tỉ lệ nghịch với bình phương điện trở dây. (Vì): Hao phí truyền tải tỉ lệ nghịch với bình phương điện trở dây. b) Đúng: Dòng điện Foucault được ứng dụng trong bếp từ và phanh điện từ. c) Sai: Máy biến áp làm việc bình thường với nguồn một chiều không đổi. d) Đúng: Máy biến áp lí tưởng thỏa U2/U1 = N2/N1.. Máy biến áp lí tưởng có U2/U1 = N2/N1 và U1I1 = U2I2. Truyền tải cùng công suất ở điện áp cao làm giảm hao phí I^2R
\item (Đúng) Dòng điện Foucault được ứng dụng trong bếp từ và phanh điện từ.
\item (Sai) Máy biến áp làm việc bình thường với nguồn một chiều không đổi.
\item (Đúng) Máy biến áp lí tưởng thỏa U2/U1 = N2/N1.
\end{itemchoice}
}
\end{ex}

% ===== Câu 71 =====
% ID: L12C3B20-13-TL
% Mức: TH
\begin{bt}
% ID: L12C3B20-13-TL
% Mức: TH
Trình bày kiến thức cốt lõi và các bước giải bài thuộc dạng tổng hợp máy biến áp, truyền tải điện năng và ứng dụng cảm ứng điện từ.
\loigiai{
Máy biến áp lí tưởng có U2/U1 = N2/N1 và U1I1 = U2I2. Truyền tải cùng công suất ở điện áp cao làm giảm hao phí I^2R. Các bước: tóm tắt dữ kiện; xác định phương, chiều và điều kiện; chọn hệ thức; đổi đơn vị; tính toán; kiểm tra kết quả.
}
\end{bt}

% ===== Câu 72 =====
% ID: L12C3B20-13-TLN
% Mức: TH
\begin{ex}
% ID: L12C3B20-13-TLN
% Mức: TH
Trong bốn phát biểu sau về tổng hợp máy biến áp, truyền tải điện năng và ứng dụng cảm ứng điện từ, có bao nhiêu phát biểu đúng? Chỉ ghi một số nguyên. (1) Máy biến áp hoạt động với dòng điện xoay chiều dựa trên cảm ứng điện từ. (2) Máy biến áp làm việc bình thường với nguồn một chiều không đổi. (3) Máy biến áp lí tưởng thỏa U2/U1 = N2/N1. (4) Máy tăng áp luôn có số vòng cuộn thứ cấp nhỏ hơn cuộn sơ cấp.
\shortans{2}
\loigiai{
Đối chiếu từng phát biểu với kiến thức: Máy biến áp lí tưởng có U2/U1 = N2/N1 và U1I1 = U2I2. Truyền tải cùng công suất ở điện áp cao làm giảm hao phí I^2R. Có 2 phát biểu đúng.
}
\end{ex}

% ===== Câu 73 =====
% ID: L12C3B20-13-TN
% Mức: NB
\dangbt{Tổng hợp lực từ tác dụng lên dây dẫn mang dòng điện}
\begin{ex}
% ID: L12C3B20-13-TN
% Mức: NB
Chọn nhận định phù hợp với kiến thức Vật lí về tổng hợp lực từ tác dụng lên dây dẫn mang dòng điện.
\choice
{Độ lớn lực từ không phụ thuộc chiều dài đoạn dây.}
{Đổi chiều cảm ứng từ nhưng giữ dòng điện thì chiều lực từ không đổi.}
{\True Khi dây dẫn song song với đường sức từ thì lực từ bằng 0.}
{Lực từ luôn cùng chiều với dòng điện.}
\loigiai{
Khi dây dẫn song song với đường sức từ thì lực từ bằng 0. Do đó chọn C. Với đoạn dây thẳng dài l mang dòng $I$ trong từ trường đều $B$, $F$ = BIl sin(alpha); chiều lực xác định bằng quy tắc bàn tay trái.
}
\end{ex}

% ===== Câu 74 =====
% ID: L12C3B20-14-DS
% Mức: VD
\dangbt{Tổng hợp máy biến áp, truyền tải điện năng và ứng dụng cảm ứng điện từ}
\begin{ex}
% ID: L12C3B20-14-DS
% Mức: VD
Xét bài toán thuộc dạng tổng hợp máy biến áp, truyền tải điện năng và ứng dụng cảm ứng điện từ. Hãy xác định tính đúng, sai của bốn phát biểu sau.
\choiceTF
{\True Dòng điện Foucault được ứng dụng trong bếp từ và phanh điện từ.}
{\True Máy biến áp hoạt động với dòng điện xoay chiều dựa trên cảm ứng điện từ.}
{Máy tăng áp luôn có số vòng cuộn thứ cấp nhỏ hơn cuộn sơ cấp.}
{Hao phí truyền tải tỉ lệ nghịch với bình phương điện trở dây.}
\loigiai{
\begin{itemchoice}
\item (Đúng) Dòng điện Foucault được ứng dụng trong bếp từ và phanh điện từ. (Vì): òng điện Foucault được ứng dụng trong bếp từ và phanh điện từ. b) Đúng: Máy biến áp hoạt động với dòng điện xoay chiều dựa trên cảm ứng điện từ. c) Sai: Máy tăng áp luôn có số vòng cuộn thứ cấp nhỏ hơn cuộn sơ cấp. d) Sai: Hao phí truyền tải tỉ lệ nghịch với bình phương điện trở dây.. Máy biến áp lí tưởng có U2/U1 = N2/N1 và U1I1 = U2I2. Truyền tải cùng công suất ở điện áp cao làm giảm hao phí I^2R
\item (Đúng) Máy biến áp hoạt động với dòng điện xoay chiều dựa trên cảm ứng điện từ.
\item (Sai) Máy tăng áp luôn có số vòng cuộn thứ cấp nhỏ hơn cuộn sơ cấp.
\item (Sai) Hao phí truyền tải tỉ lệ nghịch với bình phương điện trở dây.
\end{itemchoice}
}
\end{ex}

% ===== Câu 75 =====
% ID: L12C3B20-14-TL
% Mức: TH
\begin{bt}
% ID: L12C3B20-14-TL
% Mức: TH
Phân tích các đại lượng, phương chiều và điều kiện áp dụng trong dạng tổng hợp máy biến áp, truyền tải điện năng và ứng dụng cảm ứng điện từ.
\loigiai{
Máy biến áp lí tưởng có U2/U1 = N2/N1 và U1I1 = U2I2. Truyền tải cùng công suất ở điện áp cao làm giảm hao phí I^2R. Cần xác định rõ đại lượng vô hướng hay vectơ, quy ước chiều và điều kiện để công thức được áp dụng.
}
\end{bt}

% ===== Câu 76 =====
% ID: L12C3B20-14-TLN
% Mức: TH
\begin{ex}
% ID: L12C3B20-14-TLN
% Mức: TH
Trong bốn phát biểu sau về tổng hợp máy biến áp, truyền tải điện năng và ứng dụng cảm ứng điện từ, có bao nhiêu phát biểu đúng? Chỉ ghi một số nguyên. (1) Tăng điện áp truyền tải với cùng công suất làm giảm hao phí trên đường dây. (2) Dòng điện Foucault được ứng dụng trong bếp từ và phanh điện từ. (3) Hao phí truyền tải tỉ lệ nghịch với bình phương điện trở dây. (4) Lõi máy biến áp ghép từ một khối thép đặc để tăng dòng điện Foucault.
\shortans{2}
\loigiai{
Đối chiếu từng phát biểu với kiến thức: Máy biến áp lí tưởng có U2/U1 = N2/N1 và U1I1 = U2I2. Truyền tải cùng công suất ở điện áp cao làm giảm hao phí I^2R. Có 2 phát biểu đúng.
}
\end{ex}

% ===== Câu 77 =====
% ID: L12C3B20-14-TN
% Mức: TH
\dangbt{Tổng hợp lực từ tác dụng lên dây dẫn mang dòng điện}
\begin{ex}
% ID: L12C3B20-14-TN
% Mức: TH
Khi phân tích tổng hợp lực từ tác dụng lên dây dẫn mang dòng điện, phát biểu nào đúng?
\choice
{Đổi chiều cảm ứng từ nhưng giữ dòng điện thì chiều lực từ không đổi.}
{Lực từ luôn cùng chiều với dòng điện.}
{Khi dây dẫn vuông góc với từ trường thì lực từ bằng 0.}
{\True Đổi chiều dòng điện thì chiều lực từ đổi ngược.}
\loigiai{
Đổi chiều dòng điện thì chiều lực từ đổi ngược. Do đó chọn D. Với đoạn dây thẳng dài l mang dòng $I$ trong từ trường đều $B$, $F$ = BIl sin(alpha); chiều lực xác định bằng quy tắc bàn tay trái.
}
\end{ex}

% ===== Câu 78 =====
% ID: L12C3B20-15-DS
% Mức: VD
\dangbt{Tổng hợp máy biến áp, truyền tải điện năng và ứng dụng cảm ứng điện từ}
\begin{ex}
% ID: L12C3B20-15-DS
% Mức: VD
Xét bài toán thuộc dạng tổng hợp máy biến áp, truyền tải điện năng và ứng dụng cảm ứng điện từ. Hãy xác định tính đúng, sai của bốn phát biểu sau.
\choiceTF
{Máy biến áp làm việc bình thường với nguồn một chiều không đổi.}
{\True Máy biến áp lí tưởng thỏa U2/U1 = N2/N1.}
{\True Tăng điện áp truyền tải với cùng công suất làm giảm hao phí trên đường dây.}
{Lõi máy biến áp ghép từ một khối thép đặc để tăng dòng điện Foucault.}
\loigiai{
\begin{itemchoice}
\item (Sai) Máy biến áp làm việc bình thường với nguồn một chiều không đổi. (Vì): Máy biến áp làm việc bình thường với nguồn một chiều không đổi. b) Đúng: Máy biến áp lí tưởng thỏa U2/U1 = N2/N1. c) Đúng: Tăng điện áp truyền tải với cùng công suất làm giảm hao phí trên đường dây. d) Sai: Lõi máy biến áp ghép từ một khối thép đặc để tăng dòng điện Foucault.. Máy biến áp lí tưởng có U2/U1 = N2/N1 và U1I1 = U2I2. Truyền tải cùng công suất ở điện áp cao làm giảm hao phí I^2R
\item (Đúng) Máy biến áp lí tưởng thỏa U2/U1 = N2/N1.
\item (Đúng) Tăng điện áp truyền tải với cùng công suất làm giảm hao phí trên đường dây.
\item (Sai) Lõi máy biến áp ghép từ một khối thép đặc để tăng dòng điện Foucault.
\end{itemchoice}
}
\end{ex}

% ===== Câu 79 =====
% ID: L12C3B20-15-TL
% Mức: VD
\begin{bt}
% ID: L12C3B20-15-TL
% Mức: VD
Nêu một sai lầm thường gặp khi giải tổng hợp máy biến áp, truyền tải điện năng và ứng dụng cảm ứng điện từ và cách khắc phục.
\loigiai{
Máy biến áp lí tưởng có U2/U1 = N2/N1 và U1I1 = U2I2. Truyền tải cùng công suất ở điện áp cao làm giảm hao phí I^2R. Sai lầm thường gặp là bỏ qua phương chiều, góc hoặc đơn vị. Khắc phục bằng hình vẽ, quy ước dấu và đổi đơn vị trước khi tính.
}
\end{bt}

% ===== Câu 80 =====
% ID: L12C3B20-15-TLN
% Mức: VD
\begin{ex}
% ID: L12C3B20-15-TLN
% Mức: VD
Trong bốn phát biểu sau về tổng hợp máy biến áp, truyền tải điện năng và ứng dụng cảm ứng điện từ, có bao nhiêu phát biểu đúng? Chỉ ghi một số nguyên. (1) Máy biến áp hoạt động với dòng điện xoay chiều dựa trên cảm ứng điện từ. (2) Tăng điện áp truyền tải với cùng công suất làm giảm hao phí trên đường dây. (3) Máy tăng áp luôn có số vòng cuộn thứ cấp nhỏ hơn cuộn sơ cấp. (4) Lõi máy biến áp ghép từ một khối thép đặc để tăng dòng điện Foucault.
\shortans{2}
\loigiai{
Đối chiếu từng phát biểu với kiến thức: Máy biến áp lí tưởng có U2/U1 = N2/N1 và U1I1 = U2I2. Truyền tải cùng công suất ở điện áp cao làm giảm hao phí I^2R. Có 2 phát biểu đúng.
}
\end{ex}

% ===== Câu 81 =====
% ID: L12C3B20-15-TN
% Mức: TH
\dangbt{Tổng hợp lực từ tác dụng lên dây dẫn mang dòng điện}
\begin{ex}
% ID: L12C3B20-15-TN
% Mức: TH
Cơ sở Vật lí nào sau đây cần được sử dụng trong dạng tổng hợp lực từ tác dụng lên dây dẫn mang dòng điện?
\choice
{\True Độ lớn lực từ tác dụng lên đoạn dây thẳng là $F$ = BIl sin(alpha).}
{Lực từ luôn cùng chiều với dòng điện.}
{Khi dây dẫn vuông góc với từ trường thì lực từ bằng 0.}
{Độ lớn lực từ không phụ thuộc chiều dài đoạn dây.}
\loigiai{
Độ lớn lực từ tác dụng lên đoạn dây thẳng là $F$ = BIl sin(alpha). Do đó chọn A. Với đoạn dây thẳng dài l mang dòng $I$ trong từ trường đều $B$, $F$ = BIl sin(alpha); chiều lực xác định bằng quy tắc bàn tay trái.
}
\end{ex}

% ===== Câu 82 =====
% ID: L12C3B20-16-DS
% Mức: TH
\dangbt{Tổng hợp máy phát điện xoay chiều và đồ thị điện từ}
\begin{ex}
% ID: L12C3B20-16-DS
% Mức: TH
Xét bài toán thuộc dạng tổng hợp máy phát điện xoay chiều và đồ thị điện từ. Hãy xác định tính đúng, sai của bốn phát biểu sau.
\choiceTF
{Máy phát điện biến điện năng thành cơ năng.}
{Tăng tốc độ quay làm giảm biên độ suất điện động.}
{\True Tần số suất điện động bằng tần số quay của khung trong mô hình một cặp cực.}
{\True Suất điện động cảm ứng lệch pha pi/2 so với từ thông.}
\loigiai{
\begin{itemchoice}
\item (Sai) Máy phát điện biến điện năng thành cơ năng. (Vì): Máy phát điện biến điện năng thành cơ năng. b) Sai: Tăng tốc độ quay làm giảm biên độ suất điện động. c) Đúng: Tần số suất điện động bằng tần số quay của khung trong mô hình một cặp cực. d) Đúng: Suất điện động cảm ứng lệch pha pi/2 so với từ thông.. Máy phát điện biến cơ năng thành điện năng nhờ cảm ứng điện từ; nếu Phi = Phi0 cos(omega t + phi) thì e = E0 sin(omega t + phi), E0 = NBS omega
\item (Sai) Tăng tốc độ quay làm giảm biên độ suất điện động.
\item (Đúng) Tần số suất điện động bằng tần số quay của khung trong mô hình một cặp cực.
\item (Đúng) Suất điện động cảm ứng lệch pha pi/2 so với từ thông.
\end{itemchoice}
}
\end{ex}

% ===== Câu 83 =====
% ID: L12C3B20-16-TL
% Mức: VD
\dangbt{Tổng hợp máy biến áp, truyền tải điện năng và ứng dụng cảm ứng điện từ}
\begin{bt}
% ID: L12C3B20-16-TL
% Mức: VD
Đề xuất quy trình kiểm tra kết quả của một bài toán tổng hợp máy biến áp, truyền tải điện năng và ứng dụng cảm ứng điện từ.
\loigiai{
Máy biến áp lí tưởng có U2/U1 = N2/N1 và U1I1 = U2I2. Truyền tải cùng công suất ở điện áp cao làm giảm hao phí I^2R. Kiểm tra thứ nguyên, dấu, giới hạn đặc biệt và sự phù hợp với ý nghĩa vật lí.
}
\end{bt}

% ===== Câu 84 =====
% ID: L12C3B20-16-TLN
% Mức: VD
\begin{ex}
% ID: L12C3B20-16-TLN
% Mức: VD
Trong bốn phát biểu sau về tổng hợp máy biến áp, truyền tải điện năng và ứng dụng cảm ứng điện từ, có bao nhiêu phát biểu đúng? Chỉ ghi một số nguyên. (1) Máy biến áp hoạt động với dòng điện xoay chiều dựa trên cảm ứng điện từ. (2) Máy biến áp làm việc bình thường với nguồn một chiều không đổi. (3) Máy biến áp lí tưởng thỏa U2/U1 = N2/N1. (4) Máy tăng áp luôn có số vòng cuộn thứ cấp nhỏ hơn cuộn sơ cấp.
\shortans{2}
\loigiai{
Đối chiếu từng phát biểu với kiến thức: Máy biến áp lí tưởng có U2/U1 = N2/N1 và U1I1 = U2I2. Truyền tải cùng công suất ở điện áp cao làm giảm hao phí I^2R. Có 2 phát biểu đúng.
}
\end{ex}

% ===== Câu 85 =====
% ID: L12C3B20-16-TN
% Mức: TH
\dangbt{Tổng hợp lực từ tác dụng lên dây dẫn mang dòng điện}
\begin{ex}
% ID: L12C3B20-16-TN
% Mức: TH
Phát biểu nào sau đây đúng về tổng hợp lực từ tác dụng lên dây dẫn mang dòng điện?
\choice
{Khi dây dẫn vuông góc với từ trường thì lực từ bằng 0.}
{\True Lực từ vuông góc với cả dòng điện và vectơ cảm ứng từ.}
{Độ lớn lực từ không phụ thuộc chiều dài đoạn dây.}
{Đổi chiều cảm ứng từ nhưng giữ dòng điện thì chiều lực từ không đổi.}
\loigiai{
Lực từ vuông góc với cả dòng điện và vectơ cảm ứng từ. Do đó chọn B. Với đoạn dây thẳng dài l mang dòng $I$ trong từ trường đều $B$, $F$ = BIl sin(alpha); chiều lực xác định bằng quy tắc bàn tay trái.
}
\end{ex}

% ===== Câu 86 =====
% ID: L12C3B20-17-DS
% Mức: TH
\dangbt{Tổng hợp máy phát điện xoay chiều và đồ thị điện từ}
\begin{ex}
% ID: L12C3B20-17-DS
% Mức: TH
Xét bài toán thuộc dạng tổng hợp máy phát điện xoay chiều và đồ thị điện từ. Hãy xác định tính đúng, sai của bốn phát biểu sau.
\choiceTF
{\True Khung $N$ vòng quay đều có suất điện động cực đại E0 = NBS omega.}
{Từ thông và suất điện động luôn cùng pha.}
{\True Suất điện động cảm ứng lệch pha pi/2 so với từ thông.}
{Máy phát điện biến điện năng thành cơ năng.}
\loigiai{
\begin{itemchoice}
\item (Đúng) Khung $N$ vòng quay đều có suất điện động cực đại E0 = NBS omega. (Vì): Khung $N$ vòng quay đều có suất điện động cực đại E0 = NBS omega. b) Sai: Từ thông và suất điện động luôn cùng pha. c) Đúng: Suất điện động cảm ứng lệch pha pi/2 so với từ thông. d) Sai: Máy phát điện biến điện năng thành cơ năng.. Máy phát điện biến cơ năng thành điện năng nhờ cảm ứng điện từ; nếu Phi = Phi0 cos(omega t + phi) thì e = E0 sin(omega t + phi), E0 = NBS omega
\item (Sai) Từ thông và suất điện động luôn cùng pha.
\item (Đúng) Suất điện động cảm ứng lệch pha pi/2 so với từ thông.
\item (Sai) Máy phát điện biến điện năng thành cơ năng.
\end{itemchoice}
}
\end{ex}

% ===== Câu 87 =====
% ID: L12C3B20-17-TL
% Mức: VD
\dangbt{Tổng hợp máy biến áp, truyền tải điện năng và ứng dụng cảm ứng điện từ}
\begin{bt}
% ID: L12C3B20-17-TL
% Mức: VD
Xây dựng một tình huống thực tiễn thuộc dạng tổng hợp máy biến áp, truyền tải điện năng và ứng dụng cảm ứng điện từ và chỉ ra định luật cần dùng.
\loigiai{
Máy biến áp lí tưởng có U2/U1 = N2/N1 và U1I1 = U2I2. Truyền tải cùng công suất ở điện áp cao làm giảm hao phí I^2R. Có thể chọn thiết bị hoặc hiện tượng thực tế tương ứng, nêu dữ kiện và chỉ rõ định luật dùng để giải thích hoặc tính toán.
}
\end{bt}

% ===== Câu 88 =====
% ID: L12C3B20-17-TLN
% Mức: VD
\begin{ex}
% ID: L12C3B20-17-TLN
% Mức: VD
Trong bốn phát biểu sau về tổng hợp máy biến áp, truyền tải điện năng và ứng dụng cảm ứng điện từ, có bao nhiêu phát biểu đúng? Chỉ ghi một số nguyên. (1) Tăng điện áp truyền tải với cùng công suất làm giảm hao phí trên đường dây. (2) Dòng điện Foucault được ứng dụng trong bếp từ và phanh điện từ. (3) Hao phí truyền tải tỉ lệ nghịch với bình phương điện trở dây. (4) Lõi máy biến áp ghép từ một khối thép đặc để tăng dòng điện Foucault.
\shortans{2}
\loigiai{
Đối chiếu từng phát biểu với kiến thức: Máy biến áp lí tưởng có U2/U1 = N2/N1 và U1I1 = U2I2. Truyền tải cùng công suất ở điện áp cao làm giảm hao phí I^2R. Có 2 phát biểu đúng.
}
\end{ex}

% ===== Câu 89 =====
% ID: L12C3B20-17-TN
% Mức: TH
\dangbt{Tổng hợp lực từ tác dụng lên dây dẫn mang dòng điện}
\begin{ex}
% ID: L12C3B20-17-TN
% Mức: TH
Trong Chương III, kết luận nào đúng khi giải bài thuộc dạng tổng hợp lực từ tác dụng lên dây dẫn mang dòng điện?
\choice
{Độ lớn lực từ không phụ thuộc chiều dài đoạn dây.}
{Đổi chiều cảm ứng từ nhưng giữ dòng điện thì chiều lực từ không đổi.}
{\True Khi dây dẫn song song với đường sức từ thì lực từ bằng 0.}
{Lực từ luôn cùng chiều với dòng điện.}
\loigiai{
Khi dây dẫn song song với đường sức từ thì lực từ bằng 0. Do đó chọn C. Với đoạn dây thẳng dài l mang dòng $I$ trong từ trường đều $B$, $F$ = BIl sin(alpha); chiều lực xác định bằng quy tắc bàn tay trái.
}
\end{ex}

% ===== Câu 90 =====
% ID: L12C3B20-18-DS
% Mức: VD
\dangbt{Tổng hợp máy phát điện xoay chiều và đồ thị điện từ}
\begin{ex}
% ID: L12C3B20-18-DS
% Mức: VD
Xét bài toán thuộc dạng tổng hợp máy phát điện xoay chiều và đồ thị điện từ. Hãy xác định tính đúng, sai của bốn phát biểu sau.
\choiceTF
{Từ thông và suất điện động luôn cùng pha.}
{\True Suất điện động cảm ứng lệch pha pi/2 so với từ thông.}
{Máy phát điện biến điện năng thành cơ năng.}
{\True Khung $N$ vòng quay đều có suất điện động cực đại E0 = NBS omega.}
\loigiai{
\begin{itemchoice}
\item (Sai) Từ thông và suất điện động luôn cùng pha. (Vì): Từ thông và suất điện động luôn cùng pha. b) Đúng: Suất điện động cảm ứng lệch pha pi/2 so với từ thông. c) Sai: Máy phát điện biến điện năng thành cơ năng. d) Đúng: Khung $N$ vòng quay đều có suất điện động cực đại E0 = NBS omega.. Máy phát điện biến cơ năng thành điện năng nhờ cảm ứng điện từ; nếu Phi = Phi0 cos(omega t + phi) thì e = E0 sin(omega t + phi), E0 = NBS omega
\item (Đúng) Suất điện động cảm ứng lệch pha pi/2 so với từ thông.
\item (Sai) Máy phát điện biến điện năng thành cơ năng.
\item (Đúng) Khung $N$ vòng quay đều có suất điện động cực đại E0 = NBS omega.
\end{itemchoice}
}
\end{ex}

% ===== Câu 91 =====
% ID: L12C3B20-18-TL
% Mức: VDC
\dangbt{Tổng hợp máy biến áp, truyền tải điện năng và ứng dụng cảm ứng điện từ}
\begin{bt}
% ID: L12C3B20-18-TL
% Mức: VDC
So sánh một nhận định đúng và một nhận định sai trong dạng tổng hợp máy biến áp, truyền tải điện năng và ứng dụng cảm ứng điện từ; giải thích căn cứ Vật lí.
\loigiai{
Nhận định đúng: Máy biến áp hoạt động với dòng điện xoay chiều dựa trên cảm ứng điện từ. Nhận định sai: Máy biến áp làm việc bình thường với nguồn một chiều không đổi. Căn cứ: Máy biến áp lí tưởng có U2/U1 = N2/N1 và U1I1 = U2I2. Truyền tải cùng công suất ở điện áp cao làm giảm hao phí I^2R.
}
\end{bt}

% ===== Câu 92 =====
% ID: L12C3B20-18-TLN
% Mức: VDC
\begin{ex}
% ID: L12C3B20-18-TLN
% Mức: VDC
Trong bốn phát biểu sau về tổng hợp máy biến áp, truyền tải điện năng và ứng dụng cảm ứng điện từ, có bao nhiêu phát biểu đúng? Chỉ ghi một số nguyên. (1) Máy biến áp hoạt động với dòng điện xoay chiều dựa trên cảm ứng điện từ. (2) Tăng điện áp truyền tải với cùng công suất làm giảm hao phí trên đường dây. (3) Máy tăng áp luôn có số vòng cuộn thứ cấp nhỏ hơn cuộn sơ cấp. (4) Lõi máy biến áp ghép từ một khối thép đặc để tăng dòng điện Foucault.
\shortans{2}
\loigiai{
Đối chiếu từng phát biểu với kiến thức: Máy biến áp lí tưởng có U2/U1 = N2/N1 và U1I1 = U2I2. Truyền tải cùng công suất ở điện áp cao làm giảm hao phí I^2R. Có 2 phát biểu đúng.
}
\end{ex}

% ===== Câu 93 =====
% ID: L12C3B20-18-TN
% Mức: VD
\dangbt{Tổng hợp lực từ tác dụng lên dây dẫn mang dòng điện}
\begin{ex}
% ID: L12C3B20-18-TN
% Mức: VD
Chọn nhận định phù hợp với kiến thức Vật lí về tổng hợp lực từ tác dụng lên dây dẫn mang dòng điện.
\choice
{Đổi chiều cảm ứng từ nhưng giữ dòng điện thì chiều lực từ không đổi.}
{Lực từ luôn cùng chiều với dòng điện.}
{Khi dây dẫn vuông góc với từ trường thì lực từ bằng 0.}
{\True Đổi chiều dòng điện thì chiều lực từ đổi ngược.}
\loigiai{
Đổi chiều dòng điện thì chiều lực từ đổi ngược. Do đó chọn D. Với đoạn dây thẳng dài l mang dòng $I$ trong từ trường đều $B$, $F$ = BIl sin(alpha); chiều lực xác định bằng quy tắc bàn tay trái.
}
\end{ex}

% ===== Câu 94 =====
% ID: L12C3B20-19-DS
% Mức: VD
\dangbt{Tổng hợp máy phát điện xoay chiều và đồ thị điện từ}
\begin{ex}
% ID: L12C3B20-19-DS
% Mức: VD
Xét bài toán thuộc dạng tổng hợp máy phát điện xoay chiều và đồ thị điện từ. Hãy xác định tính đúng, sai của bốn phát biểu sau.
\choiceTF
{\True Suất điện động cảm ứng lệch pha pi/2 so với từ thông.}
{\True Máy phát điện xoay chiều hoạt động dựa trên hiện tượng cảm ứng điện từ.}
{Tăng tốc độ quay làm giảm biên độ suất điện động.}
{Từ thông và suất điện động luôn cùng pha.}
\loigiai{
\begin{itemchoice}
\item (Đúng) Suất điện động cảm ứng lệch pha pi/2 so với từ thông. (Vì): Suất điện động cảm ứng lệch pha pi/2 so với từ thông. b) Đúng: Máy phát điện xoay chiều hoạt động dựa trên hiện tượng cảm ứng điện từ. c) Sai: Tăng tốc độ quay làm giảm biên độ suất điện động. d) Sai: Từ thông và suất điện động luôn cùng pha.. Máy phát điện biến cơ năng thành điện năng nhờ cảm ứng điện từ; nếu Phi = Phi0 cos(omega t + phi) thì e = E0 sin(omega t + phi), E0 = NBS omega
\item (Đúng) Máy phát điện xoay chiều hoạt động dựa trên hiện tượng cảm ứng điện từ.
\item (Sai) Tăng tốc độ quay làm giảm biên độ suất điện động.
\item (Sai) Từ thông và suất điện động luôn cùng pha.
\end{itemchoice}
}
\end{ex}

% ===== Câu 95 =====
% ID: L12C3B20-19-TL
% Mức: TH
\begin{bt}
% ID: L12C3B20-19-TL
% Mức: TH
Trình bày kiến thức cốt lõi và các bước giải bài thuộc dạng tổng hợp máy phát điện xoay chiều và đồ thị điện từ.
\loigiai{
Máy phát điện biến cơ năng thành điện năng nhờ cảm ứng điện từ; nếu Phi = Phi0 cos(omega t + phi) thì e = E0 sin(omega t + phi), E0 = NBS omega. Các bước: tóm tắt dữ kiện; xác định phương, chiều và điều kiện; chọn hệ thức; đổi đơn vị; tính toán; kiểm tra kết quả.
}
\end{bt}

% ===== Câu 96 =====
% ID: L12C3B20-19-TLN
% Mức: TH
\begin{ex}
% ID: L12C3B20-19-TLN
% Mức: TH
Trong bốn phát biểu sau về tổng hợp máy phát điện xoay chiều và đồ thị điện từ, có bao nhiêu phát biểu đúng? Chỉ ghi một số nguyên. (1) Máy phát điện xoay chiều hoạt động dựa trên hiện tượng cảm ứng điện từ. (2) Máy phát điện biến điện năng thành cơ năng. (3) Khung $N$ vòng quay đều có suất điện động cực đại E0 = NBS omega. (4) Tăng tốc độ quay làm giảm biên độ suất điện động.
\shortans{2}
\loigiai{
Đối chiếu từng phát biểu với kiến thức: Máy phát điện biến cơ năng thành điện năng nhờ cảm ứng điện từ; nếu Phi = Phi0 cos(omega t + phi) thì e = E0 sin(omega t + phi), E0 = NBS omega. Có 2 phát biểu đúng.
}
\end{ex}

% ===== Câu 97 =====
% ID: L12C3B20-19-TN
% Mức: VD
\dangbt{Tổng hợp lực từ tác dụng lên dây dẫn mang dòng điện}
\begin{ex}
% ID: L12C3B20-19-TN
% Mức: VD
Khi phân tích tổng hợp lực từ tác dụng lên dây dẫn mang dòng điện, phát biểu nào đúng?
\choice
{\True Độ lớn lực từ tác dụng lên đoạn dây thẳng là $F$ = BIl sin(alpha).}
{Lực từ luôn cùng chiều với dòng điện.}
{Khi dây dẫn vuông góc với từ trường thì lực từ bằng 0.}
{Độ lớn lực từ không phụ thuộc chiều dài đoạn dây.}
\loigiai{
Độ lớn lực từ tác dụng lên đoạn dây thẳng là $F$ = BIl sin(alpha). Do đó chọn A. Với đoạn dây thẳng dài l mang dòng $I$ trong từ trường đều $B$, $F$ = BIl sin(alpha); chiều lực xác định bằng quy tắc bàn tay trái.
}
\end{ex}

% ===== Câu 98 =====
% ID: L12C3B20-20-DS
% Mức: VD
\dangbt{Tổng hợp máy phát điện xoay chiều và đồ thị điện từ}
\begin{ex}
% ID: L12C3B20-20-DS
% Mức: VD
Xét bài toán thuộc dạng tổng hợp máy phát điện xoay chiều và đồ thị điện từ. Hãy xác định tính đúng, sai của bốn phát biểu sau.
\choiceTF
{Máy phát điện biến điện năng thành cơ năng.}
{\True Khung $N$ vòng quay đều có suất điện động cực đại E0 = NBS omega.}
{\True Tần số suất điện động bằng tần số quay của khung trong mô hình một cặp cực.}
{Suất điện động cực đại không phụ thuộc số vòng dây.}
\loigiai{
\begin{itemchoice}
\item (Sai) Máy phát điện biến điện năng thành cơ năng. (Vì): Máy phát điện biến điện năng thành cơ năng. b) Đúng: Khung $N$ vòng quay đều có suất điện động cực đại E0 = NBS omega. c) Đúng: Tần số suất điện động bằng tần số quay của khung trong mô hình một cặp cực. d) Sai: Suất điện động cực đại không phụ thuộc số vòng dây.. Máy phát điện biến cơ năng thành điện năng nhờ cảm ứng điện từ; nếu Phi = Phi0 cos(omega t + phi) thì e = E0 sin(omega t + phi), E0 = NBS omega
\item (Đúng) Khung $N$ vòng quay đều có suất điện động cực đại E0 = NBS omega.
\item (Đúng) Tần số suất điện động bằng tần số quay của khung trong mô hình một cặp cực.
\item (Sai) Suất điện động cực đại không phụ thuộc số vòng dây.
\end{itemchoice}
}
\end{ex}

% ===== Câu 99 =====
% ID: L12C3B20-20-TL
% Mức: TH
\begin{bt}
% ID: L12C3B20-20-TL
% Mức: TH
Phân tích các đại lượng, phương chiều và điều kiện áp dụng trong dạng tổng hợp máy phát điện xoay chiều và đồ thị điện từ.
\loigiai{
Máy phát điện biến cơ năng thành điện năng nhờ cảm ứng điện từ; nếu Phi = Phi0 cos(omega t + phi) thì e = E0 sin(omega t + phi), E0 = NBS omega. Cần xác định rõ đại lượng vô hướng hay vectơ, quy ước chiều và điều kiện để công thức được áp dụng.
}
\end{bt}

% ===== Câu 100 =====
% ID: L12C3B20-20-TLN
% Mức: TH
\begin{ex}
% ID: L12C3B20-20-TLN
% Mức: TH
Trong bốn phát biểu sau về tổng hợp máy phát điện xoay chiều và đồ thị điện từ, có bao nhiêu phát biểu đúng? Chỉ ghi một số nguyên. (1) Tần số suất điện động bằng tần số quay của khung trong mô hình một cặp cực. (2) Suất điện động cảm ứng lệch pha pi/2 so với từ thông. (3) Từ thông và suất điện động luôn cùng pha. (4) Suất điện động cực đại không phụ thuộc số vòng dây.
\shortans{2}
\loigiai{
Đối chiếu từng phát biểu với kiến thức: Máy phát điện biến cơ năng thành điện năng nhờ cảm ứng điện từ; nếu Phi = Phi0 cos(omega t + phi) thì e = E0 sin(omega t + phi), E0 = NBS omega. Có 2 phát biểu đúng.
}
\end{ex}

% ===== Câu 101 =====
% ID: L12C3B20-20-TN
% Mức: VD
\dangbt{Tổng hợp lực từ tác dụng lên dây dẫn mang dòng điện}
\begin{ex}
% ID: L12C3B20-20-TN
% Mức: VD
Cơ sở Vật lí nào sau đây cần được sử dụng trong dạng tổng hợp lực từ tác dụng lên dây dẫn mang dòng điện?
\choice
{Khi dây dẫn vuông góc với từ trường thì lực từ bằng 0.}
{\True Lực từ vuông góc với cả dòng điện và vectơ cảm ứng từ.}
{Độ lớn lực từ không phụ thuộc chiều dài đoạn dây.}
{Đổi chiều cảm ứng từ nhưng giữ dòng điện thì chiều lực từ không đổi.}
\loigiai{
Lực từ vuông góc với cả dòng điện và vectơ cảm ứng từ. Do đó chọn B. Với đoạn dây thẳng dài l mang dòng $I$ trong từ trường đều $B$, $F$ = BIl sin(alpha); chiều lực xác định bằng quy tắc bàn tay trái.
}
\end{ex}

% ===== Câu 102 =====
% ID: L12C3B20-21-DS
% Mức: TH
\dangbt{Tổng hợp từ thông, định luật Lenz và suất điện động cảm ứng}
\begin{ex}
% ID: L12C3B20-21-DS
% Mức: TH
Xét bài toán thuộc dạng tổng hợp từ thông, định luật lenz và suất điện động cảm ứng. Hãy xác định tính đúng, sai của bốn phát biểu sau.
\choiceTF
{Từ thông là đại lượng vectơ.}
{Dòng điện cảm ứng luôn làm tăng từ thông ban đầu.}
{\True Độ lớn suất điện động cảm ứng bằng độ lớn tốc độ biến thiên từ thông.}
{\True Từ thông không đổi thì không xuất hiện suất điện động cảm ứng.}
\loigiai{
\begin{itemchoice}
\item (Sai) Từ thông là đại lượng vectơ. (Vì): lpha). Suất điện động cảm ứng có độ lớn |e| = |Delta Phi/Delta t|; chiều dòng điện cảm ứng tuân theo định luật Lenz
\item (Sai) Dòng điện cảm ứng luôn làm tăng từ thông ban đầu.
\item (Đúng) Độ lớn suất điện động cảm ứng bằng độ lớn tốc độ biến thiên từ thông.
\item (Đúng) Từ thông không đổi thì không xuất hiện suất điện động cảm ứng.
\end{itemchoice}
}
\end{ex}

% ===== Câu 103 =====
% ID: L12C3B20-21-TL
% Mức: VD
\dangbt{Tổng hợp máy phát điện xoay chiều và đồ thị điện từ}
\begin{bt}
% ID: L12C3B20-21-TL
% Mức: VD
Nêu một sai lầm thường gặp khi giải tổng hợp máy phát điện xoay chiều và đồ thị điện từ và cách khắc phục.
\loigiai{
Máy phát điện biến cơ năng thành điện năng nhờ cảm ứng điện từ; nếu Phi = Phi0 cos(omega t + phi) thì e = E0 sin(omega t + phi), E0 = NBS omega. Sai lầm thường gặp là bỏ qua phương chiều, góc hoặc đơn vị. Khắc phục bằng hình vẽ, quy ước dấu và đổi đơn vị trước khi tính.
}
\end{bt}

% ===== Câu 104 =====
% ID: L12C3B20-21-TLN
% Mức: VD
\begin{ex}
% ID: L12C3B20-21-TLN
% Mức: VD
Trong bốn phát biểu sau về tổng hợp máy phát điện xoay chiều và đồ thị điện từ, có bao nhiêu phát biểu đúng? Chỉ ghi một số nguyên. (1) Máy phát điện xoay chiều hoạt động dựa trên hiện tượng cảm ứng điện từ. (2) Tần số suất điện động bằng tần số quay của khung trong mô hình một cặp cực. (3) Tăng tốc độ quay làm giảm biên độ suất điện động. (4) Suất điện động cực đại không phụ thuộc số vòng dây.
\shortans{2}
\loigiai{
Đối chiếu từng phát biểu với kiến thức: Máy phát điện biến cơ năng thành điện năng nhờ cảm ứng điện từ; nếu Phi = Phi0 cos(omega t + phi) thì e = E0 sin(omega t + phi), E0 = NBS omega. Có 2 phát biểu đúng.
}
\end{ex}

% ===== Câu 105 =====
% ID: L12C3B20-21-TN
% Mức: NB
\dangbt{Tổng hợp máy biến áp, truyền tải điện năng và ứng dụng cảm ứng điện từ}
\begin{ex}
% ID: L12C3B20-21-TN
% Mức: NB
Phát biểu nào sau đây đúng về tổng hợp máy biến áp, truyền tải điện năng và ứng dụng cảm ứng điện từ?
\choice
{\True Máy biến áp hoạt động với dòng điện xoay chiều dựa trên cảm ứng điện từ.}
{Máy biến áp làm việc bình thường với nguồn một chiều không đổi.}
{Máy tăng áp luôn có số vòng cuộn thứ cấp nhỏ hơn cuộn sơ cấp.}
{Hao phí truyền tải tỉ lệ nghịch với bình phương điện trở dây.}
\loigiai{
Máy biến áp hoạt động với dòng điện xoay chiều dựa trên cảm ứng điện từ. Do đó chọn A. Máy biến áp lí tưởng có U2/U1 = N2/N1 và U1I1 = U2I2. Truyền tải cùng công suất ở điện áp cao làm giảm hao phí I^2R.
}
\end{ex}

% ===== Câu 106 =====
% ID: L12C3B20-22-DS
% Mức: TH
\dangbt{Tổng hợp từ thông, định luật Lenz và suất điện động cảm ứng}
\begin{ex}
% ID: L12C3B20-22-DS
% Mức: TH
Xét bài toán thuộc dạng tổng hợp từ thông, định luật lenz và suất điện động cảm ứng. Hãy xác định tính đúng, sai của bốn phát biểu sau.
\choiceTF
{\True Dòng điện cảm ứng có chiều chống lại sự biến thiên từ thông sinh ra nó.}
{Suất điện động cảm ứng chỉ phụ thuộc điện trở của mạch.}
{\True Từ thông không đổi thì không xuất hiện suất điện động cảm ứng.}
{Từ thông là đại lượng vectơ.}
\loigiai{
\begin{itemchoice}
\item (Đúng) Dòng điện cảm ứng có chiều chống lại sự biến thiên từ thông sinh ra nó. (Vì): lpha). Suất điện động cảm ứng có độ lớn |e| = |Delta Phi/Delta t|; chiều dòng điện cảm ứng tuân theo định luật Lenz
\item (Sai) Suất điện động cảm ứng chỉ phụ thuộc điện trở của mạch.
\item (Đúng) Từ thông không đổi thì không xuất hiện suất điện động cảm ứng.
\item (Sai) Từ thông là đại lượng vectơ.
\end{itemchoice}
}
\end{ex}

% ===== Câu 107 =====
% ID: L12C3B20-22-TL
% Mức: VD
\dangbt{Tổng hợp máy phát điện xoay chiều và đồ thị điện từ}
\begin{bt}
% ID: L12C3B20-22-TL
% Mức: VD
Đề xuất quy trình kiểm tra kết quả của một bài toán tổng hợp máy phát điện xoay chiều và đồ thị điện từ.
\loigiai{
Máy phát điện biến cơ năng thành điện năng nhờ cảm ứng điện từ; nếu Phi = Phi0 cos(omega t + phi) thì e = E0 sin(omega t + phi), E0 = NBS omega. Kiểm tra thứ nguyên, dấu, giới hạn đặc biệt và sự phù hợp với ý nghĩa vật lí.
}
\end{bt}

% ===== Câu 108 =====
% ID: L12C3B20-22-TLN
% Mức: VD
\begin{ex}
% ID: L12C3B20-22-TLN
% Mức: VD
Trong bốn phát biểu sau về tổng hợp máy phát điện xoay chiều và đồ thị điện từ, có bao nhiêu phát biểu đúng? Chỉ ghi một số nguyên. (1) Máy phát điện xoay chiều hoạt động dựa trên hiện tượng cảm ứng điện từ. (2) Máy phát điện biến điện năng thành cơ năng. (3) Khung $N$ vòng quay đều có suất điện động cực đại E0 = NBS omega. (4) Tăng tốc độ quay làm giảm biên độ suất điện động.
\shortans{2}
\loigiai{
Đối chiếu từng phát biểu với kiến thức: Máy phát điện biến cơ năng thành điện năng nhờ cảm ứng điện từ; nếu Phi = Phi0 cos(omega t + phi) thì e = E0 sin(omega t + phi), E0 = NBS omega. Có 2 phát biểu đúng.
}
\end{ex}

% ===== Câu 109 =====
% ID: L12C3B20-22-TN
% Mức: NB
\dangbt{Tổng hợp máy biến áp, truyền tải điện năng và ứng dụng cảm ứng điện từ}
\begin{ex}
% ID: L12C3B20-22-TN
% Mức: NB
Trong Chương III, kết luận nào đúng khi giải bài thuộc dạng tổng hợp máy biến áp, truyền tải điện năng và ứng dụng cảm ứng điện từ?
\choice
{Máy tăng áp luôn có số vòng cuộn thứ cấp nhỏ hơn cuộn sơ cấp.}
{\True Máy biến áp lí tưởng thỏa U2/U1 = N2/N1.}
{Hao phí truyền tải tỉ lệ nghịch với bình phương điện trở dây.}
{Lõi máy biến áp ghép từ một khối thép đặc để tăng dòng điện Foucault.}
\loigiai{
Máy biến áp lí tưởng thỏa U2/U1 = N2/N1. Do đó chọn B. Máy biến áp lí tưởng có U2/U1 = N2/N1 và U1I1 = U2I2. Truyền tải cùng công suất ở điện áp cao làm giảm hao phí I^2R.
}
\end{ex}

% ===== Câu 110 =====
% ID: L12C3B20-23-DS
% Mức: VD
\dangbt{Tổng hợp từ thông, định luật Lenz và suất điện động cảm ứng}
\begin{ex}
% ID: L12C3B20-23-DS
% Mức: VD
Xét bài toán thuộc dạng tổng hợp từ thông, định luật lenz và suất điện động cảm ứng. Hãy xác định tính đúng, sai của bốn phát biểu sau.
\choiceTF
{Suất điện động cảm ứng chỉ phụ thuộc điện trở của mạch.}
{\True Từ thông không đổi thì không xuất hiện suất điện động cảm ứng.}
{Từ thông là đại lượng vectơ.}
{\True Dòng điện cảm ứng có chiều chống lại sự biến thiên từ thông sinh ra nó.}
\loigiai{
\begin{itemchoice}
\item (Sai) Suất điện động cảm ứng chỉ phụ thuộc điện trở của mạch. (Vì): lpha). Suất điện động cảm ứng có độ lớn |e| = |Delta Phi/Delta t|; chiều dòng điện cảm ứng tuân theo định luật Lenz
\item (Đúng) Từ thông không đổi thì không xuất hiện suất điện động cảm ứng.
\item (Sai) Từ thông là đại lượng vectơ.
\item (Đúng) Dòng điện cảm ứng có chiều chống lại sự biến thiên từ thông sinh ra nó.
\end{itemchoice}
}
\end{ex}

% ===== Câu 111 =====
% ID: L12C3B20-23-TL
% Mức: VD
\dangbt{Tổng hợp máy phát điện xoay chiều và đồ thị điện từ}
\begin{bt}
% ID: L12C3B20-23-TL
% Mức: VD
Xây dựng một tình huống thực tiễn thuộc dạng tổng hợp máy phát điện xoay chiều và đồ thị điện từ và chỉ ra định luật cần dùng.
\loigiai{
Máy phát điện biến cơ năng thành điện năng nhờ cảm ứng điện từ; nếu Phi = Phi0 cos(omega t + phi) thì e = E0 sin(omega t + phi), E0 = NBS omega. Có thể chọn thiết bị hoặc hiện tượng thực tế tương ứng, nêu dữ kiện và chỉ rõ định luật dùng để giải thích hoặc tính toán.
}
\end{bt}

% ===== Câu 112 =====
% ID: L12C3B20-23-TLN
% Mức: VD
\begin{ex}
% ID: L12C3B20-23-TLN
% Mức: VD
Trong bốn phát biểu sau về tổng hợp máy phát điện xoay chiều và đồ thị điện từ, có bao nhiêu phát biểu đúng? Chỉ ghi một số nguyên. (1) Tần số suất điện động bằng tần số quay của khung trong mô hình một cặp cực. (2) Suất điện động cảm ứng lệch pha pi/2 so với từ thông. (3) Từ thông và suất điện động luôn cùng pha. (4) Suất điện động cực đại không phụ thuộc số vòng dây.
\shortans{2}
\loigiai{
Đối chiếu từng phát biểu với kiến thức: Máy phát điện biến cơ năng thành điện năng nhờ cảm ứng điện từ; nếu Phi = Phi0 cos(omega t + phi) thì e = E0 sin(omega t + phi), E0 = NBS omega. Có 2 phát biểu đúng.
}
\end{ex}

% ===== Câu 113 =====
% ID: L12C3B20-23-TN
% Mức: NB
\dangbt{Tổng hợp máy biến áp, truyền tải điện năng và ứng dụng cảm ứng điện từ}
\begin{ex}
% ID: L12C3B20-23-TN
% Mức: NB
Chọn nhận định phù hợp với kiến thức Vật lí về tổng hợp máy biến áp, truyền tải điện năng và ứng dụng cảm ứng điện từ.
\choice
{Hao phí truyền tải tỉ lệ nghịch với bình phương điện trở dây.}
{Lõi máy biến áp ghép từ một khối thép đặc để tăng dòng điện Foucault.}
{\True Tăng điện áp truyền tải với cùng công suất làm giảm hao phí trên đường dây.}
{Máy biến áp làm việc bình thường với nguồn một chiều không đổi.}
\loigiai{
Tăng điện áp truyền tải với cùng công suất làm giảm hao phí trên đường dây. Do đó chọn C. Máy biến áp lí tưởng có U2/U1 = N2/N1 và U1I1 = U2I2. Truyền tải cùng công suất ở điện áp cao làm giảm hao phí I^2R.
}
\end{ex}

% ===== Câu 114 =====
% ID: L12C3B20-24-DS
% Mức: VD
\dangbt{Tổng hợp từ thông, định luật Lenz và suất điện động cảm ứng}
\begin{ex}
% ID: L12C3B20-24-DS
% Mức: VD
Xét bài toán thuộc dạng tổng hợp từ thông, định luật lenz và suất điện động cảm ứng. Hãy xác định tính đúng, sai của bốn phát biểu sau.
\choiceTF
{\True Từ thông không đổi thì không xuất hiện suất điện động cảm ứng.}
{\True Từ thông qua khung phẳng là Phi = BS cos(alpha), với alpha là góc giữa $B$ và pháp tuyến.}
{Dòng điện cảm ứng luôn làm tăng từ thông ban đầu.}
{Suất điện động cảm ứng chỉ phụ thuộc điện trở của mạch.}
\loigiai{
\begin{itemchoice}
\item (Đúng) Từ thông không đổi thì không xuất hiện suất điện động cảm ứng. (Vì): lpha). Suất điện động cảm ứng có độ lớn |e| = |Delta Phi/Delta t|; chiều dòng điện cảm ứng tuân theo định luật Lenz
\item (Đúng) Từ thông qua khung phẳng là Phi = BS cos(alpha), với alpha là góc giữa $B$ và pháp tuyến.
\item (Sai) Dòng điện cảm ứng luôn làm tăng từ thông ban đầu.
\item (Sai) Suất điện động cảm ứng chỉ phụ thuộc điện trở của mạch.
\end{itemchoice}
}
\end{ex}

% ===== Câu 115 =====
% ID: L12C3B20-24-TL
% Mức: VDC
\dangbt{Tổng hợp máy phát điện xoay chiều và đồ thị điện từ}
\begin{bt}
% ID: L12C3B20-24-TL
% Mức: VDC
So sánh một nhận định đúng và một nhận định sai trong dạng tổng hợp máy phát điện xoay chiều và đồ thị điện từ; giải thích căn cứ Vật lí.
\loigiai{
Nhận định đúng: Máy phát điện xoay chiều hoạt động dựa trên hiện tượng cảm ứng điện từ. Nhận định sai: Máy phát điện biến điện năng thành cơ năng. Căn cứ: Máy phát điện biến cơ năng thành điện năng nhờ cảm ứng điện từ; nếu Phi = Phi0 cos(omega t + phi) thì e = E0 sin(omega t + phi), E0 = NBS omega.
}
\end{bt}

% ===== Câu 116 =====
% ID: L12C3B20-24-TLN
% Mức: VDC
\begin{ex}
% ID: L12C3B20-24-TLN
% Mức: VDC
Trong bốn phát biểu sau về tổng hợp máy phát điện xoay chiều và đồ thị điện từ, có bao nhiêu phát biểu đúng? Chỉ ghi một số nguyên. (1) Máy phát điện xoay chiều hoạt động dựa trên hiện tượng cảm ứng điện từ. (2) Tần số suất điện động bằng tần số quay của khung trong mô hình một cặp cực. (3) Tăng tốc độ quay làm giảm biên độ suất điện động. (4) Suất điện động cực đại không phụ thuộc số vòng dây.
\shortans{2}
\loigiai{
Đối chiếu từng phát biểu với kiến thức: Máy phát điện biến cơ năng thành điện năng nhờ cảm ứng điện từ; nếu Phi = Phi0 cos(omega t + phi) thì e = E0 sin(omega t + phi), E0 = NBS omega. Có 2 phát biểu đúng.
}
\end{ex}

% ===== Câu 117 =====
% ID: L12C3B20-24-TN
% Mức: TH
\dangbt{Tổng hợp máy biến áp, truyền tải điện năng và ứng dụng cảm ứng điện từ}
\begin{ex}
% ID: L12C3B20-24-TN
% Mức: TH
Khi phân tích tổng hợp máy biến áp, truyền tải điện năng và ứng dụng cảm ứng điện từ, phát biểu nào đúng?
\choice
{Lõi máy biến áp ghép từ một khối thép đặc để tăng dòng điện Foucault.}
{Máy biến áp làm việc bình thường với nguồn một chiều không đổi.}
{Máy tăng áp luôn có số vòng cuộn thứ cấp nhỏ hơn cuộn sơ cấp.}
{\True Dòng điện Foucault được ứng dụng trong bếp từ và phanh điện từ.}
\loigiai{
Dòng điện Foucault được ứng dụng trong bếp từ và phanh điện từ. Do đó chọn D. Máy biến áp lí tưởng có U2/U1 = N2/N1 và U1I1 = U2I2. Truyền tải cùng công suất ở điện áp cao làm giảm hao phí I^2R.
}
\end{ex}

% ===== Câu 118 =====
% ID: L12C3B20-25-DS
% Mức: VD
\dangbt{Tổng hợp từ thông, định luật Lenz và suất điện động cảm ứng}
\begin{ex}
% ID: L12C3B20-25-DS
% Mức: VD
Xét bài toán thuộc dạng tổng hợp từ thông, định luật lenz và suất điện động cảm ứng. Hãy xác định tính đúng, sai của bốn phát biểu sau.
\choiceTF
{Từ thông là đại lượng vectơ.}
{\True Dòng điện cảm ứng có chiều chống lại sự biến thiên từ thông sinh ra nó.}
{\True Độ lớn suất điện động cảm ứng bằng độ lớn tốc độ biến thiên từ thông.}
{Khung dây đứng yên trong từ trường đều không đổi luôn có dòng điện cảm ứng.}
\loigiai{
\begin{itemchoice}
\item (Sai) Từ thông là đại lượng vectơ. (Vì): lpha). Suất điện động cảm ứng có độ lớn |e| = |Delta Phi/Delta t|; chiều dòng điện cảm ứng tuân theo định luật Lenz
\item (Đúng) Dòng điện cảm ứng có chiều chống lại sự biến thiên từ thông sinh ra nó.
\item (Đúng) Độ lớn suất điện động cảm ứng bằng độ lớn tốc độ biến thiên từ thông.
\item (Sai) Khung dây đứng yên trong từ trường đều không đổi luôn có dòng điện cảm ứng.
\end{itemchoice}
}
\end{ex}

% ===== Câu 119 =====
% ID: L12C3B20-25-TL
% Mức: TH
\begin{bt}
% ID: L12C3B20-25-TL
% Mức: TH
Trình bày kiến thức cốt lõi và các bước giải bài thuộc dạng tổng hợp từ thông, định luật lenz và suất điện động cảm ứng.
\loigiai{
Phi = BS cos(alpha). Suất điện động cảm ứng có độ lớn |e| = |Delta Phi/Delta t|; chiều dòng điện cảm ứng tuân theo định luật Lenz. Các bước: tóm tắt dữ kiện; xác định phương, chiều và điều kiện; chọn hệ thức; đổi đơn vị; tính toán; kiểm tra kết quả.
}
\end{bt}

% ===== Câu 120 =====
% ID: L12C3B20-25-TLN
% Mức: TH
\begin{ex}
% ID: L12C3B20-25-TLN
% Mức: TH
Trong bốn phát biểu sau về tổng hợp từ thông, định luật lenz và suất điện động cảm ứng, có bao nhiêu phát biểu đúng? Chỉ ghi một số nguyên. (1) Từ thông qua khung phẳng là Phi = BS cos(alpha), với alpha là góc giữa $B$ và pháp tuyến. (2) Từ thông là đại lượng vectơ. (3) Dòng điện cảm ứng có chiều chống lại sự biến thiên từ thông sinh ra nó. (4) Dòng điện cảm ứng luôn làm tăng từ thông ban đầu.
\shortans{2}
\loigiai{
Đối chiếu từng phát biểu với kiến thức: Phi = BS cos(alpha). Suất điện động cảm ứng có độ lớn |e| = |Delta Phi/Delta t|; chiều dòng điện cảm ứng tuân theo định luật Lenz. Có 2 phát biểu đúng.
}
\end{ex}

% ===== Câu 121 =====
% ID: L12C3B20-25-TN
% Mức: TH
\dangbt{Tổng hợp máy biến áp, truyền tải điện năng và ứng dụng cảm ứng điện từ}
\begin{ex}
% ID: L12C3B20-25-TN
% Mức: TH
Cơ sở Vật lí nào sau đây cần được sử dụng trong dạng tổng hợp máy biến áp, truyền tải điện năng và ứng dụng cảm ứng điện từ?
\choice
{\True Máy biến áp hoạt động với dòng điện xoay chiều dựa trên cảm ứng điện từ.}
{Máy biến áp làm việc bình thường với nguồn một chiều không đổi.}
{Máy tăng áp luôn có số vòng cuộn thứ cấp nhỏ hơn cuộn sơ cấp.}
{Hao phí truyền tải tỉ lệ nghịch với bình phương điện trở dây.}
\loigiai{
Máy biến áp hoạt động với dòng điện xoay chiều dựa trên cảm ứng điện từ. Do đó chọn A. Máy biến áp lí tưởng có U2/U1 = N2/N1 và U1I1 = U2I2. Truyền tải cùng công suất ở điện áp cao làm giảm hao phí I^2R.
}
\end{ex}

% ===== Câu 122 =====
% ID: L12C3B20-26-DS
% Mức: TH
\dangbt{Tổng hợp đường sức từ và xác định chiều từ trường}
\begin{ex}
% ID: L12C3B20-26-DS
% Mức: TH
Xét bài toán thuộc dạng tổng hợp đường sức từ và xác định chiều từ trường. Hãy xác định tính đúng, sai của bốn phát biểu sau.
\choiceTF
{Các đường sức từ có thể cắt nhau tại một điểm.}
{Đường sức từ luôn là những đường thẳng song song.}
{\True Chiều từ trường của dòng điện được xác định bằng quy tắc nắm tay phải.}
{\True Tại một điểm, hướng của vectơ cảm ứng từ trùng với hướng mà cực Bắc của kim nam châm chỉ.}
\loigiai{
\begin{itemchoice}
\item (Sai) Các đường sức từ có thể cắt nhau tại một điểm. (Vì): Các đường sức từ có thể cắt nhau tại một điểm. b) Sai: Đường sức từ luôn là những đường thẳng song song. c) Đúng: Chiều từ trường của dòng điện được xác định bằng quy tắc nắm tay phải. d) Đúng: Tại một điểm, hướng của vectơ cảm ứng từ trùng với hướng mà cực Bắc của kim nam châm chỉ.. Đường sức từ không cắt nhau; chiều từ trường được xác định bằng kim nam châm hoặc quy tắc nắm tay phải. Bên ngoài nam châm, đường sức đi từ cực Bắc đến cực Nam
\item (Sai) Đường sức từ luôn là những đường thẳng song song.
\item (Đúng) Chiều từ trường của dòng điện được xác định bằng quy tắc nắm tay phải.
\item (Đúng) Tại một điểm, hướng của vectơ cảm ứng từ trùng với hướng mà cực Bắc của kim nam châm chỉ.
\end{itemchoice}
}
\end{ex}

% ===== Câu 123 =====
% ID: L12C3B20-26-TL
% Mức: TH
\dangbt{Tổng hợp từ thông, định luật Lenz và suất điện động cảm ứng}
\begin{bt}
% ID: L12C3B20-26-TL
% Mức: TH
Phân tích các đại lượng, phương chiều và điều kiện áp dụng trong dạng tổng hợp từ thông, định luật lenz và suất điện động cảm ứng.
\loigiai{
Phi = BS cos(alpha). Suất điện động cảm ứng có độ lớn |e| = |Delta Phi/Delta t|; chiều dòng điện cảm ứng tuân theo định luật Lenz. Cần xác định rõ đại lượng vô hướng hay vectơ, quy ước chiều và điều kiện để công thức được áp dụng.
}
\end{bt}

% ===== Câu 124 =====
% ID: L12C3B20-26-TLN
% Mức: TH
\begin{ex}
% ID: L12C3B20-26-TLN
% Mức: TH
Trong bốn phát biểu sau về tổng hợp từ thông, định luật lenz và suất điện động cảm ứng, có bao nhiêu phát biểu đúng? Chỉ ghi một số nguyên. (1) Độ lớn suất điện động cảm ứng bằng độ lớn tốc độ biến thiên từ thông. (2) Từ thông không đổi thì không xuất hiện suất điện động cảm ứng. (3) Suất điện động cảm ứng chỉ phụ thuộc điện trở của mạch. (4) Khung dây đứng yên trong từ trường đều không đổi luôn có dòng điện cảm ứng.
\shortans{2}
\loigiai{
Đối chiếu từng phát biểu với kiến thức: Phi = BS cos(alpha). Suất điện động cảm ứng có độ lớn |e| = |Delta Phi/Delta t|; chiều dòng điện cảm ứng tuân theo định luật Lenz. Có 2 phát biểu đúng.
}
\end{ex}

% ===== Câu 125 =====
% ID: L12C3B20-26-TN
% Mức: TH
\dangbt{Tổng hợp máy biến áp, truyền tải điện năng và ứng dụng cảm ứng điện từ}
\begin{ex}
% ID: L12C3B20-26-TN
% Mức: TH
Phát biểu nào sau đây đúng về tổng hợp máy biến áp, truyền tải điện năng và ứng dụng cảm ứng điện từ?
\choice
{Máy tăng áp luôn có số vòng cuộn thứ cấp nhỏ hơn cuộn sơ cấp.}
{\True Máy biến áp lí tưởng thỏa U2/U1 = N2/N1.}
{Hao phí truyền tải tỉ lệ nghịch với bình phương điện trở dây.}
{Lõi máy biến áp ghép từ một khối thép đặc để tăng dòng điện Foucault.}
\loigiai{
Máy biến áp lí tưởng thỏa U2/U1 = N2/N1. Do đó chọn B. Máy biến áp lí tưởng có U2/U1 = N2/N1 và U1I1 = U2I2. Truyền tải cùng công suất ở điện áp cao làm giảm hao phí I^2R.
}
\end{ex}

% ===== Câu 126 =====
% ID: L12C3B20-27-DS
% Mức: TH
\dangbt{Tổng hợp đường sức từ và xác định chiều từ trường}
\begin{ex}
% ID: L12C3B20-27-DS
% Mức: TH
Xét bài toán thuộc dạng tổng hợp đường sức từ và xác định chiều từ trường. Hãy xác định tính đúng, sai của bốn phát biểu sau.
\choiceTF
{\True Đường sức từ của dòng điện thẳng là các đường tròn đồng tâm nằm trong mặt phẳng vuông góc với dây dẫn.}
{Đổi chiều dòng điện không làm đổi chiều từ trường.}
{\True Tại một điểm, hướng của vectơ cảm ứng từ trùng với hướng mà cực Bắc của kim nam châm chỉ.}
{Các đường sức từ có thể cắt nhau tại một điểm.}
\loigiai{
\begin{itemchoice}
\item (Đúng) Đường sức từ của dòng điện thẳng là các đường tròn đồng tâm nằm trong mặt phẳng vuông góc với dây dẫn. (Vì): ường sức từ của dòng điện thẳng là các đường tròn đồng tâm nằm trong mặt phẳng vuông góc với dây dẫn. b) Sai: Đổi chiều dòng điện không làm đổi chiều từ trường. c) Đúng: Tại một điểm, hướng của vectơ cảm ứng từ trùng với hướng mà cực Bắc của kim nam châm chỉ. d) Sai: Các đường sức từ có thể cắt nhau tại một điểm.. Đường sức từ không cắt nhau; chiều từ trường được xác định bằng kim nam châm hoặc quy tắc nắm tay phải. Bên ngoài nam châm, đường sức đi từ cực Bắc đến cực Nam
\item (Sai) Đổi chiều dòng điện không làm đổi chiều từ trường.
\item (Đúng) Tại một điểm, hướng của vectơ cảm ứng từ trùng với hướng mà cực Bắc của kim nam châm chỉ.
\item (Sai) Các đường sức từ có thể cắt nhau tại một điểm.
\end{itemchoice}
}
\end{ex}

% ===== Câu 127 =====
% ID: L12C3B20-27-TL
% Mức: VD
\dangbt{Tổng hợp từ thông, định luật Lenz và suất điện động cảm ứng}
\begin{bt}
% ID: L12C3B20-27-TL
% Mức: VD
Nêu một sai lầm thường gặp khi giải tổng hợp từ thông, định luật lenz và suất điện động cảm ứng và cách khắc phục.
\loigiai{
Phi = BS cos(alpha). Suất điện động cảm ứng có độ lớn |e| = |Delta Phi/Delta t|; chiều dòng điện cảm ứng tuân theo định luật Lenz. Sai lầm thường gặp là bỏ qua phương chiều, góc hoặc đơn vị. Khắc phục bằng hình vẽ, quy ước dấu và đổi đơn vị trước khi tính.
}
\end{bt}

% ===== Câu 128 =====
% ID: L12C3B20-27-TLN
% Mức: VD
\begin{ex}
% ID: L12C3B20-27-TLN
% Mức: VD
Trong bốn phát biểu sau về tổng hợp từ thông, định luật lenz và suất điện động cảm ứng, có bao nhiêu phát biểu đúng? Chỉ ghi một số nguyên. (1) Từ thông qua khung phẳng là Phi = BS cos(alpha), với alpha là góc giữa $B$ và pháp tuyến. (2) Độ lớn suất điện động cảm ứng bằng độ lớn tốc độ biến thiên từ thông. (3) Dòng điện cảm ứng luôn làm tăng từ thông ban đầu. (4) Khung dây đứng yên trong từ trường đều không đổi luôn có dòng điện cảm ứng.
\shortans{2}
\loigiai{
Đối chiếu từng phát biểu với kiến thức: Phi = BS cos(alpha). Suất điện động cảm ứng có độ lớn |e| = |Delta Phi/Delta t|; chiều dòng điện cảm ứng tuân theo định luật Lenz. Có 2 phát biểu đúng.
}
\end{ex}

% ===== Câu 129 =====
% ID: L12C3B20-27-TN
% Mức: TH
\dangbt{Tổng hợp máy biến áp, truyền tải điện năng và ứng dụng cảm ứng điện từ}
\begin{ex}
% ID: L12C3B20-27-TN
% Mức: TH
Trong Chương III, kết luận nào đúng khi giải bài thuộc dạng tổng hợp máy biến áp, truyền tải điện năng và ứng dụng cảm ứng điện từ?
\choice
{Hao phí truyền tải tỉ lệ nghịch với bình phương điện trở dây.}
{Lõi máy biến áp ghép từ một khối thép đặc để tăng dòng điện Foucault.}
{\True Tăng điện áp truyền tải với cùng công suất làm giảm hao phí trên đường dây.}
{Máy biến áp làm việc bình thường với nguồn một chiều không đổi.}
\loigiai{
Tăng điện áp truyền tải với cùng công suất làm giảm hao phí trên đường dây. Do đó chọn C. Máy biến áp lí tưởng có U2/U1 = N2/N1 và U1I1 = U2I2. Truyền tải cùng công suất ở điện áp cao làm giảm hao phí I^2R.
}
\end{ex}

% ===== Câu 130 =====
% ID: L12C3B20-28-DS
% Mức: VD
\dangbt{Tổng hợp đường sức từ và xác định chiều từ trường}
\begin{ex}
% ID: L12C3B20-28-DS
% Mức: VD
Xét bài toán thuộc dạng tổng hợp đường sức từ và xác định chiều từ trường. Hãy xác định tính đúng, sai của bốn phát biểu sau.
\choiceTF
{Đổi chiều dòng điện không làm đổi chiều từ trường.}
{\True Tại một điểm, hướng của vectơ cảm ứng từ trùng với hướng mà cực Bắc của kim nam châm chỉ.}
{Các đường sức từ có thể cắt nhau tại một điểm.}
{\True Đường sức từ của dòng điện thẳng là các đường tròn đồng tâm nằm trong mặt phẳng vuông góc với dây dẫn.}
\loigiai{
\begin{itemchoice}
\item (Sai) Đổi chiều dòng điện không làm đổi chiều từ trường. (Vì): Đổi chiều dòng điện không làm đổi chiều từ trường. b) Đúng: Tại một điểm, hướng của vectơ cảm ứng từ trùng với hướng mà cực Bắc của kim nam châm chỉ. c) Sai: Các đường sức từ có thể cắt nhau tại một điểm. d) Đúng: Đường sức từ của dòng điện thẳng là các đường tròn đồng tâm nằm trong mặt phẳng vuông góc với dây dẫn.. Đường sức từ không cắt nhau; chiều từ trường được xác định bằng kim nam châm hoặc quy tắc nắm tay phải. Bên ngoài nam châm, đường sức đi từ cực Bắc đến cực Nam
\item (Đúng) Tại một điểm, hướng của vectơ cảm ứng từ trùng với hướng mà cực Bắc của kim nam châm chỉ.
\item (Sai) Các đường sức từ có thể cắt nhau tại một điểm.
\item (Đúng) Đường sức từ của dòng điện thẳng là các đường tròn đồng tâm nằm trong mặt phẳng vuông góc với dây dẫn.
\end{itemchoice}
}
\end{ex}

% ===== Câu 131 =====
% ID: L12C3B20-28-TL
% Mức: VD
\dangbt{Tổng hợp từ thông, định luật Lenz và suất điện động cảm ứng}
\begin{bt}
% ID: L12C3B20-28-TL
% Mức: VD
Đề xuất quy trình kiểm tra kết quả của một bài toán tổng hợp từ thông, định luật lenz và suất điện động cảm ứng.
\loigiai{
Phi = BS cos(alpha). Suất điện động cảm ứng có độ lớn |e| = |Delta Phi/Delta t|; chiều dòng điện cảm ứng tuân theo định luật Lenz. Kiểm tra thứ nguyên, dấu, giới hạn đặc biệt và sự phù hợp với ý nghĩa vật lí.
}
\end{bt}

% ===== Câu 132 =====
% ID: L12C3B20-28-TLN
% Mức: VD
\begin{ex}
% ID: L12C3B20-28-TLN
% Mức: VD
Trong bốn phát biểu sau về tổng hợp từ thông, định luật lenz và suất điện động cảm ứng, có bao nhiêu phát biểu đúng? Chỉ ghi một số nguyên. (1) Từ thông qua khung phẳng là Phi = BS cos(alpha), với alpha là góc giữa $B$ và pháp tuyến. (2) Từ thông là đại lượng vectơ. (3) Dòng điện cảm ứng có chiều chống lại sự biến thiên từ thông sinh ra nó. (4) Dòng điện cảm ứng luôn làm tăng từ thông ban đầu.
\shortans{2}
\loigiai{
Đối chiếu từng phát biểu với kiến thức: Phi = BS cos(alpha). Suất điện động cảm ứng có độ lớn |e| = |Delta Phi/Delta t|; chiều dòng điện cảm ứng tuân theo định luật Lenz. Có 2 phát biểu đúng.
}
\end{ex}

% ===== Câu 133 =====
% ID: L12C3B20-28-TN
% Mức: VD
\dangbt{Tổng hợp máy biến áp, truyền tải điện năng và ứng dụng cảm ứng điện từ}
\begin{ex}
% ID: L12C3B20-28-TN
% Mức: VD
Chọn nhận định phù hợp với kiến thức Vật lí về tổng hợp máy biến áp, truyền tải điện năng và ứng dụng cảm ứng điện từ.
\choice
{Lõi máy biến áp ghép từ một khối thép đặc để tăng dòng điện Foucault.}
{Máy biến áp làm việc bình thường với nguồn một chiều không đổi.}
{Máy tăng áp luôn có số vòng cuộn thứ cấp nhỏ hơn cuộn sơ cấp.}
{\True Dòng điện Foucault được ứng dụng trong bếp từ và phanh điện từ.}
\loigiai{
Dòng điện Foucault được ứng dụng trong bếp từ và phanh điện từ. Do đó chọn D. Máy biến áp lí tưởng có U2/U1 = N2/N1 và U1I1 = U2I2. Truyền tải cùng công suất ở điện áp cao làm giảm hao phí I^2R.
}
\end{ex}

% ===== Câu 134 =====
% ID: L12C3B20-29-DS
% Mức: VD
\dangbt{Tổng hợp đường sức từ và xác định chiều từ trường}
\begin{ex}
% ID: L12C3B20-29-DS
% Mức: VD
Xét bài toán thuộc dạng tổng hợp đường sức từ và xác định chiều từ trường. Hãy xác định tính đúng, sai của bốn phát biểu sau.
\choiceTF
{\True Tại một điểm, hướng của vectơ cảm ứng từ trùng với hướng mà cực Bắc của kim nam châm chỉ.}
{\True Bên ngoài nam châm, đường sức từ có chiều đi ra từ cực Bắc và đi vào cực Nam.}
{Đường sức từ luôn là những đường thẳng song song.}
{Đổi chiều dòng điện không làm đổi chiều từ trường.}
\loigiai{
\begin{itemchoice}
\item (Đúng) Tại một điểm, hướng của vectơ cảm ứng từ trùng với hướng mà cực Bắc của kim nam châm chỉ. (Vì): Tại một điểm, hướng của vectơ cảm ứng từ trùng với hướng mà cực Bắc của kim nam châm chỉ. b) Đúng: Bên ngoài nam châm, đường sức từ có chiều đi ra từ cực Bắc và đi vào cực Nam. c) Sai: Đường sức từ luôn là những đường thẳng song song. d) Sai: Đổi chiều dòng điện không làm đổi chiều từ trường.. Đường sức từ không cắt nhau; chiều từ trường được xác định bằng kim nam châm hoặc quy tắc nắm tay phải. Bên ngoài nam châm, đường sức đi từ cực Bắc đến cực Nam
\item (Đúng) Bên ngoài nam châm, đường sức từ có chiều đi ra từ cực Bắc và đi vào cực Nam.
\item (Sai) Đường sức từ luôn là những đường thẳng song song.
\item (Sai) Đổi chiều dòng điện không làm đổi chiều từ trường.
\end{itemchoice}
}
\end{ex}

% ===== Câu 135 =====
% ID: L12C3B20-29-TL
% Mức: VD
\dangbt{Tổng hợp từ thông, định luật Lenz và suất điện động cảm ứng}
\begin{bt}
% ID: L12C3B20-29-TL
% Mức: VD
Xây dựng một tình huống thực tiễn thuộc dạng tổng hợp từ thông, định luật lenz và suất điện động cảm ứng và chỉ ra định luật cần dùng.
\loigiai{
Phi = BS cos(alpha). Suất điện động cảm ứng có độ lớn |e| = |Delta Phi/Delta t|; chiều dòng điện cảm ứng tuân theo định luật Lenz. Có thể chọn thiết bị hoặc hiện tượng thực tế tương ứng, nêu dữ kiện và chỉ rõ định luật dùng để giải thích hoặc tính toán.
}
\end{bt}

% ===== Câu 136 =====
% ID: L12C3B20-29-TLN
% Mức: VD
\begin{ex}
% ID: L12C3B20-29-TLN
% Mức: VD
Trong bốn phát biểu sau về tổng hợp từ thông, định luật lenz và suất điện động cảm ứng, có bao nhiêu phát biểu đúng? Chỉ ghi một số nguyên. (1) Độ lớn suất điện động cảm ứng bằng độ lớn tốc độ biến thiên từ thông. (2) Từ thông không đổi thì không xuất hiện suất điện động cảm ứng. (3) Suất điện động cảm ứng chỉ phụ thuộc điện trở của mạch. (4) Khung dây đứng yên trong từ trường đều không đổi luôn có dòng điện cảm ứng.
\shortans{2}
\loigiai{
Đối chiếu từng phát biểu với kiến thức: Phi = BS cos(alpha). Suất điện động cảm ứng có độ lớn |e| = |Delta Phi/Delta t|; chiều dòng điện cảm ứng tuân theo định luật Lenz. Có 2 phát biểu đúng.
}
\end{ex}

% ===== Câu 137 =====
% ID: L12C3B20-29-TN
% Mức: VD
\dangbt{Tổng hợp máy biến áp, truyền tải điện năng và ứng dụng cảm ứng điện từ}
\begin{ex}
% ID: L12C3B20-29-TN
% Mức: VD
Khi phân tích tổng hợp máy biến áp, truyền tải điện năng và ứng dụng cảm ứng điện từ, phát biểu nào đúng?
\choice
{\True Máy biến áp hoạt động với dòng điện xoay chiều dựa trên cảm ứng điện từ.}
{Máy biến áp làm việc bình thường với nguồn một chiều không đổi.}
{Máy tăng áp luôn có số vòng cuộn thứ cấp nhỏ hơn cuộn sơ cấp.}
{Hao phí truyền tải tỉ lệ nghịch với bình phương điện trở dây.}
\loigiai{
Máy biến áp hoạt động với dòng điện xoay chiều dựa trên cảm ứng điện từ. Do đó chọn A. Máy biến áp lí tưởng có U2/U1 = N2/N1 và U1I1 = U2I2. Truyền tải cùng công suất ở điện áp cao làm giảm hao phí I^2R.
}
\end{ex}

% ===== Câu 138 =====
% ID: L12C3B20-30-DS
% Mức: VD
\dangbt{Tổng hợp đường sức từ và xác định chiều từ trường}
\begin{ex}
% ID: L12C3B20-30-DS
% Mức: VD
Xét bài toán thuộc dạng tổng hợp đường sức từ và xác định chiều từ trường. Hãy xác định tính đúng, sai của bốn phát biểu sau.
\choiceTF
{Các đường sức từ có thể cắt nhau tại một điểm.}
{\True Đường sức từ của dòng điện thẳng là các đường tròn đồng tâm nằm trong mặt phẳng vuông góc với dây dẫn.}
{\True Chiều từ trường của dòng điện được xác định bằng quy tắc nắm tay phải.}
{Từ trường đều có các đường sức cong và mật độ thay đổi.}
\loigiai{
\begin{itemchoice}
\item (Sai) Các đường sức từ có thể cắt nhau tại một điểm. (Vì): Các đường sức từ có thể cắt nhau tại một điểm. b) Đúng: Đường sức từ của dòng điện thẳng là các đường tròn đồng tâm nằm trong mặt phẳng vuông góc với dây dẫn. c) Đúng: Chiều từ trường của dòng điện được xác định bằng quy tắc nắm tay phải. d) Sai: Từ trường đều có các đường sức cong và mật độ thay đổi.. Đường sức từ không cắt nhau; chiều từ trường được xác định bằng kim nam châm hoặc quy tắc nắm tay phải. Bên ngoài nam châm, đường sức đi từ cực Bắc đến cực Nam
\item (Đúng) Đường sức từ của dòng điện thẳng là các đường tròn đồng tâm nằm trong mặt phẳng vuông góc với dây dẫn.
\item (Đúng) Chiều từ trường của dòng điện được xác định bằng quy tắc nắm tay phải.
\item (Sai) Từ trường đều có các đường sức cong và mật độ thay đổi.
\end{itemchoice}
}
\end{ex}

% ===== Câu 139 =====
% ID: L12C3B20-30-TL
% Mức: VDC
\dangbt{Tổng hợp từ thông, định luật Lenz và suất điện động cảm ứng}
\begin{bt}
% ID: L12C3B20-30-TL
% Mức: VDC
So sánh một nhận định đúng và một nhận định sai trong dạng tổng hợp từ thông, định luật lenz và suất điện động cảm ứng; giải thích căn cứ Vật lí.
\loigiai{
Nhận định đúng: Từ thông qua khung phẳng là Phi = BS cos(alpha), với alpha là góc giữa $B$ và pháp tuyến. Nhận định sai: Từ thông là đại lượng vectơ. Căn cứ: Phi = BS cos(alpha). Suất điện động cảm ứng có độ lớn |e| = |Delta Phi/Delta t|; chiều dòng điện cảm ứng tuân theo định luật Lenz.
}
\end{bt}

% ===== Câu 140 =====
% ID: L12C3B20-30-TLN
% Mức: VDC
\begin{ex}
% ID: L12C3B20-30-TLN
% Mức: VDC
Trong bốn phát biểu sau về tổng hợp từ thông, định luật lenz và suất điện động cảm ứng, có bao nhiêu phát biểu đúng? Chỉ ghi một số nguyên. (1) Từ thông qua khung phẳng là Phi = BS cos(alpha), với alpha là góc giữa $B$ và pháp tuyến. (2) Độ lớn suất điện động cảm ứng bằng độ lớn tốc độ biến thiên từ thông. (3) Dòng điện cảm ứng luôn làm tăng từ thông ban đầu. (4) Khung dây đứng yên trong từ trường đều không đổi luôn có dòng điện cảm ứng.
\shortans{2}
\loigiai{
Đối chiếu từng phát biểu với kiến thức: Phi = BS cos(alpha). Suất điện động cảm ứng có độ lớn |e| = |Delta Phi/Delta t|; chiều dòng điện cảm ứng tuân theo định luật Lenz. Có 2 phát biểu đúng.
}
\end{ex}

% ===== Câu 141 =====
% ID: L12C3B20-30-TN
% Mức: VD
\dangbt{Tổng hợp máy biến áp, truyền tải điện năng và ứng dụng cảm ứng điện từ}
\begin{ex}
% ID: L12C3B20-30-TN
% Mức: VD
Cơ sở Vật lí nào sau đây cần được sử dụng trong dạng tổng hợp máy biến áp, truyền tải điện năng và ứng dụng cảm ứng điện từ?
\choice
{Máy tăng áp luôn có số vòng cuộn thứ cấp nhỏ hơn cuộn sơ cấp.}
{\True Máy biến áp lí tưởng thỏa U2/U1 = N2/N1.}
{Hao phí truyền tải tỉ lệ nghịch với bình phương điện trở dây.}
{Lõi máy biến áp ghép từ một khối thép đặc để tăng dòng điện Foucault.}
\loigiai{
Máy biến áp lí tưởng thỏa U2/U1 = N2/N1. Do đó chọn B. Máy biến áp lí tưởng có U2/U1 = N2/N1 và U1I1 = U2I2. Truyền tải cùng công suất ở điện áp cao làm giảm hao phí I^2R.
}
\end{ex}

% ===== Câu 142 =====
% ID: L12C3B20-31-DS
% Mức: TH
\dangbt{Tổng hợp đường sức từ và xác định chiều từ trường}
\begin{ex}
% ID: L12C3B20-31-DS
% Mức: TH
Xét bài toán thuộc dạng tổng hợp đường sức từ và xác định chiều từ trường. Hãy xác định tính đúng, sai của bốn phát biểu sau.
\choiceTF
{Các đường sức từ có thể cắt nhau tại một điểm.}
{Đường sức từ luôn là những đường thẳng song song.}
{\True Chiều từ trường của dòng điện được xác định bằng quy tắc nắm tay phải.}
{\True Tại một điểm, hướng của vectơ cảm ứng từ trùng với hướng mà cực Bắc của kim nam châm chỉ.}
\loigiai{
\begin{itemchoice}
\item (Sai) Các đường sức từ có thể cắt nhau tại một điểm. (Vì): Các đường sức từ có thể cắt nhau tại một điểm. b) Sai: Đường sức từ luôn là những đường thẳng song song. c) Đúng: Chiều từ trường của dòng điện được xác định bằng quy tắc nắm tay phải. d) Đúng: Tại một điểm, hướng của vectơ cảm ứng từ trùng với hướng mà cực Bắc của kim nam châm chỉ.. Đường sức từ không cắt nhau; chiều từ trường được xác định bằng kim nam châm hoặc quy tắc nắm tay phải. Bên ngoài nam châm, đường sức đi từ cực Bắc đến cực Nam
\item (Sai) Đường sức từ luôn là những đường thẳng song song.
\item (Đúng) Chiều từ trường của dòng điện được xác định bằng quy tắc nắm tay phải.
\item (Đúng) Tại một điểm, hướng của vectơ cảm ứng từ trùng với hướng mà cực Bắc của kim nam châm chỉ.
\end{itemchoice}
}
\end{ex}

% ===== Câu 143 =====
% ID: L12C3B20-31-TL
% Mức: TH
\begin{bt}
% ID: L12C3B20-31-TL
% Mức: TH
Trình bày kiến thức cốt lõi và các bước giải bài thuộc dạng tổng hợp đường sức từ và xác định chiều từ trường.
\loigiai{
Đường sức từ không cắt nhau; chiều từ trường được xác định bằng kim nam châm hoặc quy tắc nắm tay phải. Bên ngoài nam châm, đường sức đi từ cực Bắc đến cực Nam. Các bước: tóm tắt dữ kiện; xác định phương, chiều và điều kiện; chọn hệ thức; đổi đơn vị; tính toán; kiểm tra kết quả.
}
\end{bt}

% ===== Câu 144 =====
% ID: L12C3B20-31-TLN
% Mức: TH
\begin{ex}
% ID: L12C3B20-31-TLN
% Mức: TH
Trong bốn phát biểu sau về tổng hợp đường sức từ và xác định chiều từ trường, có bao nhiêu phát biểu đúng? Chỉ ghi một số nguyên. (1) Bên ngoài nam châm, đường sức từ có chiều đi ra từ cực Bắc và đi vào cực Nam. (2) Các đường sức từ có thể cắt nhau tại một điểm. (3) Đường sức từ của dòng điện thẳng là các đường tròn đồng tâm nằm trong mặt phẳng vuông góc với dây dẫn. (4) Đường sức từ luôn là những đường thẳng song song.
\shortans{2}
\loigiai{
Đối chiếu từng phát biểu với kiến thức: Đường sức từ không cắt nhau; chiều từ trường được xác định bằng kim nam châm hoặc quy tắc nắm tay phải. Bên ngoài nam châm, đường sức đi từ cực Bắc đến cực Nam. Có 2 phát biểu đúng.
}
\end{ex}

% ===== Câu 145 =====
% ID: L12C3B20-31-TN
% Mức: NB
\dangbt{Tổng hợp máy phát điện xoay chiều và đồ thị điện từ}
\begin{ex}
% ID: L12C3B20-31-TN
% Mức: NB
Phát biểu nào sau đây đúng về tổng hợp máy phát điện xoay chiều và đồ thị điện từ?
\choice
{\True Máy phát điện xoay chiều hoạt động dựa trên hiện tượng cảm ứng điện từ.}
{Máy phát điện biến điện năng thành cơ năng.}
{Tăng tốc độ quay làm giảm biên độ suất điện động.}
{Từ thông và suất điện động luôn cùng pha.}
\loigiai{
Máy phát điện xoay chiều hoạt động dựa trên hiện tượng cảm ứng điện từ. Do đó chọn A. Máy phát điện biến cơ năng thành điện năng nhờ cảm ứng điện từ; nếu Phi = Phi0 cos(omega t + phi) thì e = E0 sin(omega t + phi), E0 = NBS omega.
}
\end{ex}

% ===== Câu 146 =====
% ID: L12C3B20-32-DS
% Mức: TH
\dangbt{Tổng hợp đường sức từ và xác định chiều từ trường}
\begin{ex}
% ID: L12C3B20-32-DS
% Mức: TH
Xét bài toán thuộc dạng tổng hợp đường sức từ và xác định chiều từ trường. Hãy xác định tính đúng, sai của bốn phát biểu sau.
\choiceTF
{\True Đường sức từ của dòng điện thẳng là các đường tròn đồng tâm nằm trong mặt phẳng vuông góc với dây dẫn.}
{Đổi chiều dòng điện không làm đổi chiều từ trường.}
{\True Tại một điểm, hướng của vectơ cảm ứng từ trùng với hướng mà cực Bắc của kim nam châm chỉ.}
{Các đường sức từ có thể cắt nhau tại một điểm.}
\loigiai{
\begin{itemchoice}
\item (Đúng) Đường sức từ của dòng điện thẳng là các đường tròn đồng tâm nằm trong mặt phẳng vuông góc với dây dẫn. (Vì): ường sức từ của dòng điện thẳng là các đường tròn đồng tâm nằm trong mặt phẳng vuông góc với dây dẫn. b) Sai: Đổi chiều dòng điện không làm đổi chiều từ trường. c) Đúng: Tại một điểm, hướng của vectơ cảm ứng từ trùng với hướng mà cực Bắc của kim nam châm chỉ. d) Sai: Các đường sức từ có thể cắt nhau tại một điểm.. Đường sức từ không cắt nhau; chiều từ trường được xác định bằng kim nam châm hoặc quy tắc nắm tay phải. Bên ngoài nam châm, đường sức đi từ cực Bắc đến cực Nam
\item (Sai) Đổi chiều dòng điện không làm đổi chiều từ trường.
\item (Đúng) Tại một điểm, hướng của vectơ cảm ứng từ trùng với hướng mà cực Bắc của kim nam châm chỉ.
\item (Sai) Các đường sức từ có thể cắt nhau tại một điểm.
\end{itemchoice}
}
\end{ex}

% ===== Câu 147 =====
% ID: L12C3B20-32-TL
% Mức: TH
\begin{bt}
% ID: L12C3B20-32-TL
% Mức: TH
Phân tích các đại lượng, phương chiều và điều kiện áp dụng trong dạng tổng hợp đường sức từ và xác định chiều từ trường.
\loigiai{
Đường sức từ không cắt nhau; chiều từ trường được xác định bằng kim nam châm hoặc quy tắc nắm tay phải. Bên ngoài nam châm, đường sức đi từ cực Bắc đến cực Nam. Cần xác định rõ đại lượng vô hướng hay vectơ, quy ước chiều và điều kiện để công thức được áp dụng.
}
\end{bt}

% ===== Câu 148 =====
% ID: L12C3B20-32-TLN
% Mức: TH
\begin{ex}
% ID: L12C3B20-32-TLN
% Mức: TH
Trong bốn phát biểu sau về tổng hợp đường sức từ và xác định chiều từ trường, có bao nhiêu phát biểu đúng? Chỉ ghi một số nguyên. (1) Chiều từ trường của dòng điện được xác định bằng quy tắc nắm tay phải. (2) Tại một điểm, hướng của vectơ cảm ứng từ trùng với hướng mà cực Bắc của kim nam châm chỉ. (3) Đổi chiều dòng điện không làm đổi chiều từ trường. (4) Từ trường đều có các đường sức cong và mật độ thay đổi.
\shortans{2}
\loigiai{
Đối chiếu từng phát biểu với kiến thức: Đường sức từ không cắt nhau; chiều từ trường được xác định bằng kim nam châm hoặc quy tắc nắm tay phải. Bên ngoài nam châm, đường sức đi từ cực Bắc đến cực Nam. Có 2 phát biểu đúng.
}
\end{ex}

% ===== Câu 149 =====
% ID: L12C3B20-32-TN
% Mức: NB
\dangbt{Tổng hợp máy phát điện xoay chiều và đồ thị điện từ}
\begin{ex}
% ID: L12C3B20-32-TN
% Mức: NB
Trong Chương III, kết luận nào đúng khi giải bài thuộc dạng tổng hợp máy phát điện xoay chiều và đồ thị điện từ?
\choice
{Tăng tốc độ quay làm giảm biên độ suất điện động.}
{\True Khung $N$ vòng quay đều có suất điện động cực đại E0 = NBS omega.}
{Từ thông và suất điện động luôn cùng pha.}
{Suất điện động cực đại không phụ thuộc số vòng dây.}
\loigiai{
Khung $N$ vòng quay đều có suất điện động cực đại E0 = NBS omega. Do đó chọn B. Máy phát điện biến cơ năng thành điện năng nhờ cảm ứng điện từ; nếu Phi = Phi0 cos(omega t + phi) thì e = E0 sin(omega t + phi), E0 = NBS omega.
}
\end{ex}

% ===== Câu 150 =====
% ID: L12C3B20-33-DS
% Mức: VD
\dangbt{Tổng hợp đường sức từ và xác định chiều từ trường}
\begin{ex}
% ID: L12C3B20-33-DS
% Mức: VD
Xét bài toán thuộc dạng tổng hợp đường sức từ và xác định chiều từ trường. Hãy xác định tính đúng, sai của bốn phát biểu sau.
\choiceTF
{Đổi chiều dòng điện không làm đổi chiều từ trường.}
{\True Tại một điểm, hướng của vectơ cảm ứng từ trùng với hướng mà cực Bắc của kim nam châm chỉ.}
{Các đường sức từ có thể cắt nhau tại một điểm.}
{\True Đường sức từ của dòng điện thẳng là các đường tròn đồng tâm nằm trong mặt phẳng vuông góc với dây dẫn.}
\loigiai{
\begin{itemchoice}
\item (Sai) Đổi chiều dòng điện không làm đổi chiều từ trường. (Vì): Đổi chiều dòng điện không làm đổi chiều từ trường. b) Đúng: Tại một điểm, hướng của vectơ cảm ứng từ trùng với hướng mà cực Bắc của kim nam châm chỉ. c) Sai: Các đường sức từ có thể cắt nhau tại một điểm. d) Đúng: Đường sức từ của dòng điện thẳng là các đường tròn đồng tâm nằm trong mặt phẳng vuông góc với dây dẫn.. Đường sức từ không cắt nhau; chiều từ trường được xác định bằng kim nam châm hoặc quy tắc nắm tay phải. Bên ngoài nam châm, đường sức đi từ cực Bắc đến cực Nam
\item (Đúng) Tại một điểm, hướng của vectơ cảm ứng từ trùng với hướng mà cực Bắc của kim nam châm chỉ.
\item (Sai) Các đường sức từ có thể cắt nhau tại một điểm.
\item (Đúng) Đường sức từ của dòng điện thẳng là các đường tròn đồng tâm nằm trong mặt phẳng vuông góc với dây dẫn.
\end{itemchoice}
}
\end{ex}

% ===== Câu 151 =====
% ID: L12C3B20-33-TL
% Mức: VD
\begin{bt}
% ID: L12C3B20-33-TL
% Mức: VD
Nêu một sai lầm thường gặp khi giải tổng hợp đường sức từ và xác định chiều từ trường và cách khắc phục.
\loigiai{
Đường sức từ không cắt nhau; chiều từ trường được xác định bằng kim nam châm hoặc quy tắc nắm tay phải. Bên ngoài nam châm, đường sức đi từ cực Bắc đến cực Nam. Sai lầm thường gặp là bỏ qua phương chiều, góc hoặc đơn vị. Khắc phục bằng hình vẽ, quy ước dấu và đổi đơn vị trước khi tính.
}
\end{bt}

% ===== Câu 152 =====
% ID: L12C3B20-33-TLN
% Mức: VD
\begin{ex}
% ID: L12C3B20-33-TLN
% Mức: VD
Trong bốn phát biểu sau về tổng hợp đường sức từ và xác định chiều từ trường, có bao nhiêu phát biểu đúng? Chỉ ghi một số nguyên. (1) Bên ngoài nam châm, đường sức từ có chiều đi ra từ cực Bắc và đi vào cực Nam. (2) Chiều từ trường của dòng điện được xác định bằng quy tắc nắm tay phải. (3) Đường sức từ luôn là những đường thẳng song song. (4) Từ trường đều có các đường sức cong và mật độ thay đổi.
\shortans{2}
\loigiai{
Đối chiếu từng phát biểu với kiến thức: Đường sức từ không cắt nhau; chiều từ trường được xác định bằng kim nam châm hoặc quy tắc nắm tay phải. Bên ngoài nam châm, đường sức đi từ cực Bắc đến cực Nam. Có 2 phát biểu đúng.
}
\end{ex}

% ===== Câu 153 =====
% ID: L12C3B20-33-TN
% Mức: NB
\dangbt{Tổng hợp máy phát điện xoay chiều và đồ thị điện từ}
\begin{ex}
% ID: L12C3B20-33-TN
% Mức: NB
Chọn nhận định phù hợp với kiến thức Vật lí về tổng hợp máy phát điện xoay chiều và đồ thị điện từ.
\choice
{Từ thông và suất điện động luôn cùng pha.}
{Suất điện động cực đại không phụ thuộc số vòng dây.}
{\True Tần số suất điện động bằng tần số quay của khung trong mô hình một cặp cực.}
{Máy phát điện biến điện năng thành cơ năng.}
\loigiai{
Tần số suất điện động bằng tần số quay của khung trong mô hình một cặp cực. Do đó chọn C. Máy phát điện biến cơ năng thành điện năng nhờ cảm ứng điện từ; nếu Phi = Phi0 cos(omega t + phi) thì e = E0 sin(omega t + phi), E0 = NBS omega.
}
\end{ex}

% ===== Câu 154 =====
% ID: L12C3B20-34-DS
% Mức: VD
\dangbt{Tổng hợp đường sức từ và xác định chiều từ trường}
\begin{ex}
% ID: L12C3B20-34-DS
% Mức: VD
Xét bài toán thuộc dạng tổng hợp đường sức từ và xác định chiều từ trường. Hãy xác định tính đúng, sai của bốn phát biểu sau.
\choiceTF
{\True Tại một điểm, hướng của vectơ cảm ứng từ trùng với hướng mà cực Bắc của kim nam châm chỉ.}
{\True Bên ngoài nam châm, đường sức từ có chiều đi ra từ cực Bắc và đi vào cực Nam.}
{Đường sức từ luôn là những đường thẳng song song.}
{Đổi chiều dòng điện không làm đổi chiều từ trường.}
\loigiai{
\begin{itemchoice}
\item (Đúng) Tại một điểm, hướng của vectơ cảm ứng từ trùng với hướng mà cực Bắc của kim nam châm chỉ. (Vì): Tại một điểm, hướng của vectơ cảm ứng từ trùng với hướng mà cực Bắc của kim nam châm chỉ. b) Đúng: Bên ngoài nam châm, đường sức từ có chiều đi ra từ cực Bắc và đi vào cực Nam. c) Sai: Đường sức từ luôn là những đường thẳng song song. d) Sai: Đổi chiều dòng điện không làm đổi chiều từ trường.. Đường sức từ không cắt nhau; chiều từ trường được xác định bằng kim nam châm hoặc quy tắc nắm tay phải. Bên ngoài nam châm, đường sức đi từ cực Bắc đến cực Nam
\item (Đúng) Bên ngoài nam châm, đường sức từ có chiều đi ra từ cực Bắc và đi vào cực Nam.
\item (Sai) Đường sức từ luôn là những đường thẳng song song.
\item (Sai) Đổi chiều dòng điện không làm đổi chiều từ trường.
\end{itemchoice}
}
\end{ex}

% ===== Câu 155 =====
% ID: L12C3B20-34-TL
% Mức: VD
\begin{bt}
% ID: L12C3B20-34-TL
% Mức: VD
Đề xuất quy trình kiểm tra kết quả của một bài toán tổng hợp đường sức từ và xác định chiều từ trường.
\loigiai{
Đường sức từ không cắt nhau; chiều từ trường được xác định bằng kim nam châm hoặc quy tắc nắm tay phải. Bên ngoài nam châm, đường sức đi từ cực Bắc đến cực Nam. Kiểm tra thứ nguyên, dấu, giới hạn đặc biệt và sự phù hợp với ý nghĩa vật lí.
}
\end{bt}

% ===== Câu 156 =====
% ID: L12C3B20-34-TLN
% Mức: VD
\begin{ex}
% ID: L12C3B20-34-TLN
% Mức: VD
Trong bốn phát biểu sau về tổng hợp đường sức từ và xác định chiều từ trường, có bao nhiêu phát biểu đúng? Chỉ ghi một số nguyên. (1) Bên ngoài nam châm, đường sức từ có chiều đi ra từ cực Bắc và đi vào cực Nam. (2) Các đường sức từ có thể cắt nhau tại một điểm. (3) Đường sức từ của dòng điện thẳng là các đường tròn đồng tâm nằm trong mặt phẳng vuông góc với dây dẫn. (4) Đường sức từ luôn là những đường thẳng song song.
\shortans{2}
\loigiai{
Đối chiếu từng phát biểu với kiến thức: Đường sức từ không cắt nhau; chiều từ trường được xác định bằng kim nam châm hoặc quy tắc nắm tay phải. Bên ngoài nam châm, đường sức đi từ cực Bắc đến cực Nam. Có 2 phát biểu đúng.
}
\end{ex}

% ===== Câu 157 =====
% ID: L12C3B20-34-TN
% Mức: TH
\dangbt{Tổng hợp máy phát điện xoay chiều và đồ thị điện từ}
\begin{ex}
% ID: L12C3B20-34-TN
% Mức: TH
Khi phân tích tổng hợp máy phát điện xoay chiều và đồ thị điện từ, phát biểu nào đúng?
\choice
{Suất điện động cực đại không phụ thuộc số vòng dây.}
{Máy phát điện biến điện năng thành cơ năng.}
{Tăng tốc độ quay làm giảm biên độ suất điện động.}
{\True Suất điện động cảm ứng lệch pha pi/2 so với từ thông.}
\loigiai{
Suất điện động cảm ứng lệch pha pi/2 so với từ thông. Do đó chọn D. Máy phát điện biến cơ năng thành điện năng nhờ cảm ứng điện từ; nếu Phi = Phi0 cos(omega t + phi) thì e = E0 sin(omega t + phi), E0 = NBS omega.
}
\end{ex}

% ===== Câu 158 =====
% ID: L12C3B20-35-DS
% Mức: VD
\dangbt{Tổng hợp đường sức từ và xác định chiều từ trường}
\begin{ex}
% ID: L12C3B20-35-DS
% Mức: VD
Xét bài toán thuộc dạng tổng hợp đường sức từ và xác định chiều từ trường. Hãy xác định tính đúng, sai của bốn phát biểu sau.
\choiceTF
{Các đường sức từ có thể cắt nhau tại một điểm.}
{\True Đường sức từ của dòng điện thẳng là các đường tròn đồng tâm nằm trong mặt phẳng vuông góc với dây dẫn.}
{\True Chiều từ trường của dòng điện được xác định bằng quy tắc nắm tay phải.}
{Từ trường đều có các đường sức cong và mật độ thay đổi.}
\loigiai{
\begin{itemchoice}
\item (Sai) Các đường sức từ có thể cắt nhau tại một điểm. (Vì): Các đường sức từ có thể cắt nhau tại một điểm. b) Đúng: Đường sức từ của dòng điện thẳng là các đường tròn đồng tâm nằm trong mặt phẳng vuông góc với dây dẫn. c) Đúng: Chiều từ trường của dòng điện được xác định bằng quy tắc nắm tay phải. d) Sai: Từ trường đều có các đường sức cong và mật độ thay đổi.. Đường sức từ không cắt nhau; chiều từ trường được xác định bằng kim nam châm hoặc quy tắc nắm tay phải. Bên ngoài nam châm, đường sức đi từ cực Bắc đến cực Nam
\item (Đúng) Đường sức từ của dòng điện thẳng là các đường tròn đồng tâm nằm trong mặt phẳng vuông góc với dây dẫn.
\item (Đúng) Chiều từ trường của dòng điện được xác định bằng quy tắc nắm tay phải.
\item (Sai) Từ trường đều có các đường sức cong và mật độ thay đổi.
\end{itemchoice}
}
\end{ex}

% ===== Câu 159 =====
% ID: L12C3B20-35-TL
% Mức: VD
\begin{bt}
% ID: L12C3B20-35-TL
% Mức: VD
Xây dựng một tình huống thực tiễn thuộc dạng tổng hợp đường sức từ và xác định chiều từ trường và chỉ ra định luật cần dùng.
\loigiai{
Đường sức từ không cắt nhau; chiều từ trường được xác định bằng kim nam châm hoặc quy tắc nắm tay phải. Bên ngoài nam châm, đường sức đi từ cực Bắc đến cực Nam. Có thể chọn thiết bị hoặc hiện tượng thực tế tương ứng, nêu dữ kiện và chỉ rõ định luật dùng để giải thích hoặc tính toán.
}
\end{bt}

% ===== Câu 160 =====
% ID: L12C3B20-35-TLN
% Mức: VD
\begin{ex}
% ID: L12C3B20-35-TLN
% Mức: VD
Trong bốn phát biểu sau về tổng hợp đường sức từ và xác định chiều từ trường, có bao nhiêu phát biểu đúng? Chỉ ghi một số nguyên. (1) Chiều từ trường của dòng điện được xác định bằng quy tắc nắm tay phải. (2) Tại một điểm, hướng của vectơ cảm ứng từ trùng với hướng mà cực Bắc của kim nam châm chỉ. (3) Đổi chiều dòng điện không làm đổi chiều từ trường. (4) Từ trường đều có các đường sức cong và mật độ thay đổi.
\shortans{2}
\loigiai{
Đối chiếu từng phát biểu với kiến thức: Đường sức từ không cắt nhau; chiều từ trường được xác định bằng kim nam châm hoặc quy tắc nắm tay phải. Bên ngoài nam châm, đường sức đi từ cực Bắc đến cực Nam. Có 2 phát biểu đúng.
}
\end{ex}

% ===== Câu 161 =====
% ID: L12C3B20-35-TN
% Mức: TH
\dangbt{Tổng hợp máy phát điện xoay chiều và đồ thị điện từ}
\begin{ex}
% ID: L12C3B20-35-TN
% Mức: TH
Cơ sở Vật lí nào sau đây cần được sử dụng trong dạng tổng hợp máy phát điện xoay chiều và đồ thị điện từ?
\choice
{\True Máy phát điện xoay chiều hoạt động dựa trên hiện tượng cảm ứng điện từ.}
{Máy phát điện biến điện năng thành cơ năng.}
{Tăng tốc độ quay làm giảm biên độ suất điện động.}
{Từ thông và suất điện động luôn cùng pha.}
\loigiai{
Máy phát điện xoay chiều hoạt động dựa trên hiện tượng cảm ứng điện từ. Do đó chọn A. Máy phát điện biến cơ năng thành điện năng nhờ cảm ứng điện từ; nếu Phi = Phi0 cos(omega t + phi) thì e = E0 sin(omega t + phi), E0 = NBS omega.
}
\end{ex}

% ===== Câu 162 =====
% ID: L12C3B20-36-DS
% Mức: TH
\dangbt{Tổng hợp lực từ tác dụng lên dây dẫn mang dòng điện}
\begin{ex}
% ID: L12C3B20-36-DS
% Mức: TH
Xét bài toán thuộc dạng tổng hợp lực từ tác dụng lên dây dẫn mang dòng điện. Hãy xác định tính đúng, sai của bốn phát biểu sau.
\choiceTF
{Lực từ luôn cùng chiều với dòng điện.}
{Khi dây dẫn vuông góc với từ trường thì lực từ bằng 0.}
{\True Khi dây dẫn song song với đường sức từ thì lực từ bằng 0.}
{\True Đổi chiều dòng điện thì chiều lực từ đổi ngược.}
\loigiai{
\begin{itemchoice}
\item (Sai) Lực từ luôn cùng chiều với dòng điện. (Vì): lpha); chiều lực xác định bằng quy tắc bàn tay trái
\item (Sai) Khi dây dẫn vuông góc với từ trường thì lực từ bằng 0.
\item (Đúng) Khi dây dẫn song song với đường sức từ thì lực từ bằng 0.
\item (Đúng) Đổi chiều dòng điện thì chiều lực từ đổi ngược.
\end{itemchoice}
}
\end{ex}

% ===== Câu 163 =====
% ID: L12C3B20-36-TL
% Mức: VDC
\dangbt{Tổng hợp đường sức từ và xác định chiều từ trường}
\begin{bt}
% ID: L12C3B20-36-TL
% Mức: VDC
So sánh một nhận định đúng và một nhận định sai trong dạng tổng hợp đường sức từ và xác định chiều từ trường; giải thích căn cứ Vật lí.
\loigiai{
Nhận định đúng: Bên ngoài nam châm, đường sức từ có chiều đi ra từ cực Bắc và đi vào cực Nam. Nhận định sai: Các đường sức từ có thể cắt nhau tại một điểm. Căn cứ: Đường sức từ không cắt nhau; chiều từ trường được xác định bằng kim nam châm hoặc quy tắc nắm tay phải. Bên ngoài nam châm, đường sức đi từ cực Bắc đến cực Nam.
}
\end{bt}

% ===== Câu 164 =====
% ID: L12C3B20-36-TLN
% Mức: VDC
\begin{ex}
% ID: L12C3B20-36-TLN
% Mức: VDC
Trong bốn phát biểu sau về tổng hợp đường sức từ và xác định chiều từ trường, có bao nhiêu phát biểu đúng? Chỉ ghi một số nguyên. (1) Bên ngoài nam châm, đường sức từ có chiều đi ra từ cực Bắc và đi vào cực Nam. (2) Chiều từ trường của dòng điện được xác định bằng quy tắc nắm tay phải. (3) Đường sức từ luôn là những đường thẳng song song. (4) Từ trường đều có các đường sức cong và mật độ thay đổi.
\shortans{2}
\loigiai{
Đối chiếu từng phát biểu với kiến thức: Đường sức từ không cắt nhau; chiều từ trường được xác định bằng kim nam châm hoặc quy tắc nắm tay phải. Bên ngoài nam châm, đường sức đi từ cực Bắc đến cực Nam. Có 2 phát biểu đúng.
}
\end{ex}

% ===== Câu 165 =====
% ID: L12C3B20-36-TN
% Mức: TH
\dangbt{Tổng hợp máy phát điện xoay chiều và đồ thị điện từ}
\begin{ex}
% ID: L12C3B20-36-TN
% Mức: TH
Phát biểu nào sau đây đúng về tổng hợp máy phát điện xoay chiều và đồ thị điện từ?
\choice
{Tăng tốc độ quay làm giảm biên độ suất điện động.}
{\True Khung $N$ vòng quay đều có suất điện động cực đại E0 = NBS omega.}
{Từ thông và suất điện động luôn cùng pha.}
{Suất điện động cực đại không phụ thuộc số vòng dây.}
\loigiai{
Khung $N$ vòng quay đều có suất điện động cực đại E0 = NBS omega. Do đó chọn B. Máy phát điện biến cơ năng thành điện năng nhờ cảm ứng điện từ; nếu Phi = Phi0 cos(omega t + phi) thì e = E0 sin(omega t + phi), E0 = NBS omega.
}
\end{ex}

% ===== Câu 166 =====
% ID: L12C3B20-37-DS
% Mức: TH
\dangbt{Tổng hợp lực từ tác dụng lên dây dẫn mang dòng điện}
\begin{ex}
% ID: L12C3B20-37-DS
% Mức: TH
Xét bài toán thuộc dạng tổng hợp lực từ tác dụng lên dây dẫn mang dòng điện. Hãy xác định tính đúng, sai của bốn phát biểu sau.
\choiceTF
{\True Lực từ vuông góc với cả dòng điện và vectơ cảm ứng từ.}
{Độ lớn lực từ không phụ thuộc chiều dài đoạn dây.}
{\True Đổi chiều dòng điện thì chiều lực từ đổi ngược.}
{Lực từ luôn cùng chiều với dòng điện.}
\loigiai{
\begin{itemchoice}
\item (Đúng) Lực từ vuông góc với cả dòng điện và vectơ cảm ứng từ. (Vì): lpha); chiều lực xác định bằng quy tắc bàn tay trái
\item (Sai) Độ lớn lực từ không phụ thuộc chiều dài đoạn dây.
\item (Đúng) Đổi chiều dòng điện thì chiều lực từ đổi ngược.
\item (Sai) Lực từ luôn cùng chiều với dòng điện.
\end{itemchoice}
}
\end{ex}

% ===== Câu 167 =====
% ID: L12C3B20-37-TL
% Mức: TH
\dangbt{Tổng hợp đường sức từ và xác định chiều từ trường}
\begin{bt}
% ID: L12C3B20-37-TL
% Mức: TH
Trình bày kiến thức cốt lõi và các bước giải bài thuộc dạng tổng hợp đường sức từ và xác định chiều từ trường.
\loigiai{
Đường sức từ không cắt nhau; chiều từ trường được xác định bằng kim nam châm hoặc quy tắc nắm tay phải. Bên ngoài nam châm, đường sức đi từ cực Bắc đến cực Nam. Các bước: tóm tắt dữ kiện; xác định phương, chiều và điều kiện; chọn hệ thức; đổi đơn vị; tính toán; kiểm tra kết quả.
}
\end{bt}

% ===== Câu 168 =====
% ID: L12C3B20-37-TLN
% Mức: TH
\begin{ex}
% ID: L12C3B20-37-TLN
% Mức: TH
Trong bốn phát biểu sau về tổng hợp đường sức từ và xác định chiều từ trường, có bao nhiêu phát biểu đúng? Chỉ ghi một số nguyên. (1) Bên ngoài nam châm, đường sức từ có chiều đi ra từ cực Bắc và đi vào cực Nam. (2) Các đường sức từ có thể cắt nhau tại một điểm. (3) Đường sức từ của dòng điện thẳng là các đường tròn đồng tâm nằm trong mặt phẳng vuông góc với dây dẫn. (4) Đường sức từ luôn là những đường thẳng song song.
\shortans{2}
\loigiai{
Đối chiếu từng phát biểu với kiến thức: Đường sức từ không cắt nhau; chiều từ trường được xác định bằng kim nam châm hoặc quy tắc nắm tay phải. Bên ngoài nam châm, đường sức đi từ cực Bắc đến cực Nam. Có 2 phát biểu đúng.
}
\end{ex}

% ===== Câu 169 =====
% ID: L12C3B20-37-TN
% Mức: TH
\dangbt{Tổng hợp máy phát điện xoay chiều và đồ thị điện từ}
\begin{ex}
% ID: L12C3B20-37-TN
% Mức: TH
Trong Chương III, kết luận nào đúng khi giải bài thuộc dạng tổng hợp máy phát điện xoay chiều và đồ thị điện từ?
\choice
{Từ thông và suất điện động luôn cùng pha.}
{Suất điện động cực đại không phụ thuộc số vòng dây.}
{\True Tần số suất điện động bằng tần số quay của khung trong mô hình một cặp cực.}
{Máy phát điện biến điện năng thành cơ năng.}
\loigiai{
Tần số suất điện động bằng tần số quay của khung trong mô hình một cặp cực. Do đó chọn C. Máy phát điện biến cơ năng thành điện năng nhờ cảm ứng điện từ; nếu Phi = Phi0 cos(omega t + phi) thì e = E0 sin(omega t + phi), E0 = NBS omega.
}
\end{ex}

% ===== Câu 170 =====
% ID: L12C3B20-38-DS
% Mức: VD
\dangbt{Tổng hợp lực từ tác dụng lên dây dẫn mang dòng điện}
\begin{ex}
% ID: L12C3B20-38-DS
% Mức: VD
Xét bài toán thuộc dạng tổng hợp lực từ tác dụng lên dây dẫn mang dòng điện. Hãy xác định tính đúng, sai của bốn phát biểu sau.
\choiceTF
{Độ lớn lực từ không phụ thuộc chiều dài đoạn dây.}
{\True Đổi chiều dòng điện thì chiều lực từ đổi ngược.}
{Lực từ luôn cùng chiều với dòng điện.}
{\True Lực từ vuông góc với cả dòng điện và vectơ cảm ứng từ.}
\loigiai{
\begin{itemchoice}
\item (Sai) Độ lớn lực từ không phụ thuộc chiều dài đoạn dây. (Vì): lpha); chiều lực xác định bằng quy tắc bàn tay trái
\item (Đúng) Đổi chiều dòng điện thì chiều lực từ đổi ngược.
\item (Sai) Lực từ luôn cùng chiều với dòng điện.
\item (Đúng) Lực từ vuông góc với cả dòng điện và vectơ cảm ứng từ.
\end{itemchoice}
}
\end{ex}

% ===== Câu 171 =====
% ID: L12C3B20-38-TL
% Mức: TH
\dangbt{Tổng hợp đường sức từ và xác định chiều từ trường}
\begin{bt}
% ID: L12C3B20-38-TL
% Mức: TH
Phân tích các đại lượng, phương chiều và điều kiện áp dụng trong dạng tổng hợp đường sức từ và xác định chiều từ trường.
\loigiai{
Đường sức từ không cắt nhau; chiều từ trường được xác định bằng kim nam châm hoặc quy tắc nắm tay phải. Bên ngoài nam châm, đường sức đi từ cực Bắc đến cực Nam. Cần xác định rõ đại lượng vô hướng hay vectơ, quy ước chiều và điều kiện để công thức được áp dụng.
}
\end{bt}

% ===== Câu 172 =====
% ID: L12C3B20-38-TLN
% Mức: TH
\begin{ex}
% ID: L12C3B20-38-TLN
% Mức: TH
Trong bốn phát biểu sau về tổng hợp đường sức từ và xác định chiều từ trường, có bao nhiêu phát biểu đúng? Chỉ ghi một số nguyên. (1) Chiều từ trường của dòng điện được xác định bằng quy tắc nắm tay phải. (2) Tại một điểm, hướng của vectơ cảm ứng từ trùng với hướng mà cực Bắc của kim nam châm chỉ. (3) Đổi chiều dòng điện không làm đổi chiều từ trường. (4) Từ trường đều có các đường sức cong và mật độ thay đổi.
\shortans{2}
\loigiai{
Đối chiếu từng phát biểu với kiến thức: Đường sức từ không cắt nhau; chiều từ trường được xác định bằng kim nam châm hoặc quy tắc nắm tay phải. Bên ngoài nam châm, đường sức đi từ cực Bắc đến cực Nam. Có 2 phát biểu đúng.
}
\end{ex}

% ===== Câu 173 =====
% ID: L12C3B20-38-TN
% Mức: VD
\dangbt{Tổng hợp máy phát điện xoay chiều và đồ thị điện từ}
\begin{ex}
% ID: L12C3B20-38-TN
% Mức: VD
Chọn nhận định phù hợp với kiến thức Vật lí về tổng hợp máy phát điện xoay chiều và đồ thị điện từ.
\choice
{Suất điện động cực đại không phụ thuộc số vòng dây.}
{Máy phát điện biến điện năng thành cơ năng.}
{Tăng tốc độ quay làm giảm biên độ suất điện động.}
{\True Suất điện động cảm ứng lệch pha pi/2 so với từ thông.}
\loigiai{
Suất điện động cảm ứng lệch pha pi/2 so với từ thông. Do đó chọn D. Máy phát điện biến cơ năng thành điện năng nhờ cảm ứng điện từ; nếu Phi = Phi0 cos(omega t + phi) thì e = E0 sin(omega t + phi), E0 = NBS omega.
}
\end{ex}

% ===== Câu 174 =====
% ID: L12C3B20-39-DS
% Mức: VD
\dangbt{Tổng hợp lực từ tác dụng lên dây dẫn mang dòng điện}
\begin{ex}
% ID: L12C3B20-39-DS
% Mức: VD
Xét bài toán thuộc dạng tổng hợp lực từ tác dụng lên dây dẫn mang dòng điện. Hãy xác định tính đúng, sai của bốn phát biểu sau.
\choiceTF
{\True Đổi chiều dòng điện thì chiều lực từ đổi ngược.}
{\True Độ lớn lực từ tác dụng lên đoạn dây thẳng là $F$ = BIl sin(alpha).}
{Khi dây dẫn vuông góc với từ trường thì lực từ bằng 0.}
{Độ lớn lực từ không phụ thuộc chiều dài đoạn dây.}
\loigiai{
\begin{itemchoice}
\item (Đúng) Đổi chiều dòng điện thì chiều lực từ đổi ngược. (Vì): lpha); chiều lực xác định bằng quy tắc bàn tay trái
\item (Đúng) Độ lớn lực từ tác dụng lên đoạn dây thẳng là $F$ = BIl sin(alpha).
\item (Sai) Khi dây dẫn vuông góc với từ trường thì lực từ bằng 0.
\item (Sai) Độ lớn lực từ không phụ thuộc chiều dài đoạn dây.
\end{itemchoice}
}
\end{ex}

% ===== Câu 175 =====
% ID: L12C3B20-39-TL
% Mức: VD
\dangbt{Tổng hợp đường sức từ và xác định chiều từ trường}
\begin{bt}
% ID: L12C3B20-39-TL
% Mức: VD
Nêu một sai lầm thường gặp khi giải tổng hợp đường sức từ và xác định chiều từ trường và cách khắc phục.
\loigiai{
Đường sức từ không cắt nhau; chiều từ trường được xác định bằng kim nam châm hoặc quy tắc nắm tay phải. Bên ngoài nam châm, đường sức đi từ cực Bắc đến cực Nam. Sai lầm thường gặp là bỏ qua phương chiều, góc hoặc đơn vị. Khắc phục bằng hình vẽ, quy ước dấu và đổi đơn vị trước khi tính.
}
\end{bt}

% ===== Câu 176 =====
% ID: L12C3B20-39-TLN
% Mức: VD
\begin{ex}
% ID: L12C3B20-39-TLN
% Mức: VD
Trong bốn phát biểu sau về tổng hợp đường sức từ và xác định chiều từ trường, có bao nhiêu phát biểu đúng? Chỉ ghi một số nguyên. (1) Bên ngoài nam châm, đường sức từ có chiều đi ra từ cực Bắc và đi vào cực Nam. (2) Chiều từ trường của dòng điện được xác định bằng quy tắc nắm tay phải. (3) Đường sức từ luôn là những đường thẳng song song. (4) Từ trường đều có các đường sức cong và mật độ thay đổi.
\shortans{2}
\loigiai{
Đối chiếu từng phát biểu với kiến thức: Đường sức từ không cắt nhau; chiều từ trường được xác định bằng kim nam châm hoặc quy tắc nắm tay phải. Bên ngoài nam châm, đường sức đi từ cực Bắc đến cực Nam. Có 2 phát biểu đúng.
}
\end{ex}

% ===== Câu 177 =====
% ID: L12C3B20-39-TN
% Mức: VD
\dangbt{Tổng hợp máy phát điện xoay chiều và đồ thị điện từ}
\begin{ex}
% ID: L12C3B20-39-TN
% Mức: VD
Khi phân tích tổng hợp máy phát điện xoay chiều và đồ thị điện từ, phát biểu nào đúng?
\choice
{\True Máy phát điện xoay chiều hoạt động dựa trên hiện tượng cảm ứng điện từ.}
{Máy phát điện biến điện năng thành cơ năng.}
{Tăng tốc độ quay làm giảm biên độ suất điện động.}
{Từ thông và suất điện động luôn cùng pha.}
\loigiai{
Máy phát điện xoay chiều hoạt động dựa trên hiện tượng cảm ứng điện từ. Do đó chọn A. Máy phát điện biến cơ năng thành điện năng nhờ cảm ứng điện từ; nếu Phi = Phi0 cos(omega t + phi) thì e = E0 sin(omega t + phi), E0 = NBS omega.
}
\end{ex}

% ===== Câu 178 =====
% ID: L12C3B20-40-DS
% Mức: VD
\dangbt{Tổng hợp lực từ tác dụng lên dây dẫn mang dòng điện}
\begin{ex}
% ID: L12C3B20-40-DS
% Mức: VD
Xét bài toán thuộc dạng tổng hợp lực từ tác dụng lên dây dẫn mang dòng điện. Hãy xác định tính đúng, sai của bốn phát biểu sau.
\choiceTF
{Lực từ luôn cùng chiều với dòng điện.}
{\True Lực từ vuông góc với cả dòng điện và vectơ cảm ứng từ.}
{\True Khi dây dẫn song song với đường sức từ thì lực từ bằng 0.}
{Đổi chiều cảm ứng từ nhưng giữ dòng điện thì chiều lực từ không đổi.}
\loigiai{
\begin{itemchoice}
\item (Sai) Lực từ luôn cùng chiều với dòng điện. (Vì): lpha); chiều lực xác định bằng quy tắc bàn tay trái
\item (Đúng) Lực từ vuông góc với cả dòng điện và vectơ cảm ứng từ.
\item (Đúng) Khi dây dẫn song song với đường sức từ thì lực từ bằng 0.
\item (Sai) Đổi chiều cảm ứng từ nhưng giữ dòng điện thì chiều lực từ không đổi.
\end{itemchoice}
}
\end{ex}

% ===== Câu 179 =====
% ID: L12C3B20-40-TL
% Mức: VD
\dangbt{Tổng hợp đường sức từ và xác định chiều từ trường}
\begin{bt}
% ID: L12C3B20-40-TL
% Mức: VD
Đề xuất quy trình kiểm tra kết quả của một bài toán tổng hợp đường sức từ và xác định chiều từ trường.
\loigiai{
Đường sức từ không cắt nhau; chiều từ trường được xác định bằng kim nam châm hoặc quy tắc nắm tay phải. Bên ngoài nam châm, đường sức đi từ cực Bắc đến cực Nam. Kiểm tra thứ nguyên, dấu, giới hạn đặc biệt và sự phù hợp với ý nghĩa vật lí.
}
\end{bt}

% ===== Câu 180 =====
% ID: L12C3B20-40-TLN
% Mức: VD
\begin{ex}
% ID: L12C3B20-40-TLN
% Mức: VD
Trong bốn phát biểu sau về tổng hợp đường sức từ và xác định chiều từ trường, có bao nhiêu phát biểu đúng? Chỉ ghi một số nguyên. (1) Bên ngoài nam châm, đường sức từ có chiều đi ra từ cực Bắc và đi vào cực Nam. (2) Các đường sức từ có thể cắt nhau tại một điểm. (3) Đường sức từ của dòng điện thẳng là các đường tròn đồng tâm nằm trong mặt phẳng vuông góc với dây dẫn. (4) Đường sức từ luôn là những đường thẳng song song.
\shortans{2}
\loigiai{
Đối chiếu từng phát biểu với kiến thức: Đường sức từ không cắt nhau; chiều từ trường được xác định bằng kim nam châm hoặc quy tắc nắm tay phải. Bên ngoài nam châm, đường sức đi từ cực Bắc đến cực Nam. Có 2 phát biểu đúng.
}
\end{ex}

% ===== Câu 181 =====
% ID: L12C3B20-40-TN
% Mức: VD
\dangbt{Tổng hợp máy phát điện xoay chiều và đồ thị điện từ}
\begin{ex}
% ID: L12C3B20-40-TN
% Mức: VD
Cơ sở Vật lí nào sau đây cần được sử dụng trong dạng tổng hợp máy phát điện xoay chiều và đồ thị điện từ?
\choice
{Tăng tốc độ quay làm giảm biên độ suất điện động.}
{\True Khung $N$ vòng quay đều có suất điện động cực đại E0 = NBS omega.}
{Từ thông và suất điện động luôn cùng pha.}
{Suất điện động cực đại không phụ thuộc số vòng dây.}
\loigiai{
Khung $N$ vòng quay đều có suất điện động cực đại E0 = NBS omega. Do đó chọn B. Máy phát điện biến cơ năng thành điện năng nhờ cảm ứng điện từ; nếu Phi = Phi0 cos(omega t + phi) thì e = E0 sin(omega t + phi), E0 = NBS omega.
}
\end{ex}

% ===== Câu 182 =====
% ID: L12C3B20-41-DS
% Mức: TH
\dangbt{Tổng hợp từ thông, định luật Lenz và suất điện động cảm ứng}
\begin{ex}
% ID: L12C3B20-41-DS
% Mức: TH
Xét bài toán thuộc dạng tổng hợp từ thông, định luật lenz và suất điện động cảm ứng. Hãy xác định tính đúng, sai của bốn phát biểu sau.
\choiceTF
{Từ thông là đại lượng vectơ.}
{Dòng điện cảm ứng luôn làm tăng từ thông ban đầu.}
{\True Độ lớn suất điện động cảm ứng bằng độ lớn tốc độ biến thiên từ thông.}
{\True Từ thông không đổi thì không xuất hiện suất điện động cảm ứng.}
\loigiai{
\begin{itemchoice}
\item (Sai) Từ thông là đại lượng vectơ. (Vì): lpha). Suất điện động cảm ứng có độ lớn |e| = |Delta Phi/Delta t|; chiều dòng điện cảm ứng tuân theo định luật Lenz
\item (Sai) Dòng điện cảm ứng luôn làm tăng từ thông ban đầu.
\item (Đúng) Độ lớn suất điện động cảm ứng bằng độ lớn tốc độ biến thiên từ thông.
\item (Đúng) Từ thông không đổi thì không xuất hiện suất điện động cảm ứng.
\end{itemchoice}
}
\end{ex}

% ===== Câu 183 =====
% ID: L12C3B20-41-TL
% Mức: VD
\dangbt{Tổng hợp đường sức từ và xác định chiều từ trường}
\begin{bt}
% ID: L12C3B20-41-TL
% Mức: VD
Xây dựng một tình huống thực tiễn thuộc dạng tổng hợp đường sức từ và xác định chiều từ trường và chỉ ra định luật cần dùng.
\loigiai{
Đường sức từ không cắt nhau; chiều từ trường được xác định bằng kim nam châm hoặc quy tắc nắm tay phải. Bên ngoài nam châm, đường sức đi từ cực Bắc đến cực Nam. Có thể chọn thiết bị hoặc hiện tượng thực tế tương ứng, nêu dữ kiện và chỉ rõ định luật dùng để giải thích hoặc tính toán.
}
\end{bt}

% ===== Câu 184 =====
% ID: L12C3B20-41-TLN
% Mức: VD
\begin{ex}
% ID: L12C3B20-41-TLN
% Mức: VD
Trong bốn phát biểu sau về tổng hợp đường sức từ và xác định chiều từ trường, có bao nhiêu phát biểu đúng? Chỉ ghi một số nguyên. (1) Chiều từ trường của dòng điện được xác định bằng quy tắc nắm tay phải. (2) Tại một điểm, hướng của vectơ cảm ứng từ trùng với hướng mà cực Bắc của kim nam châm chỉ. (3) Đổi chiều dòng điện không làm đổi chiều từ trường. (4) Từ trường đều có các đường sức cong và mật độ thay đổi.
\shortans{2}
\loigiai{
Đối chiếu từng phát biểu với kiến thức: Đường sức từ không cắt nhau; chiều từ trường được xác định bằng kim nam châm hoặc quy tắc nắm tay phải. Bên ngoài nam châm, đường sức đi từ cực Bắc đến cực Nam. Có 2 phát biểu đúng.
}
\end{ex}

% ===== Câu 185 =====
% ID: L12C3B20-41-TN
% Mức: NB
\dangbt{Tổng hợp từ thông, định luật Lenz và suất điện động cảm ứng}
\begin{ex}
% ID: L12C3B20-41-TN
% Mức: NB
Phát biểu nào sau đây đúng về tổng hợp từ thông, định luật lenz và suất điện động cảm ứng?
\choice
{\True Từ thông qua khung phẳng là Phi = BS cos(alpha), với alpha là góc giữa $B$ và pháp tuyến.}
{Từ thông là đại lượng vectơ.}
{Dòng điện cảm ứng luôn làm tăng từ thông ban đầu.}
{Suất điện động cảm ứng chỉ phụ thuộc điện trở của mạch.}
\loigiai{
Từ thông qua khung phẳng là Phi = BS cos(alpha), với alpha là góc giữa $B$ và pháp tuyến. Do đó chọn A. Phi = BS cos(alpha). Suất điện động cảm ứng có độ lớn |e| = |Delta Phi/Delta t|; chiều dòng điện cảm ứng tuân theo định luật Lenz.
}
\end{ex}

% ===== Câu 186 =====
% ID: L12C3B20-42-DS
% Mức: TH
\begin{ex}
% ID: L12C3B20-42-DS
% Mức: TH
Xét bài toán thuộc dạng tổng hợp từ thông, định luật lenz và suất điện động cảm ứng. Hãy xác định tính đúng, sai của bốn phát biểu sau.
\choiceTF
{\True Dòng điện cảm ứng có chiều chống lại sự biến thiên từ thông sinh ra nó.}
{Suất điện động cảm ứng chỉ phụ thuộc điện trở của mạch.}
{\True Từ thông không đổi thì không xuất hiện suất điện động cảm ứng.}
{Từ thông là đại lượng vectơ.}
\loigiai{
\begin{itemchoice}
\item (Đúng) Dòng điện cảm ứng có chiều chống lại sự biến thiên từ thông sinh ra nó. (Vì): lpha). Suất điện động cảm ứng có độ lớn |e| = |Delta Phi/Delta t|; chiều dòng điện cảm ứng tuân theo định luật Lenz
\item (Sai) Suất điện động cảm ứng chỉ phụ thuộc điện trở của mạch.
\item (Đúng) Từ thông không đổi thì không xuất hiện suất điện động cảm ứng.
\item (Sai) Từ thông là đại lượng vectơ.
\end{itemchoice}
}
\end{ex}

% ===== Câu 187 =====
% ID: L12C3B20-42-TL
% Mức: VDC
\dangbt{Tổng hợp đường sức từ và xác định chiều từ trường}
\begin{bt}
% ID: L12C3B20-42-TL
% Mức: VDC
So sánh một nhận định đúng và một nhận định sai trong dạng tổng hợp đường sức từ và xác định chiều từ trường; giải thích căn cứ Vật lí.
\loigiai{
Nhận định đúng: Bên ngoài nam châm, đường sức từ có chiều đi ra từ cực Bắc và đi vào cực Nam. Nhận định sai: Các đường sức từ có thể cắt nhau tại một điểm. Căn cứ: Đường sức từ không cắt nhau; chiều từ trường được xác định bằng kim nam châm hoặc quy tắc nắm tay phải. Bên ngoài nam châm, đường sức đi từ cực Bắc đến cực Nam.
}
\end{bt}

% ===== Câu 188 =====
% ID: L12C3B20-42-TLN
% Mức: VDC
\begin{ex}
% ID: L12C3B20-42-TLN
% Mức: VDC
Trong bốn phát biểu sau về tổng hợp đường sức từ và xác định chiều từ trường, có bao nhiêu phát biểu đúng? Chỉ ghi một số nguyên. (1) Bên ngoài nam châm, đường sức từ có chiều đi ra từ cực Bắc và đi vào cực Nam. (2) Chiều từ trường của dòng điện được xác định bằng quy tắc nắm tay phải. (3) Đường sức từ luôn là những đường thẳng song song. (4) Từ trường đều có các đường sức cong và mật độ thay đổi.
\shortans{2}
\loigiai{
Đối chiếu từng phát biểu với kiến thức: Đường sức từ không cắt nhau; chiều từ trường được xác định bằng kim nam châm hoặc quy tắc nắm tay phải. Bên ngoài nam châm, đường sức đi từ cực Bắc đến cực Nam. Có 2 phát biểu đúng.
}
\end{ex}

% ===== Câu 189 =====
% ID: L12C3B20-42-TN
% Mức: NB
\dangbt{Tổng hợp từ thông, định luật Lenz và suất điện động cảm ứng}
\begin{ex}
% ID: L12C3B20-42-TN
% Mức: NB
Trong Chương III, kết luận nào đúng khi giải bài thuộc dạng tổng hợp từ thông, định luật lenz và suất điện động cảm ứng?
\choice
{Dòng điện cảm ứng luôn làm tăng từ thông ban đầu.}
{\True Dòng điện cảm ứng có chiều chống lại sự biến thiên từ thông sinh ra nó.}
{Suất điện động cảm ứng chỉ phụ thuộc điện trở của mạch.}
{Khung dây đứng yên trong từ trường đều không đổi luôn có dòng điện cảm ứng.}
\loigiai{
Dòng điện cảm ứng có chiều chống lại sự biến thiên từ thông sinh ra nó. Do đó chọn B. Phi = BS cos(alpha). Suất điện động cảm ứng có độ lớn |e| = |Delta Phi/Delta t|; chiều dòng điện cảm ứng tuân theo định luật Lenz.
}
\end{ex}

% ===== Câu 190 =====
% ID: L12C3B20-43-DS
% Mức: VD
\begin{ex}
% ID: L12C3B20-43-DS
% Mức: VD
Xét bài toán thuộc dạng tổng hợp từ thông, định luật lenz và suất điện động cảm ứng. Hãy xác định tính đúng, sai của bốn phát biểu sau.
\choiceTF
{Suất điện động cảm ứng chỉ phụ thuộc điện trở của mạch.}
{\True Từ thông không đổi thì không xuất hiện suất điện động cảm ứng.}
{Từ thông là đại lượng vectơ.}
{\True Dòng điện cảm ứng có chiều chống lại sự biến thiên từ thông sinh ra nó.}
\loigiai{
\begin{itemchoice}
\item (Sai) Suất điện động cảm ứng chỉ phụ thuộc điện trở của mạch. (Vì): lpha). Suất điện động cảm ứng có độ lớn |e| = |Delta Phi/Delta t|; chiều dòng điện cảm ứng tuân theo định luật Lenz
\item (Đúng) Từ thông không đổi thì không xuất hiện suất điện động cảm ứng.
\item (Sai) Từ thông là đại lượng vectơ.
\item (Đúng) Dòng điện cảm ứng có chiều chống lại sự biến thiên từ thông sinh ra nó.
\end{itemchoice}
}
\end{ex}

% ===== Câu 191 =====
% ID: L12C3B20-43-TL
% Mức: TH
\dangbt{Tổng hợp lực từ tác dụng lên dây dẫn mang dòng điện}
\begin{bt}
% ID: L12C3B20-43-TL
% Mức: TH
Trình bày kiến thức cốt lõi và các bước giải bài thuộc dạng tổng hợp lực từ tác dụng lên dây dẫn mang dòng điện.
\loigiai{
Với đoạn dây thẳng dài l mang dòng $I$ trong từ trường đều $B$, $F$ = BIl sin(alpha); chiều lực xác định bằng quy tắc bàn tay trái. Các bước: tóm tắt dữ kiện; xác định phương, chiều và điều kiện; chọn hệ thức; đổi đơn vị; tính toán; kiểm tra kết quả.
}
\end{bt}

% ===== Câu 192 =====
% ID: L12C3B20-43-TLN
% Mức: TH
\begin{ex}
% ID: L12C3B20-43-TLN
% Mức: TH
Trong bốn phát biểu sau về tổng hợp lực từ tác dụng lên dây dẫn mang dòng điện, có bao nhiêu phát biểu đúng? Chỉ ghi một số nguyên. (1) Độ lớn lực từ tác dụng lên đoạn dây thẳng là $F$ = BIl sin(alpha). (2) Lực từ luôn cùng chiều với dòng điện. (3) Lực từ vuông góc với cả dòng điện và vectơ cảm ứng từ. (4) Khi dây dẫn vuông góc với từ trường thì lực từ bằng 0.
\shortans{2}
\loigiai{
Đối chiếu từng phát biểu với kiến thức: Với đoạn dây thẳng dài l mang dòng $I$ trong từ trường đều $B$, $F$ = BIl sin(alpha); chiều lực xác định bằng quy tắc bàn tay trái. Có 2 phát biểu đúng.
}
\end{ex}

% ===== Câu 193 =====
% ID: L12C3B20-43-TN
% Mức: NB
\dangbt{Tổng hợp từ thông, định luật Lenz và suất điện động cảm ứng}
\begin{ex}
% ID: L12C3B20-43-TN
% Mức: NB
Chọn nhận định phù hợp với kiến thức Vật lí về tổng hợp từ thông, định luật lenz và suất điện động cảm ứng.
\choice
{Suất điện động cảm ứng chỉ phụ thuộc điện trở của mạch.}
{Khung dây đứng yên trong từ trường đều không đổi luôn có dòng điện cảm ứng.}
{\True Độ lớn suất điện động cảm ứng bằng độ lớn tốc độ biến thiên từ thông.}
{Từ thông là đại lượng vectơ.}
\loigiai{
Độ lớn suất điện động cảm ứng bằng độ lớn tốc độ biến thiên từ thông. Do đó chọn C. Phi = BS cos(alpha). Suất điện động cảm ứng có độ lớn |e| = |Delta Phi/Delta t|; chiều dòng điện cảm ứng tuân theo định luật Lenz.
}
\end{ex}

% ===== Câu 194 =====
% ID: L12C3B20-44-DS
% Mức: VD
\begin{ex}
% ID: L12C3B20-44-DS
% Mức: VD
Xét bài toán thuộc dạng tổng hợp từ thông, định luật lenz và suất điện động cảm ứng. Hãy xác định tính đúng, sai của bốn phát biểu sau.
\choiceTF
{\True Từ thông không đổi thì không xuất hiện suất điện động cảm ứng.}
{\True Từ thông qua khung phẳng là Phi = BS cos(alpha), với alpha là góc giữa $B$ và pháp tuyến.}
{Dòng điện cảm ứng luôn làm tăng từ thông ban đầu.}
{Suất điện động cảm ứng chỉ phụ thuộc điện trở của mạch.}
\loigiai{
\begin{itemchoice}
\item (Đúng) Từ thông không đổi thì không xuất hiện suất điện động cảm ứng. (Vì): lpha). Suất điện động cảm ứng có độ lớn |e| = |Delta Phi/Delta t|; chiều dòng điện cảm ứng tuân theo định luật Lenz
\item (Đúng) Từ thông qua khung phẳng là Phi = BS cos(alpha), với alpha là góc giữa $B$ và pháp tuyến.
\item (Sai) Dòng điện cảm ứng luôn làm tăng từ thông ban đầu.
\item (Sai) Suất điện động cảm ứng chỉ phụ thuộc điện trở của mạch.
\end{itemchoice}
}
\end{ex}

% ===== Câu 195 =====
% ID: L12C3B20-44-TL
% Mức: TH
\dangbt{Tổng hợp lực từ tác dụng lên dây dẫn mang dòng điện}
\begin{bt}
% ID: L12C3B20-44-TL
% Mức: TH
Phân tích các đại lượng, phương chiều và điều kiện áp dụng trong dạng tổng hợp lực từ tác dụng lên dây dẫn mang dòng điện.
\loigiai{
Với đoạn dây thẳng dài l mang dòng $I$ trong từ trường đều $B$, $F$ = BIl sin(alpha); chiều lực xác định bằng quy tắc bàn tay trái. Cần xác định rõ đại lượng vô hướng hay vectơ, quy ước chiều và điều kiện để công thức được áp dụng.
}
\end{bt}

% ===== Câu 196 =====
% ID: L12C3B20-44-TLN
% Mức: TH
\begin{ex}
% ID: L12C3B20-44-TLN
% Mức: TH
Trong bốn phát biểu sau về tổng hợp lực từ tác dụng lên dây dẫn mang dòng điện, có bao nhiêu phát biểu đúng? Chỉ ghi một số nguyên. (1) Khi dây dẫn song song với đường sức từ thì lực từ bằng 0. (2) Đổi chiều dòng điện thì chiều lực từ đổi ngược. (3) Độ lớn lực từ không phụ thuộc chiều dài đoạn dây. (4) Đổi chiều cảm ứng từ nhưng giữ dòng điện thì chiều lực từ không đổi.
\shortans{2}
\loigiai{
Đối chiếu từng phát biểu với kiến thức: Với đoạn dây thẳng dài l mang dòng $I$ trong từ trường đều $B$, $F$ = BIl sin(alpha); chiều lực xác định bằng quy tắc bàn tay trái. Có 2 phát biểu đúng.
}
\end{ex}

% ===== Câu 197 =====
% ID: L12C3B20-44-TN
% Mức: TH
\dangbt{Tổng hợp từ thông, định luật Lenz và suất điện động cảm ứng}
\begin{ex}
% ID: L12C3B20-44-TN
% Mức: TH
Khi phân tích tổng hợp từ thông, định luật lenz và suất điện động cảm ứng, phát biểu nào đúng?
\choice
{Khung dây đứng yên trong từ trường đều không đổi luôn có dòng điện cảm ứng.}
{Từ thông là đại lượng vectơ.}
{Dòng điện cảm ứng luôn làm tăng từ thông ban đầu.}
{\True Từ thông không đổi thì không xuất hiện suất điện động cảm ứng.}
\loigiai{
Từ thông không đổi thì không xuất hiện suất điện động cảm ứng. Do đó chọn D. Phi = BS cos(alpha). Suất điện động cảm ứng có độ lớn |e| = |Delta Phi/Delta t|; chiều dòng điện cảm ứng tuân theo định luật Lenz.
}
\end{ex}

% ===== Câu 198 =====
% ID: L12C3B20-45-DS
% Mức: VD
\begin{ex}
% ID: L12C3B20-45-DS
% Mức: VD
Xét bài toán thuộc dạng tổng hợp từ thông, định luật lenz và suất điện động cảm ứng. Hãy xác định tính đúng, sai của bốn phát biểu sau.
\choiceTF
{Từ thông là đại lượng vectơ.}
{\True Dòng điện cảm ứng có chiều chống lại sự biến thiên từ thông sinh ra nó.}
{\True Độ lớn suất điện động cảm ứng bằng độ lớn tốc độ biến thiên từ thông.}
{Khung dây đứng yên trong từ trường đều không đổi luôn có dòng điện cảm ứng.}
\loigiai{
\begin{itemchoice}
\item (Sai) Từ thông là đại lượng vectơ. (Vì): lpha). Suất điện động cảm ứng có độ lớn |e| = |Delta Phi/Delta t|; chiều dòng điện cảm ứng tuân theo định luật Lenz
\item (Đúng) Dòng điện cảm ứng có chiều chống lại sự biến thiên từ thông sinh ra nó.
\item (Đúng) Độ lớn suất điện động cảm ứng bằng độ lớn tốc độ biến thiên từ thông.
\item (Sai) Khung dây đứng yên trong từ trường đều không đổi luôn có dòng điện cảm ứng.
\end{itemchoice}
}
\end{ex}

% ===== Câu 199 =====
% ID: L12C3B20-45-TL
% Mức: VD
\dangbt{Tổng hợp lực từ tác dụng lên dây dẫn mang dòng điện}
\begin{bt}
% ID: L12C3B20-45-TL
% Mức: VD
Nêu một sai lầm thường gặp khi giải tổng hợp lực từ tác dụng lên dây dẫn mang dòng điện và cách khắc phục.
\loigiai{
Với đoạn dây thẳng dài l mang dòng $I$ trong từ trường đều $B$, $F$ = BIl sin(alpha); chiều lực xác định bằng quy tắc bàn tay trái. Sai lầm thường gặp là bỏ qua phương chiều, góc hoặc đơn vị. Khắc phục bằng hình vẽ, quy ước dấu và đổi đơn vị trước khi tính.
}
\end{bt}

% ===== Câu 200 =====
% ID: L12C3B20-45-TLN
% Mức: VD
\begin{ex}
% ID: L12C3B20-45-TLN
% Mức: VD
Trong bốn phát biểu sau về tổng hợp lực từ tác dụng lên dây dẫn mang dòng điện, có bao nhiêu phát biểu đúng? Chỉ ghi một số nguyên. (1) Độ lớn lực từ tác dụng lên đoạn dây thẳng là $F$ = BIl sin(alpha). (2) Khi dây dẫn song song với đường sức từ thì lực từ bằng 0. (3) Khi dây dẫn vuông góc với từ trường thì lực từ bằng 0. (4) Đổi chiều cảm ứng từ nhưng giữ dòng điện thì chiều lực từ không đổi.
\shortans{2}
\loigiai{
Đối chiếu từng phát biểu với kiến thức: Với đoạn dây thẳng dài l mang dòng $I$ trong từ trường đều $B$, $F$ = BIl sin(alpha); chiều lực xác định bằng quy tắc bàn tay trái. Có 2 phát biểu đúng.
}
\end{ex}

% ===== Câu 201 =====
% ID: L12C3B20-45-TN
% Mức: TH
\dangbt{Tổng hợp từ thông, định luật Lenz và suất điện động cảm ứng}
\begin{ex}
% ID: L12C3B20-45-TN
% Mức: TH
Cơ sở Vật lí nào sau đây cần được sử dụng trong dạng tổng hợp từ thông, định luật lenz và suất điện động cảm ứng?
\choice
{\True Từ thông qua khung phẳng là Phi = BS cos(alpha), với alpha là góc giữa $B$ và pháp tuyến.}
{Từ thông là đại lượng vectơ.}
{Dòng điện cảm ứng luôn làm tăng từ thông ban đầu.}
{Suất điện động cảm ứng chỉ phụ thuộc điện trở của mạch.}
\loigiai{
Từ thông qua khung phẳng là Phi = BS cos(alpha), với alpha là góc giữa $B$ và pháp tuyến. Do đó chọn A. Phi = BS cos(alpha). Suất điện động cảm ứng có độ lớn |e| = |Delta Phi/Delta t|; chiều dòng điện cảm ứng tuân theo định luật Lenz.
}
\end{ex}

% ===== Câu 202 =====
% ID: L12C3B20-46-DS
% Mức: TH
\dangbt{Tổng hợp máy phát điện xoay chiều và đồ thị điện từ}
\begin{ex}
% ID: L12C3B20-46-DS
% Mức: TH
Xét bài toán thuộc dạng tổng hợp máy phát điện xoay chiều và đồ thị điện từ. Hãy xác định tính đúng, sai của bốn phát biểu sau.
\choiceTF
{Máy phát điện biến điện năng thành cơ năng.}
{Tăng tốc độ quay làm giảm biên độ suất điện động.}
{\True Tần số suất điện động bằng tần số quay của khung trong mô hình một cặp cực.}
{\True Suất điện động cảm ứng lệch pha pi/2 so với từ thông.}
\loigiai{
\begin{itemchoice}
\item (Sai) Máy phát điện biến điện năng thành cơ năng. (Vì): Máy phát điện biến điện năng thành cơ năng. b) Sai: Tăng tốc độ quay làm giảm biên độ suất điện động. c) Đúng: Tần số suất điện động bằng tần số quay của khung trong mô hình một cặp cực. d) Đúng: Suất điện động cảm ứng lệch pha pi/2 so với từ thông.. Máy phát điện biến cơ năng thành điện năng nhờ cảm ứng điện từ; nếu Phi = Phi0 cos(omega t + phi) thì e = E0 sin(omega t + phi), E0 = NBS omega
\item (Sai) Tăng tốc độ quay làm giảm biên độ suất điện động.
\item (Đúng) Tần số suất điện động bằng tần số quay của khung trong mô hình một cặp cực.
\item (Đúng) Suất điện động cảm ứng lệch pha pi/2 so với từ thông.
\end{itemchoice}
}
\end{ex}

% ===== Câu 203 =====
% ID: L12C3B20-46-TL
% Mức: VD
\dangbt{Tổng hợp lực từ tác dụng lên dây dẫn mang dòng điện}
\begin{bt}
% ID: L12C3B20-46-TL
% Mức: VD
Đề xuất quy trình kiểm tra kết quả của một bài toán tổng hợp lực từ tác dụng lên dây dẫn mang dòng điện.
\loigiai{
Với đoạn dây thẳng dài l mang dòng $I$ trong từ trường đều $B$, $F$ = BIl sin(alpha); chiều lực xác định bằng quy tắc bàn tay trái. Kiểm tra thứ nguyên, dấu, giới hạn đặc biệt và sự phù hợp với ý nghĩa vật lí.
}
\end{bt}

% ===== Câu 204 =====
% ID: L12C3B20-46-TLN
% Mức: VD
\begin{ex}
% ID: L12C3B20-46-TLN
% Mức: VD
Trong bốn phát biểu sau về tổng hợp lực từ tác dụng lên dây dẫn mang dòng điện, có bao nhiêu phát biểu đúng? Chỉ ghi một số nguyên. (1) Độ lớn lực từ tác dụng lên đoạn dây thẳng là $F$ = BIl sin(alpha). (2) Lực từ luôn cùng chiều với dòng điện. (3) Lực từ vuông góc với cả dòng điện và vectơ cảm ứng từ. (4) Khi dây dẫn vuông góc với từ trường thì lực từ bằng 0.
\shortans{2}
\loigiai{
Đối chiếu từng phát biểu với kiến thức: Với đoạn dây thẳng dài l mang dòng $I$ trong từ trường đều $B$, $F$ = BIl sin(alpha); chiều lực xác định bằng quy tắc bàn tay trái. Có 2 phát biểu đúng.
}
\end{ex}

% ===== Câu 205 =====
% ID: L12C3B20-46-TN
% Mức: TH
\dangbt{Tổng hợp từ thông, định luật Lenz và suất điện động cảm ứng}
\begin{ex}
% ID: L12C3B20-46-TN
% Mức: TH
Phát biểu nào sau đây đúng về tổng hợp từ thông, định luật lenz và suất điện động cảm ứng?
\choice
{Dòng điện cảm ứng luôn làm tăng từ thông ban đầu.}
{\True Dòng điện cảm ứng có chiều chống lại sự biến thiên từ thông sinh ra nó.}
{Suất điện động cảm ứng chỉ phụ thuộc điện trở của mạch.}
{Khung dây đứng yên trong từ trường đều không đổi luôn có dòng điện cảm ứng.}
\loigiai{
Dòng điện cảm ứng có chiều chống lại sự biến thiên từ thông sinh ra nó. Do đó chọn B. Phi = BS cos(alpha). Suất điện động cảm ứng có độ lớn |e| = |Delta Phi/Delta t|; chiều dòng điện cảm ứng tuân theo định luật Lenz.
}
\end{ex}

% ===== Câu 206 =====
% ID: L12C3B20-47-DS
% Mức: TH
\dangbt{Tổng hợp máy phát điện xoay chiều và đồ thị điện từ}
\begin{ex}
% ID: L12C3B20-47-DS
% Mức: TH
Xét bài toán thuộc dạng tổng hợp máy phát điện xoay chiều và đồ thị điện từ. Hãy xác định tính đúng, sai của bốn phát biểu sau.
\choiceTF
{\True Khung $N$ vòng quay đều có suất điện động cực đại E0 = NBS omega.}
{Từ thông và suất điện động luôn cùng pha.}
{\True Suất điện động cảm ứng lệch pha pi/2 so với từ thông.}
{Máy phát điện biến điện năng thành cơ năng.}
\loigiai{
\begin{itemchoice}
\item (Đúng) Khung $N$ vòng quay đều có suất điện động cực đại E0 = NBS omega. (Vì): Khung $N$ vòng quay đều có suất điện động cực đại E0 = NBS omega. b) Sai: Từ thông và suất điện động luôn cùng pha. c) Đúng: Suất điện động cảm ứng lệch pha pi/2 so với từ thông. d) Sai: Máy phát điện biến điện năng thành cơ năng.. Máy phát điện biến cơ năng thành điện năng nhờ cảm ứng điện từ; nếu Phi = Phi0 cos(omega t + phi) thì e = E0 sin(omega t + phi), E0 = NBS omega
\item (Sai) Từ thông và suất điện động luôn cùng pha.
\item (Đúng) Suất điện động cảm ứng lệch pha pi/2 so với từ thông.
\item (Sai) Máy phát điện biến điện năng thành cơ năng.
\end{itemchoice}
}
\end{ex}

% ===== Câu 207 =====
% ID: L12C3B20-47-TL
% Mức: VD
\dangbt{Tổng hợp lực từ tác dụng lên dây dẫn mang dòng điện}
\begin{bt}
% ID: L12C3B20-47-TL
% Mức: VD
Xây dựng một tình huống thực tiễn thuộc dạng tổng hợp lực từ tác dụng lên dây dẫn mang dòng điện và chỉ ra định luật cần dùng.
\loigiai{
Với đoạn dây thẳng dài l mang dòng $I$ trong từ trường đều $B$, $F$ = BIl sin(alpha); chiều lực xác định bằng quy tắc bàn tay trái. Có thể chọn thiết bị hoặc hiện tượng thực tế tương ứng, nêu dữ kiện và chỉ rõ định luật dùng để giải thích hoặc tính toán.
}
\end{bt}

% ===== Câu 208 =====
% ID: L12C3B20-47-TLN
% Mức: VD
\begin{ex}
% ID: L12C3B20-47-TLN
% Mức: VD
Trong bốn phát biểu sau về tổng hợp lực từ tác dụng lên dây dẫn mang dòng điện, có bao nhiêu phát biểu đúng? Chỉ ghi một số nguyên. (1) Khi dây dẫn song song với đường sức từ thì lực từ bằng 0. (2) Đổi chiều dòng điện thì chiều lực từ đổi ngược. (3) Độ lớn lực từ không phụ thuộc chiều dài đoạn dây. (4) Đổi chiều cảm ứng từ nhưng giữ dòng điện thì chiều lực từ không đổi.
\shortans{2}
\loigiai{
Đối chiếu từng phát biểu với kiến thức: Với đoạn dây thẳng dài l mang dòng $I$ trong từ trường đều $B$, $F$ = BIl sin(alpha); chiều lực xác định bằng quy tắc bàn tay trái. Có 2 phát biểu đúng.
}
\end{ex}

% ===== Câu 209 =====
% ID: L12C3B20-47-TN
% Mức: TH
\dangbt{Tổng hợp từ thông, định luật Lenz và suất điện động cảm ứng}
\begin{ex}
% ID: L12C3B20-47-TN
% Mức: TH
Trong Chương III, kết luận nào đúng khi giải bài thuộc dạng tổng hợp từ thông, định luật lenz và suất điện động cảm ứng?
\choice
{Suất điện động cảm ứng chỉ phụ thuộc điện trở của mạch.}
{Khung dây đứng yên trong từ trường đều không đổi luôn có dòng điện cảm ứng.}
{\True Độ lớn suất điện động cảm ứng bằng độ lớn tốc độ biến thiên từ thông.}
{Từ thông là đại lượng vectơ.}
\loigiai{
Độ lớn suất điện động cảm ứng bằng độ lớn tốc độ biến thiên từ thông. Do đó chọn C. Phi = BS cos(alpha). Suất điện động cảm ứng có độ lớn |e| = |Delta Phi/Delta t|; chiều dòng điện cảm ứng tuân theo định luật Lenz.
}
\end{ex}

% ===== Câu 210 =====
% ID: L12C3B20-48-DS
% Mức: VD
\dangbt{Tổng hợp máy phát điện xoay chiều và đồ thị điện từ}
\begin{ex}
% ID: L12C3B20-48-DS
% Mức: VD
Xét bài toán thuộc dạng tổng hợp máy phát điện xoay chiều và đồ thị điện từ. Hãy xác định tính đúng, sai của bốn phát biểu sau.
\choiceTF
{Từ thông và suất điện động luôn cùng pha.}
{\True Suất điện động cảm ứng lệch pha pi/2 so với từ thông.}
{Máy phát điện biến điện năng thành cơ năng.}
{\True Khung $N$ vòng quay đều có suất điện động cực đại E0 = NBS omega.}
\loigiai{
\begin{itemchoice}
\item (Sai) Từ thông và suất điện động luôn cùng pha. (Vì): Từ thông và suất điện động luôn cùng pha. b) Đúng: Suất điện động cảm ứng lệch pha pi/2 so với từ thông. c) Sai: Máy phát điện biến điện năng thành cơ năng. d) Đúng: Khung $N$ vòng quay đều có suất điện động cực đại E0 = NBS omega.. Máy phát điện biến cơ năng thành điện năng nhờ cảm ứng điện từ; nếu Phi = Phi0 cos(omega t + phi) thì e = E0 sin(omega t + phi), E0 = NBS omega
\item (Đúng) Suất điện động cảm ứng lệch pha pi/2 so với từ thông.
\item (Sai) Máy phát điện biến điện năng thành cơ năng.
\item (Đúng) Khung $N$ vòng quay đều có suất điện động cực đại E0 = NBS omega.
\end{itemchoice}
}
\end{ex}

% ===== Câu 211 =====
% ID: L12C3B20-48-TL
% Mức: VDC
\dangbt{Tổng hợp lực từ tác dụng lên dây dẫn mang dòng điện}
\begin{bt}
% ID: L12C3B20-48-TL
% Mức: VDC
So sánh một nhận định đúng và một nhận định sai trong dạng tổng hợp lực từ tác dụng lên dây dẫn mang dòng điện; giải thích căn cứ Vật lí.
\loigiai{
Nhận định đúng: Độ lớn lực từ tác dụng lên đoạn dây thẳng là $F$ = BIl sin(alpha). Nhận định sai: Lực từ luôn cùng chiều với dòng điện. Căn cứ: Với đoạn dây thẳng dài l mang dòng $I$ trong từ trường đều $B$, $F$ = BIl sin(alpha); chiều lực xác định bằng quy tắc bàn tay trái.
}
\end{bt}

% ===== Câu 212 =====
% ID: L12C3B20-48-TLN
% Mức: VDC
\begin{ex}
% ID: L12C3B20-48-TLN
% Mức: VDC
Trong bốn phát biểu sau về tổng hợp lực từ tác dụng lên dây dẫn mang dòng điện, có bao nhiêu phát biểu đúng? Chỉ ghi một số nguyên. (1) Độ lớn lực từ tác dụng lên đoạn dây thẳng là $F$ = BIl sin(alpha). (2) Khi dây dẫn song song với đường sức từ thì lực từ bằng 0. (3) Khi dây dẫn vuông góc với từ trường thì lực từ bằng 0. (4) Đổi chiều cảm ứng từ nhưng giữ dòng điện thì chiều lực từ không đổi.
\shortans{2}
\loigiai{
Đối chiếu từng phát biểu với kiến thức: Với đoạn dây thẳng dài l mang dòng $I$ trong từ trường đều $B$, $F$ = BIl sin(alpha); chiều lực xác định bằng quy tắc bàn tay trái. Có 2 phát biểu đúng.
}
\end{ex}

% ===== Câu 213 =====
% ID: L12C3B20-48-TN
% Mức: VD
\dangbt{Tổng hợp từ thông, định luật Lenz và suất điện động cảm ứng}
\begin{ex}
% ID: L12C3B20-48-TN
% Mức: VD
Chọn nhận định phù hợp với kiến thức Vật lí về tổng hợp từ thông, định luật lenz và suất điện động cảm ứng.
\choice
{Khung dây đứng yên trong từ trường đều không đổi luôn có dòng điện cảm ứng.}
{Từ thông là đại lượng vectơ.}
{Dòng điện cảm ứng luôn làm tăng từ thông ban đầu.}
{\True Từ thông không đổi thì không xuất hiện suất điện động cảm ứng.}
\loigiai{
Từ thông không đổi thì không xuất hiện suất điện động cảm ứng. Do đó chọn D. Phi = BS cos(alpha). Suất điện động cảm ứng có độ lớn |e| = |Delta Phi/Delta t|; chiều dòng điện cảm ứng tuân theo định luật Lenz.
}
\end{ex}

% ===== Câu 214 =====
% ID: L12C3B20-49-DS
% Mức: VD
\dangbt{Tổng hợp máy phát điện xoay chiều và đồ thị điện từ}
\begin{ex}
% ID: L12C3B20-49-DS
% Mức: VD
Xét bài toán thuộc dạng tổng hợp máy phát điện xoay chiều và đồ thị điện từ. Hãy xác định tính đúng, sai của bốn phát biểu sau.
\choiceTF
{\True Suất điện động cảm ứng lệch pha pi/2 so với từ thông.}
{\True Máy phát điện xoay chiều hoạt động dựa trên hiện tượng cảm ứng điện từ.}
{Tăng tốc độ quay làm giảm biên độ suất điện động.}
{Từ thông và suất điện động luôn cùng pha.}
\loigiai{
\begin{itemchoice}
\item (Đúng) Suất điện động cảm ứng lệch pha pi/2 so với từ thông. (Vì): Suất điện động cảm ứng lệch pha pi/2 so với từ thông. b) Đúng: Máy phát điện xoay chiều hoạt động dựa trên hiện tượng cảm ứng điện từ. c) Sai: Tăng tốc độ quay làm giảm biên độ suất điện động. d) Sai: Từ thông và suất điện động luôn cùng pha.. Máy phát điện biến cơ năng thành điện năng nhờ cảm ứng điện từ; nếu Phi = Phi0 cos(omega t + phi) thì e = E0 sin(omega t + phi), E0 = NBS omega
\item (Đúng) Máy phát điện xoay chiều hoạt động dựa trên hiện tượng cảm ứng điện từ.
\item (Sai) Tăng tốc độ quay làm giảm biên độ suất điện động.
\item (Sai) Từ thông và suất điện động luôn cùng pha.
\end{itemchoice}
}
\end{ex}

% ===== Câu 215 =====
% ID: L12C3B20-49-TL
% Mức: TH
\dangbt{Tổng hợp từ thông, định luật Lenz và suất điện động cảm ứng}
\begin{bt}
% ID: L12C3B20-49-TL
% Mức: TH
Trình bày kiến thức cốt lõi và các bước giải bài thuộc dạng tổng hợp từ thông, định luật lenz và suất điện động cảm ứng.
\loigiai{
Phi = BS cos(alpha). Suất điện động cảm ứng có độ lớn |e| = |Delta Phi/Delta t|; chiều dòng điện cảm ứng tuân theo định luật Lenz. Các bước: tóm tắt dữ kiện; xác định phương, chiều và điều kiện; chọn hệ thức; đổi đơn vị; tính toán; kiểm tra kết quả.
}
\end{bt}

% ===== Câu 216 =====
% ID: L12C3B20-49-TLN
% Mức: TH
\begin{ex}
% ID: L12C3B20-49-TLN
% Mức: TH
Trong bốn phát biểu sau về tổng hợp từ thông, định luật lenz và suất điện động cảm ứng, có bao nhiêu phát biểu đúng? Chỉ ghi một số nguyên. (1) Từ thông qua khung phẳng là Phi = BS cos(alpha), với alpha là góc giữa $B$ và pháp tuyến. (2) Từ thông là đại lượng vectơ. (3) Dòng điện cảm ứng có chiều chống lại sự biến thiên từ thông sinh ra nó. (4) Dòng điện cảm ứng luôn làm tăng từ thông ban đầu.
\shortans{2}
\loigiai{
Đối chiếu từng phát biểu với kiến thức: Phi = BS cos(alpha). Suất điện động cảm ứng có độ lớn |e| = |Delta Phi/Delta t|; chiều dòng điện cảm ứng tuân theo định luật Lenz. Có 2 phát biểu đúng.
}
\end{ex}

% ===== Câu 217 =====
% ID: L12C3B20-49-TN
% Mức: VD
\begin{ex}
% ID: L12C3B20-49-TN
% Mức: VD
Khi phân tích tổng hợp từ thông, định luật lenz và suất điện động cảm ứng, phát biểu nào đúng?
\choice
{\True Từ thông qua khung phẳng là Phi = BS cos(alpha), với alpha là góc giữa $B$ và pháp tuyến.}
{Từ thông là đại lượng vectơ.}
{Dòng điện cảm ứng luôn làm tăng từ thông ban đầu.}
{Suất điện động cảm ứng chỉ phụ thuộc điện trở của mạch.}
\loigiai{
Từ thông qua khung phẳng là Phi = BS cos(alpha), với alpha là góc giữa $B$ và pháp tuyến. Do đó chọn A. Phi = BS cos(alpha). Suất điện động cảm ứng có độ lớn |e| = |Delta Phi/Delta t|; chiều dòng điện cảm ứng tuân theo định luật Lenz.
}
\end{ex}

% ===== Câu 218 =====
% ID: L12C3B20-50-DS
% Mức: VD
\dangbt{Tổng hợp máy phát điện xoay chiều và đồ thị điện từ}
\begin{ex}
% ID: L12C3B20-50-DS
% Mức: VD
Xét bài toán thuộc dạng tổng hợp máy phát điện xoay chiều và đồ thị điện từ. Hãy xác định tính đúng, sai của bốn phát biểu sau.
\choiceTF
{Máy phát điện biến điện năng thành cơ năng.}
{\True Khung $N$ vòng quay đều có suất điện động cực đại E0 = NBS omega.}
{\True Tần số suất điện động bằng tần số quay của khung trong mô hình một cặp cực.}
{Suất điện động cực đại không phụ thuộc số vòng dây.}
\loigiai{
\begin{itemchoice}
\item (Sai) Máy phát điện biến điện năng thành cơ năng. (Vì): Máy phát điện biến điện năng thành cơ năng. b) Đúng: Khung $N$ vòng quay đều có suất điện động cực đại E0 = NBS omega. c) Đúng: Tần số suất điện động bằng tần số quay của khung trong mô hình một cặp cực. d) Sai: Suất điện động cực đại không phụ thuộc số vòng dây.. Máy phát điện biến cơ năng thành điện năng nhờ cảm ứng điện từ; nếu Phi = Phi0 cos(omega t + phi) thì e = E0 sin(omega t + phi), E0 = NBS omega
\item (Đúng) Khung $N$ vòng quay đều có suất điện động cực đại E0 = NBS omega.
\item (Đúng) Tần số suất điện động bằng tần số quay của khung trong mô hình một cặp cực.
\item (Sai) Suất điện động cực đại không phụ thuộc số vòng dây.
\end{itemchoice}
}
\end{ex}

% ===== Câu 219 =====
% ID: L12C3B20-50-TL
% Mức: TH
\dangbt{Tổng hợp từ thông, định luật Lenz và suất điện động cảm ứng}
\begin{bt}
% ID: L12C3B20-50-TL
% Mức: TH
Phân tích các đại lượng, phương chiều và điều kiện áp dụng trong dạng tổng hợp từ thông, định luật lenz và suất điện động cảm ứng.
\loigiai{
Phi = BS cos(alpha). Suất điện động cảm ứng có độ lớn |e| = |Delta Phi/Delta t|; chiều dòng điện cảm ứng tuân theo định luật Lenz. Cần xác định rõ đại lượng vô hướng hay vectơ, quy ước chiều và điều kiện để công thức được áp dụng.
}
\end{bt}

% ===== Câu 220 =====
% ID: L12C3B20-50-TLN
% Mức: TH
\begin{ex}
% ID: L12C3B20-50-TLN
% Mức: TH
Trong bốn phát biểu sau về tổng hợp từ thông, định luật lenz và suất điện động cảm ứng, có bao nhiêu phát biểu đúng? Chỉ ghi một số nguyên. (1) Độ lớn suất điện động cảm ứng bằng độ lớn tốc độ biến thiên từ thông. (2) Từ thông không đổi thì không xuất hiện suất điện động cảm ứng. (3) Suất điện động cảm ứng chỉ phụ thuộc điện trở của mạch. (4) Khung dây đứng yên trong từ trường đều không đổi luôn có dòng điện cảm ứng.
\shortans{2}
\loigiai{
Đối chiếu từng phát biểu với kiến thức: Phi = BS cos(alpha). Suất điện động cảm ứng có độ lớn |e| = |Delta Phi/Delta t|; chiều dòng điện cảm ứng tuân theo định luật Lenz. Có 2 phát biểu đúng.
}
\end{ex}

% ===== Câu 221 =====
% ID: L12C3B20-50-TN
% Mức: VD
\begin{ex}
% ID: L12C3B20-50-TN
% Mức: VD
Cơ sở Vật lí nào sau đây cần được sử dụng trong dạng tổng hợp từ thông, định luật lenz và suất điện động cảm ứng?
\choice
{Dòng điện cảm ứng luôn làm tăng từ thông ban đầu.}
{\True Dòng điện cảm ứng có chiều chống lại sự biến thiên từ thông sinh ra nó.}
{Suất điện động cảm ứng chỉ phụ thuộc điện trở của mạch.}
{Khung dây đứng yên trong từ trường đều không đổi luôn có dòng điện cảm ứng.}
\loigiai{
Dòng điện cảm ứng có chiều chống lại sự biến thiên từ thông sinh ra nó. Do đó chọn B. Phi = BS cos(alpha). Suất điện động cảm ứng có độ lớn |e| = |Delta Phi/Delta t|; chiều dòng điện cảm ứng tuân theo định luật Lenz.
}
\end{ex}

% ===== Câu 222 =====
% ID: L12C3B20-51-DS
% Mức: TH
\dangbt{Tổng hợp máy biến áp, truyền tải điện năng và ứng dụng cảm ứng điện từ}
\begin{ex}
% ID: L12C3B20-51-DS
% Mức: TH
Xét bài toán thuộc dạng tổng hợp máy biến áp, truyền tải điện năng và ứng dụng cảm ứng điện từ. Hãy xác định tính đúng, sai của bốn phát biểu sau.
\choiceTF
{Máy biến áp làm việc bình thường với nguồn một chiều không đổi.}
{Máy tăng áp luôn có số vòng cuộn thứ cấp nhỏ hơn cuộn sơ cấp.}
{\True Tăng điện áp truyền tải với cùng công suất làm giảm hao phí trên đường dây.}
{\True Dòng điện Foucault được ứng dụng trong bếp từ và phanh điện từ.}
\loigiai{
\begin{itemchoice}
\item (Sai) Máy biến áp làm việc bình thường với nguồn một chiều không đổi. (Vì): Máy biến áp làm việc bình thường với nguồn một chiều không đổi. b) Sai: Máy tăng áp luôn có số vòng cuộn thứ cấp nhỏ hơn cuộn sơ cấp. c) Đúng: Tăng điện áp truyền tải với cùng công suất làm giảm hao phí trên đường dây. d) Đúng: Dòng điện Foucault được ứng dụng trong bếp từ và phanh điện từ.. Máy biến áp lí tưởng có U2/U1 = N2/N1 và U1I1 = U2I2. Truyền tải cùng công suất ở điện áp cao làm giảm hao phí I^2R
\item (Sai) Máy tăng áp luôn có số vòng cuộn thứ cấp nhỏ hơn cuộn sơ cấp.
\item (Đúng) Tăng điện áp truyền tải với cùng công suất làm giảm hao phí trên đường dây.
\item (Đúng) Dòng điện Foucault được ứng dụng trong bếp từ và phanh điện từ.
\end{itemchoice}
}
\end{ex}

% ===== Câu 223 =====
% ID: L12C3B20-51-TL
% Mức: VD
\dangbt{Tổng hợp từ thông, định luật Lenz và suất điện động cảm ứng}
\begin{bt}
% ID: L12C3B20-51-TL
% Mức: VD
Nêu một sai lầm thường gặp khi giải tổng hợp từ thông, định luật lenz và suất điện động cảm ứng và cách khắc phục.
\loigiai{
Phi = BS cos(alpha). Suất điện động cảm ứng có độ lớn |e| = |Delta Phi/Delta t|; chiều dòng điện cảm ứng tuân theo định luật Lenz. Sai lầm thường gặp là bỏ qua phương chiều, góc hoặc đơn vị. Khắc phục bằng hình vẽ, quy ước dấu và đổi đơn vị trước khi tính.
}
\end{bt}

% ===== Câu 224 =====
% ID: L12C3B20-51-TLN
% Mức: VD
\begin{ex}
% ID: L12C3B20-51-TLN
% Mức: VD
Trong bốn phát biểu sau về tổng hợp từ thông, định luật lenz và suất điện động cảm ứng, có bao nhiêu phát biểu đúng? Chỉ ghi một số nguyên. (1) Từ thông qua khung phẳng là Phi = BS cos(alpha), với alpha là góc giữa $B$ và pháp tuyến. (2) Độ lớn suất điện động cảm ứng bằng độ lớn tốc độ biến thiên từ thông. (3) Dòng điện cảm ứng luôn làm tăng từ thông ban đầu. (4) Khung dây đứng yên trong từ trường đều không đổi luôn có dòng điện cảm ứng.
\shortans{2}
\loigiai{
Đối chiếu từng phát biểu với kiến thức: Phi = BS cos(alpha). Suất điện động cảm ứng có độ lớn |e| = |Delta Phi/Delta t|; chiều dòng điện cảm ứng tuân theo định luật Lenz. Có 2 phát biểu đúng.
}
\end{ex}

% ===== Câu 225 =====
% ID: L12C3B20-51-TN
% Mức: NB
\dangbt{Tổng hợp đường sức từ và xác định chiều từ trường}
\begin{ex}
% ID: L12C3B20-51-TN
% Mức: NB
Phát biểu nào sau đây đúng về tổng hợp đường sức từ và xác định chiều từ trường?
\choice
{\True Bên ngoài nam châm, đường sức từ có chiều đi ra từ cực Bắc và đi vào cực Nam.}
{Các đường sức từ có thể cắt nhau tại một điểm.}
{Đường sức từ luôn là những đường thẳng song song.}
{Đổi chiều dòng điện không làm đổi chiều từ trường.}
\loigiai{
Bên ngoài nam châm, đường sức từ có chiều đi ra từ cực Bắc và đi vào cực Nam. Do đó chọn A. Đường sức từ không cắt nhau; chiều từ trường được xác định bằng kim nam châm hoặc quy tắc nắm tay phải. Bên ngoài nam châm, đường sức đi từ cực Bắc đến cực Nam.
}
\end{ex}

% ===== Câu 226 =====
% ID: L12C3B20-52-DS
% Mức: TH
\dangbt{Tổng hợp máy biến áp, truyền tải điện năng và ứng dụng cảm ứng điện từ}
\begin{ex}
% ID: L12C3B20-52-DS
% Mức: TH
Xét bài toán thuộc dạng tổng hợp máy biến áp, truyền tải điện năng và ứng dụng cảm ứng điện từ. Hãy xác định tính đúng, sai của bốn phát biểu sau.
\choiceTF
{\True Máy biến áp lí tưởng thỏa U2/U1 = N2/N1.}
{Hao phí truyền tải tỉ lệ nghịch với bình phương điện trở dây.}
{\True Dòng điện Foucault được ứng dụng trong bếp từ và phanh điện từ.}
{Máy biến áp làm việc bình thường với nguồn một chiều không đổi.}
\loigiai{
\begin{itemchoice}
\item (Đúng) Máy biến áp lí tưởng thỏa U2/U1 = N2/N1. (Vì): Máy biến áp lí tưởng thỏa U2/U1 = N2/N1. b) Sai: Hao phí truyền tải tỉ lệ nghịch với bình phương điện trở dây. c) Đúng: Dòng điện Foucault được ứng dụng trong bếp từ và phanh điện từ. d) Sai: Máy biến áp làm việc bình thường với nguồn một chiều không đổi.. Máy biến áp lí tưởng có U2/U1 = N2/N1 và U1I1 = U2I2. Truyền tải cùng công suất ở điện áp cao làm giảm hao phí I^2R
\item (Sai) Hao phí truyền tải tỉ lệ nghịch với bình phương điện trở dây.
\item (Đúng) Dòng điện Foucault được ứng dụng trong bếp từ và phanh điện từ.
\item (Sai) Máy biến áp làm việc bình thường với nguồn một chiều không đổi.
\end{itemchoice}
}
\end{ex}

% ===== Câu 227 =====
% ID: L12C3B20-52-TL
% Mức: VD
\dangbt{Tổng hợp từ thông, định luật Lenz và suất điện động cảm ứng}
\begin{bt}
% ID: L12C3B20-52-TL
% Mức: VD
Đề xuất quy trình kiểm tra kết quả của một bài toán tổng hợp từ thông, định luật lenz và suất điện động cảm ứng.
\loigiai{
Phi = BS cos(alpha). Suất điện động cảm ứng có độ lớn |e| = |Delta Phi/Delta t|; chiều dòng điện cảm ứng tuân theo định luật Lenz. Kiểm tra thứ nguyên, dấu, giới hạn đặc biệt và sự phù hợp với ý nghĩa vật lí.
}
\end{bt}

% ===== Câu 228 =====
% ID: L12C3B20-52-TLN
% Mức: VD
\begin{ex}
% ID: L12C3B20-52-TLN
% Mức: VD
Trong bốn phát biểu sau về tổng hợp từ thông, định luật lenz và suất điện động cảm ứng, có bao nhiêu phát biểu đúng? Chỉ ghi một số nguyên. (1) Từ thông qua khung phẳng là Phi = BS cos(alpha), với alpha là góc giữa $B$ và pháp tuyến. (2) Từ thông là đại lượng vectơ. (3) Dòng điện cảm ứng có chiều chống lại sự biến thiên từ thông sinh ra nó. (4) Dòng điện cảm ứng luôn làm tăng từ thông ban đầu.
\shortans{2}
\loigiai{
Đối chiếu từng phát biểu với kiến thức: Phi = BS cos(alpha). Suất điện động cảm ứng có độ lớn |e| = |Delta Phi/Delta t|; chiều dòng điện cảm ứng tuân theo định luật Lenz. Có 2 phát biểu đúng.
}
\end{ex}

% ===== Câu 229 =====
% ID: L12C3B20-52-TN
% Mức: NB
\dangbt{Tổng hợp đường sức từ và xác định chiều từ trường}
\begin{ex}
% ID: L12C3B20-52-TN
% Mức: NB
Trong Chương III, kết luận nào đúng khi giải bài thuộc dạng tổng hợp đường sức từ và xác định chiều từ trường?
\choice
{Đường sức từ luôn là những đường thẳng song song.}
{\True Đường sức từ của dòng điện thẳng là các đường tròn đồng tâm nằm trong mặt phẳng vuông góc với dây dẫn.}
{Đổi chiều dòng điện không làm đổi chiều từ trường.}
{Từ trường đều có các đường sức cong và mật độ thay đổi.}
\loigiai{
Đường sức từ của dòng điện thẳng là các đường tròn đồng tâm nằm trong mặt phẳng vuông góc với dây dẫn. Do đó chọn B. Đường sức từ không cắt nhau; chiều từ trường được xác định bằng kim nam châm hoặc quy tắc nắm tay phải. Bên ngoài nam châm, đường sức đi từ cực Bắc đến cực Nam.
}
\end{ex}

% ===== Câu 230 =====
% ID: L12C3B20-53-DS
% Mức: VD
\dangbt{Tổng hợp máy biến áp, truyền tải điện năng và ứng dụng cảm ứng điện từ}
\begin{ex}
% ID: L12C3B20-53-DS
% Mức: VD
Xét bài toán thuộc dạng tổng hợp máy biến áp, truyền tải điện năng và ứng dụng cảm ứng điện từ. Hãy xác định tính đúng, sai của bốn phát biểu sau.
\choiceTF
{Hao phí truyền tải tỉ lệ nghịch với bình phương điện trở dây.}
{\True Dòng điện Foucault được ứng dụng trong bếp từ và phanh điện từ.}
{Máy biến áp làm việc bình thường với nguồn một chiều không đổi.}
{\True Máy biến áp lí tưởng thỏa U2/U1 = N2/N1.}
\loigiai{
\begin{itemchoice}
\item (Sai) Hao phí truyền tải tỉ lệ nghịch với bình phương điện trở dây. (Vì): Hao phí truyền tải tỉ lệ nghịch với bình phương điện trở dây. b) Đúng: Dòng điện Foucault được ứng dụng trong bếp từ và phanh điện từ. c) Sai: Máy biến áp làm việc bình thường với nguồn một chiều không đổi. d) Đúng: Máy biến áp lí tưởng thỏa U2/U1 = N2/N1.. Máy biến áp lí tưởng có U2/U1 = N2/N1 và U1I1 = U2I2. Truyền tải cùng công suất ở điện áp cao làm giảm hao phí I^2R
\item (Đúng) Dòng điện Foucault được ứng dụng trong bếp từ và phanh điện từ.
\item (Sai) Máy biến áp làm việc bình thường với nguồn một chiều không đổi.
\item (Đúng) Máy biến áp lí tưởng thỏa U2/U1 = N2/N1.
\end{itemchoice}
}
\end{ex}

% ===== Câu 231 =====
% ID: L12C3B20-53-TL
% Mức: VD
\dangbt{Tổng hợp từ thông, định luật Lenz và suất điện động cảm ứng}
\begin{bt}
% ID: L12C3B20-53-TL
% Mức: VD
Xây dựng một tình huống thực tiễn thuộc dạng tổng hợp từ thông, định luật lenz và suất điện động cảm ứng và chỉ ra định luật cần dùng.
\loigiai{
Phi = BS cos(alpha). Suất điện động cảm ứng có độ lớn |e| = |Delta Phi/Delta t|; chiều dòng điện cảm ứng tuân theo định luật Lenz. Có thể chọn thiết bị hoặc hiện tượng thực tế tương ứng, nêu dữ kiện và chỉ rõ định luật dùng để giải thích hoặc tính toán.
}
\end{bt}

% ===== Câu 232 =====
% ID: L12C3B20-53-TLN
% Mức: VD
\begin{ex}
% ID: L12C3B20-53-TLN
% Mức: VD
Trong bốn phát biểu sau về tổng hợp từ thông, định luật lenz và suất điện động cảm ứng, có bao nhiêu phát biểu đúng? Chỉ ghi một số nguyên. (1) Độ lớn suất điện động cảm ứng bằng độ lớn tốc độ biến thiên từ thông. (2) Từ thông không đổi thì không xuất hiện suất điện động cảm ứng. (3) Suất điện động cảm ứng chỉ phụ thuộc điện trở của mạch. (4) Khung dây đứng yên trong từ trường đều không đổi luôn có dòng điện cảm ứng.
\shortans{2}
\loigiai{
Đối chiếu từng phát biểu với kiến thức: Phi = BS cos(alpha). Suất điện động cảm ứng có độ lớn |e| = |Delta Phi/Delta t|; chiều dòng điện cảm ứng tuân theo định luật Lenz. Có 2 phát biểu đúng.
}
\end{ex}

% ===== Câu 233 =====
% ID: L12C3B20-53-TN
% Mức: NB
\dangbt{Tổng hợp đường sức từ và xác định chiều từ trường}
\begin{ex}
% ID: L12C3B20-53-TN
% Mức: NB
Chọn nhận định phù hợp với kiến thức Vật lí về tổng hợp đường sức từ và xác định chiều từ trường.
\choice
{Đổi chiều dòng điện không làm đổi chiều từ trường.}
{Từ trường đều có các đường sức cong và mật độ thay đổi.}
{\True Chiều từ trường của dòng điện được xác định bằng quy tắc nắm tay phải.}
{Các đường sức từ có thể cắt nhau tại một điểm.}
\loigiai{
Chiều từ trường của dòng điện được xác định bằng quy tắc nắm tay phải. Do đó chọn C. Đường sức từ không cắt nhau; chiều từ trường được xác định bằng kim nam châm hoặc quy tắc nắm tay phải. Bên ngoài nam châm, đường sức đi từ cực Bắc đến cực Nam.
}
\end{ex}

% ===== Câu 234 =====
% ID: L12C3B20-54-DS
% Mức: VD
\dangbt{Tổng hợp máy biến áp, truyền tải điện năng và ứng dụng cảm ứng điện từ}
\begin{ex}
% ID: L12C3B20-54-DS
% Mức: VD
Xét bài toán thuộc dạng tổng hợp máy biến áp, truyền tải điện năng và ứng dụng cảm ứng điện từ. Hãy xác định tính đúng, sai của bốn phát biểu sau.
\choiceTF
{\True Dòng điện Foucault được ứng dụng trong bếp từ và phanh điện từ.}
{\True Máy biến áp hoạt động với dòng điện xoay chiều dựa trên cảm ứng điện từ.}
{Máy tăng áp luôn có số vòng cuộn thứ cấp nhỏ hơn cuộn sơ cấp.}
{Hao phí truyền tải tỉ lệ nghịch với bình phương điện trở dây.}
\loigiai{
\begin{itemchoice}
\item (Đúng) Dòng điện Foucault được ứng dụng trong bếp từ và phanh điện từ. (Vì): òng điện Foucault được ứng dụng trong bếp từ và phanh điện từ. b) Đúng: Máy biến áp hoạt động với dòng điện xoay chiều dựa trên cảm ứng điện từ. c) Sai: Máy tăng áp luôn có số vòng cuộn thứ cấp nhỏ hơn cuộn sơ cấp. d) Sai: Hao phí truyền tải tỉ lệ nghịch với bình phương điện trở dây.. Máy biến áp lí tưởng có U2/U1 = N2/N1 và U1I1 = U2I2. Truyền tải cùng công suất ở điện áp cao làm giảm hao phí I^2R
\item (Đúng) Máy biến áp hoạt động với dòng điện xoay chiều dựa trên cảm ứng điện từ.
\item (Sai) Máy tăng áp luôn có số vòng cuộn thứ cấp nhỏ hơn cuộn sơ cấp.
\item (Sai) Hao phí truyền tải tỉ lệ nghịch với bình phương điện trở dây.
\end{itemchoice}
}
\end{ex}

% ===== Câu 235 =====
% ID: L12C3B20-54-TL
% Mức: VDC
\dangbt{Tổng hợp từ thông, định luật Lenz và suất điện động cảm ứng}
\begin{bt}
% ID: L12C3B20-54-TL
% Mức: VDC
So sánh một nhận định đúng và một nhận định sai trong dạng tổng hợp từ thông, định luật lenz và suất điện động cảm ứng; giải thích căn cứ Vật lí.
\loigiai{
Nhận định đúng: Từ thông qua khung phẳng là Phi = BS cos(alpha), với alpha là góc giữa $B$ và pháp tuyến. Nhận định sai: Từ thông là đại lượng vectơ. Căn cứ: Phi = BS cos(alpha). Suất điện động cảm ứng có độ lớn |e| = |Delta Phi/Delta t|; chiều dòng điện cảm ứng tuân theo định luật Lenz.
}
\end{bt}

% ===== Câu 236 =====
% ID: L12C3B20-54-TLN
% Mức: VDC
\begin{ex}
% ID: L12C3B20-54-TLN
% Mức: VDC
Trong bốn phát biểu sau về tổng hợp từ thông, định luật lenz và suất điện động cảm ứng, có bao nhiêu phát biểu đúng? Chỉ ghi một số nguyên. (1) Từ thông qua khung phẳng là Phi = BS cos(alpha), với alpha là góc giữa $B$ và pháp tuyến. (2) Độ lớn suất điện động cảm ứng bằng độ lớn tốc độ biến thiên từ thông. (3) Dòng điện cảm ứng luôn làm tăng từ thông ban đầu. (4) Khung dây đứng yên trong từ trường đều không đổi luôn có dòng điện cảm ứng.
\shortans{2}
\loigiai{
Đối chiếu từng phát biểu với kiến thức: Phi = BS cos(alpha). Suất điện động cảm ứng có độ lớn |e| = |Delta Phi/Delta t|; chiều dòng điện cảm ứng tuân theo định luật Lenz. Có 2 phát biểu đúng.
}
\end{ex}

% ===== Câu 237 =====
% ID: L12C3B20-54-TN
% Mức: TH
\dangbt{Tổng hợp đường sức từ và xác định chiều từ trường}
\begin{ex}
% ID: L12C3B20-54-TN
% Mức: TH
Khi phân tích tổng hợp đường sức từ và xác định chiều từ trường, phát biểu nào đúng?
\choice
{Từ trường đều có các đường sức cong và mật độ thay đổi.}
{Các đường sức từ có thể cắt nhau tại một điểm.}
{Đường sức từ luôn là những đường thẳng song song.}
{\True Tại một điểm, hướng của vectơ cảm ứng từ trùng với hướng mà cực Bắc của kim nam châm chỉ.}
\loigiai{
Tại một điểm, hướng của vectơ cảm ứng từ trùng với hướng mà cực Bắc của kim nam châm chỉ. Do đó chọn D. Đường sức từ không cắt nhau; chiều từ trường được xác định bằng kim nam châm hoặc quy tắc nắm tay phải. Bên ngoài nam châm, đường sức đi từ cực Bắc đến cực Nam.
}
\end{ex}

% ===== Câu 238 =====
% ID: L12C3B20-55-DS
% Mức: VD
\dangbt{Tổng hợp máy biến áp, truyền tải điện năng và ứng dụng cảm ứng điện từ}
\begin{ex}
% ID: L12C3B20-55-DS
% Mức: VD
Xét bài toán thuộc dạng tổng hợp máy biến áp, truyền tải điện năng và ứng dụng cảm ứng điện từ. Hãy xác định tính đúng, sai của bốn phát biểu sau.
\choiceTF
{Máy biến áp làm việc bình thường với nguồn một chiều không đổi.}
{\True Máy biến áp lí tưởng thỏa U2/U1 = N2/N1.}
{\True Tăng điện áp truyền tải với cùng công suất làm giảm hao phí trên đường dây.}
{Lõi máy biến áp ghép từ một khối thép đặc để tăng dòng điện Foucault.}
\loigiai{
\begin{itemchoice}
\item (Sai) Máy biến áp làm việc bình thường với nguồn một chiều không đổi. (Vì): Máy biến áp làm việc bình thường với nguồn một chiều không đổi. b) Đúng: Máy biến áp lí tưởng thỏa U2/U1 = N2/N1. c) Đúng: Tăng điện áp truyền tải với cùng công suất làm giảm hao phí trên đường dây. d) Sai: Lõi máy biến áp ghép từ một khối thép đặc để tăng dòng điện Foucault.. Máy biến áp lí tưởng có U2/U1 = N2/N1 và U1I1 = U2I2. Truyền tải cùng công suất ở điện áp cao làm giảm hao phí I^2R
\item (Đúng) Máy biến áp lí tưởng thỏa U2/U1 = N2/N1.
\item (Đúng) Tăng điện áp truyền tải với cùng công suất làm giảm hao phí trên đường dây.
\item (Sai) Lõi máy biến áp ghép từ một khối thép đặc để tăng dòng điện Foucault.
\end{itemchoice}
}
\end{ex}

% ===== Câu 239 =====
% ID: L12C3B20-55-TL
% Mức: TH
\dangbt{Tổng hợp máy phát điện xoay chiều và đồ thị điện từ}
\begin{bt}
% ID: L12C3B20-55-TL
% Mức: TH
Trình bày kiến thức cốt lõi và các bước giải bài thuộc dạng tổng hợp máy phát điện xoay chiều và đồ thị điện từ.
\loigiai{
Máy phát điện biến cơ năng thành điện năng nhờ cảm ứng điện từ; nếu Phi = Phi0 cos(omega t + phi) thì e = E0 sin(omega t + phi), E0 = NBS omega. Các bước: tóm tắt dữ kiện; xác định phương, chiều và điều kiện; chọn hệ thức; đổi đơn vị; tính toán; kiểm tra kết quả.
}
\end{bt}

% ===== Câu 240 =====
% ID: L12C3B20-55-TLN
% Mức: TH
\begin{ex}
% ID: L12C3B20-55-TLN
% Mức: TH
Trong bốn phát biểu sau về tổng hợp máy phát điện xoay chiều và đồ thị điện từ, có bao nhiêu phát biểu đúng? Chỉ ghi một số nguyên. (1) Máy phát điện xoay chiều hoạt động dựa trên hiện tượng cảm ứng điện từ. (2) Máy phát điện biến điện năng thành cơ năng. (3) Khung $N$ vòng quay đều có suất điện động cực đại E0 = NBS omega. (4) Tăng tốc độ quay làm giảm biên độ suất điện động.
\shortans{2}
\loigiai{
Đối chiếu từng phát biểu với kiến thức: Máy phát điện biến cơ năng thành điện năng nhờ cảm ứng điện từ; nếu Phi = Phi0 cos(omega t + phi) thì e = E0 sin(omega t + phi), E0 = NBS omega. Có 2 phát biểu đúng.
}
\end{ex}

% ===== Câu 241 =====
% ID: L12C3B20-55-TN
% Mức: TH
\dangbt{Tổng hợp đường sức từ và xác định chiều từ trường}
\begin{ex}
% ID: L12C3B20-55-TN
% Mức: TH
Cơ sở Vật lí nào sau đây cần được sử dụng trong dạng tổng hợp đường sức từ và xác định chiều từ trường?
\choice
{\True Bên ngoài nam châm, đường sức từ có chiều đi ra từ cực Bắc và đi vào cực Nam.}
{Các đường sức từ có thể cắt nhau tại một điểm.}
{Đường sức từ luôn là những đường thẳng song song.}
{Đổi chiều dòng điện không làm đổi chiều từ trường.}
\loigiai{
Bên ngoài nam châm, đường sức từ có chiều đi ra từ cực Bắc và đi vào cực Nam. Do đó chọn A. Đường sức từ không cắt nhau; chiều từ trường được xác định bằng kim nam châm hoặc quy tắc nắm tay phải. Bên ngoài nam châm, đường sức đi từ cực Bắc đến cực Nam.
}
\end{ex}

% ===== Câu 242 =====
% ID: L12C3B20-56-DS
% Mức: TH
\dangbt{Bài toán vận dụng cao kết hợp nhiều nội dung của Chương III}
\begin{ex}
% ID: L12C3B20-56-DS
% Mức: TH
Xét bài toán thuộc dạng bài toán vận dụng cao kết hợp nhiều nội dung của chương iii. Hãy xác định tính đúng, sai của bốn phát biểu sau.
\choiceTF
{Có thể cộng đại số mọi vectơ mà không xét phương chiều.}
{Có thể dùng trực tiếp số đo góc theo độ trong mọi biểu thức máy tính mà không kiểm tra chế độ.}
{\True Định luật Lenz quyết định chiều, còn định luật Faraday xác định độ lớn suất điện động cảm ứng.}
{\True Kết quả phải được kiểm tra về đơn vị và ý nghĩa vật lí.}
\loigiai{
\begin{itemchoice}
\item (Sai) Có thể cộng đại số mọi vectơ mà không xét phương chiều. (Vì): Có thể cộng đại số mọi vectơ mà không xét phương chiều. b) Sai: Có thể dùng trực tiếp số đo góc theo độ trong mọi biểu thức máy tính mà không kiểm tra chế độ. c) Đúng: Định luật Lenz quyết định chiều, còn định luật Faraday xác định độ lớn suất điện động cảm ứng. d) Đúng: Kết quả phải được kiểm tra về đơn vị và ý nghĩa vật lí.. Bài toán tổng hợp phải tách các giai đoạn, xác định đúng phương chiều, chọn định luật phù hợp, đổi đơn vị thống nhất và kiểm tra ý nghĩa vật lí của kết quả
\item (Sai) Có thể dùng trực tiếp số đo góc theo độ trong mọi biểu thức máy tính mà không kiểm tra chế độ.
\item (Đúng) Định luật Lenz quyết định chiều, còn định luật Faraday xác định độ lớn suất điện động cảm ứng.
\item (Đúng) Kết quả phải được kiểm tra về đơn vị và ý nghĩa vật lí.
\end{itemchoice}
}
\end{ex}

% ===== Câu 243 =====
% ID: L12C3B20-56-TL
% Mức: TH
\dangbt{Tổng hợp máy phát điện xoay chiều và đồ thị điện từ}
\begin{bt}
% ID: L12C3B20-56-TL
% Mức: TH
Phân tích các đại lượng, phương chiều và điều kiện áp dụng trong dạng tổng hợp máy phát điện xoay chiều và đồ thị điện từ.
\loigiai{
Máy phát điện biến cơ năng thành điện năng nhờ cảm ứng điện từ; nếu Phi = Phi0 cos(omega t + phi) thì e = E0 sin(omega t + phi), E0 = NBS omega. Cần xác định rõ đại lượng vô hướng hay vectơ, quy ước chiều và điều kiện để công thức được áp dụng.
}
\end{bt}

% ===== Câu 244 =====
% ID: L12C3B20-56-TLN
% Mức: TH
\begin{ex}
% ID: L12C3B20-56-TLN
% Mức: TH
Trong bốn phát biểu sau về tổng hợp máy phát điện xoay chiều và đồ thị điện từ, có bao nhiêu phát biểu đúng? Chỉ ghi một số nguyên. (1) Tần số suất điện động bằng tần số quay của khung trong mô hình một cặp cực. (2) Suất điện động cảm ứng lệch pha pi/2 so với từ thông. (3) Từ thông và suất điện động luôn cùng pha. (4) Suất điện động cực đại không phụ thuộc số vòng dây.
\shortans{2}
\loigiai{
Đối chiếu từng phát biểu với kiến thức: Máy phát điện biến cơ năng thành điện năng nhờ cảm ứng điện từ; nếu Phi = Phi0 cos(omega t + phi) thì e = E0 sin(omega t + phi), E0 = NBS omega. Có 2 phát biểu đúng.
}
\end{ex}

% ===== Câu 245 =====
% ID: L12C3B20-56-TN
% Mức: TH
\dangbt{Tổng hợp đường sức từ và xác định chiều từ trường}
\begin{ex}
% ID: L12C3B20-56-TN
% Mức: TH
Phát biểu nào sau đây đúng về tổng hợp đường sức từ và xác định chiều từ trường?
\choice
{Đường sức từ luôn là những đường thẳng song song.}
{\True Đường sức từ của dòng điện thẳng là các đường tròn đồng tâm nằm trong mặt phẳng vuông góc với dây dẫn.}
{Đổi chiều dòng điện không làm đổi chiều từ trường.}
{Từ trường đều có các đường sức cong và mật độ thay đổi.}
\loigiai{
Đường sức từ của dòng điện thẳng là các đường tròn đồng tâm nằm trong mặt phẳng vuông góc với dây dẫn. Do đó chọn B. Đường sức từ không cắt nhau; chiều từ trường được xác định bằng kim nam châm hoặc quy tắc nắm tay phải. Bên ngoài nam châm, đường sức đi từ cực Bắc đến cực Nam.
}
\end{ex}

% ===== Câu 246 =====
% ID: L12C3B20-57-DS
% Mức: TH
\dangbt{Bài toán vận dụng cao kết hợp nhiều nội dung của Chương III}
\begin{ex}
% ID: L12C3B20-57-DS
% Mức: TH
Xét bài toán thuộc dạng bài toán vận dụng cao kết hợp nhiều nội dung của chương iii. Hãy xác định tính đúng, sai của bốn phát biểu sau.
\choiceTF
{\True Các đại lượng phải được đổi về cùng hệ đơn vị trước khi thay số.}
{Hiệu suất của mọi thiết bị điện luôn bằng 100 phần trăm.}
{\True Kết quả phải được kiểm tra về đơn vị và ý nghĩa vật lí.}
{Có thể cộng đại số mọi vectơ mà không xét phương chiều.}
\loigiai{
\begin{itemchoice}
\item (Đúng) Các đại lượng phải được đổi về cùng hệ đơn vị trước khi thay số. (Vì): Các đại lượng phải được đổi về cùng hệ đơn vị trước khi thay số. b) Sai: Hiệu suất của mọi thiết bị điện luôn bằng 100 phần trăm. c) Đúng: Kết quả phải được kiểm tra về đơn vị và ý nghĩa vật lí. d) Sai: Có thể cộng đại số mọi vectơ mà không xét phương chiều.. Bài toán tổng hợp phải tách các giai đoạn, xác định đúng phương chiều, chọn định luật phù hợp, đổi đơn vị thống nhất và kiểm tra ý nghĩa vật lí của kết quả
\item (Sai) Hiệu suất của mọi thiết bị điện luôn bằng 100 phần trăm.
\item (Đúng) Kết quả phải được kiểm tra về đơn vị và ý nghĩa vật lí.
\item (Sai) Có thể cộng đại số mọi vectơ mà không xét phương chiều.
\end{itemchoice}
}
\end{ex}

% ===== Câu 247 =====
% ID: L12C3B20-57-TL
% Mức: VD
\dangbt{Tổng hợp máy phát điện xoay chiều và đồ thị điện từ}
\begin{bt}
% ID: L12C3B20-57-TL
% Mức: VD
Nêu một sai lầm thường gặp khi giải tổng hợp máy phát điện xoay chiều và đồ thị điện từ và cách khắc phục.
\loigiai{
Máy phát điện biến cơ năng thành điện năng nhờ cảm ứng điện từ; nếu Phi = Phi0 cos(omega t + phi) thì e = E0 sin(omega t + phi), E0 = NBS omega. Sai lầm thường gặp là bỏ qua phương chiều, góc hoặc đơn vị. Khắc phục bằng hình vẽ, quy ước dấu và đổi đơn vị trước khi tính.
}
\end{bt}

% ===== Câu 248 =====
% ID: L12C3B20-57-TLN
% Mức: VD
\begin{ex}
% ID: L12C3B20-57-TLN
% Mức: VD
Trong bốn phát biểu sau về tổng hợp máy phát điện xoay chiều và đồ thị điện từ, có bao nhiêu phát biểu đúng? Chỉ ghi một số nguyên. (1) Máy phát điện xoay chiều hoạt động dựa trên hiện tượng cảm ứng điện từ. (2) Tần số suất điện động bằng tần số quay của khung trong mô hình một cặp cực. (3) Tăng tốc độ quay làm giảm biên độ suất điện động. (4) Suất điện động cực đại không phụ thuộc số vòng dây.
\shortans{2}
\loigiai{
Đối chiếu từng phát biểu với kiến thức: Máy phát điện biến cơ năng thành điện năng nhờ cảm ứng điện từ; nếu Phi = Phi0 cos(omega t + phi) thì e = E0 sin(omega t + phi), E0 = NBS omega. Có 2 phát biểu đúng.
}
\end{ex}

% ===== Câu 249 =====
% ID: L12C3B20-57-TN
% Mức: TH
\dangbt{Tổng hợp đường sức từ và xác định chiều từ trường}
\begin{ex}
% ID: L12C3B20-57-TN
% Mức: TH
Trong Chương III, kết luận nào đúng khi giải bài thuộc dạng tổng hợp đường sức từ và xác định chiều từ trường?
\choice
{Đổi chiều dòng điện không làm đổi chiều từ trường.}
{Từ trường đều có các đường sức cong và mật độ thay đổi.}
{\True Chiều từ trường của dòng điện được xác định bằng quy tắc nắm tay phải.}
{Các đường sức từ có thể cắt nhau tại một điểm.}
\loigiai{
Chiều từ trường của dòng điện được xác định bằng quy tắc nắm tay phải. Do đó chọn C. Đường sức từ không cắt nhau; chiều từ trường được xác định bằng kim nam châm hoặc quy tắc nắm tay phải. Bên ngoài nam châm, đường sức đi từ cực Bắc đến cực Nam.
}
\end{ex}

% ===== Câu 250 =====
% ID: L12C3B20-58-DS
% Mức: VD
\dangbt{Bài toán vận dụng cao kết hợp nhiều nội dung của Chương III}
\begin{ex}
% ID: L12C3B20-58-DS
% Mức: VD
Xét bài toán thuộc dạng bài toán vận dụng cao kết hợp nhiều nội dung của chương iii. Hãy xác định tính đúng, sai của bốn phát biểu sau.
\choiceTF
{Hiệu suất của mọi thiết bị điện luôn bằng 100 phần trăm.}
{\True Kết quả phải được kiểm tra về đơn vị và ý nghĩa vật lí.}
{Có thể cộng đại số mọi vectơ mà không xét phương chiều.}
{\True Các đại lượng phải được đổi về cùng hệ đơn vị trước khi thay số.}
\loigiai{
\begin{itemchoice}
\item (Sai) Hiệu suất của mọi thiết bị điện luôn bằng 100 phần trăm. (Vì): Hiệu suất của mọi thiết bị điện luôn bằng 100 phần trăm. b) Đúng: Kết quả phải được kiểm tra về đơn vị và ý nghĩa vật lí. c) Sai: Có thể cộng đại số mọi vectơ mà không xét phương chiều. d) Đúng: Các đại lượng phải được đổi về cùng hệ đơn vị trước khi thay số.. Bài toán tổng hợp phải tách các giai đoạn, xác định đúng phương chiều, chọn định luật phù hợp, đổi đơn vị thống nhất và kiểm tra ý nghĩa vật lí của kết quả
\item (Đúng) Kết quả phải được kiểm tra về đơn vị và ý nghĩa vật lí.
\item (Sai) Có thể cộng đại số mọi vectơ mà không xét phương chiều.
\item (Đúng) Các đại lượng phải được đổi về cùng hệ đơn vị trước khi thay số.
\end{itemchoice}
}
\end{ex}

% ===== Câu 251 =====
% ID: L12C3B20-58-TL
% Mức: VD
\dangbt{Tổng hợp máy phát điện xoay chiều và đồ thị điện từ}
\begin{bt}
% ID: L12C3B20-58-TL
% Mức: VD
Đề xuất quy trình kiểm tra kết quả của một bài toán tổng hợp máy phát điện xoay chiều và đồ thị điện từ.
\loigiai{
Máy phát điện biến cơ năng thành điện năng nhờ cảm ứng điện từ; nếu Phi = Phi0 cos(omega t + phi) thì e = E0 sin(omega t + phi), E0 = NBS omega. Kiểm tra thứ nguyên, dấu, giới hạn đặc biệt và sự phù hợp với ý nghĩa vật lí.
}
\end{bt}

% ===== Câu 252 =====
% ID: L12C3B20-58-TLN
% Mức: VD
\begin{ex}
% ID: L12C3B20-58-TLN
% Mức: VD
Trong bốn phát biểu sau về tổng hợp máy phát điện xoay chiều và đồ thị điện từ, có bao nhiêu phát biểu đúng? Chỉ ghi một số nguyên. (1) Máy phát điện xoay chiều hoạt động dựa trên hiện tượng cảm ứng điện từ. (2) Máy phát điện biến điện năng thành cơ năng. (3) Khung $N$ vòng quay đều có suất điện động cực đại E0 = NBS omega. (4) Tăng tốc độ quay làm giảm biên độ suất điện động.
\shortans{2}
\loigiai{
Đối chiếu từng phát biểu với kiến thức: Máy phát điện biến cơ năng thành điện năng nhờ cảm ứng điện từ; nếu Phi = Phi0 cos(omega t + phi) thì e = E0 sin(omega t + phi), E0 = NBS omega. Có 2 phát biểu đúng.
}
\end{ex}

% ===== Câu 253 =====
% ID: L12C3B20-58-TN
% Mức: VD
\dangbt{Tổng hợp đường sức từ và xác định chiều từ trường}
\begin{ex}
% ID: L12C3B20-58-TN
% Mức: VD
Chọn nhận định phù hợp với kiến thức Vật lí về tổng hợp đường sức từ và xác định chiều từ trường.
\choice
{Từ trường đều có các đường sức cong và mật độ thay đổi.}
{Các đường sức từ có thể cắt nhau tại một điểm.}
{Đường sức từ luôn là những đường thẳng song song.}
{\True Tại một điểm, hướng của vectơ cảm ứng từ trùng với hướng mà cực Bắc của kim nam châm chỉ.}
\loigiai{
Tại một điểm, hướng của vectơ cảm ứng từ trùng với hướng mà cực Bắc của kim nam châm chỉ. Do đó chọn D. Đường sức từ không cắt nhau; chiều từ trường được xác định bằng kim nam châm hoặc quy tắc nắm tay phải. Bên ngoài nam châm, đường sức đi từ cực Bắc đến cực Nam.
}
\end{ex}

% ===== Câu 254 =====
% ID: L12C3B20-59-DS
% Mức: VD
\dangbt{Bài toán vận dụng cao kết hợp nhiều nội dung của Chương III}
\begin{ex}
% ID: L12C3B20-59-DS
% Mức: VD
Xét bài toán thuộc dạng bài toán vận dụng cao kết hợp nhiều nội dung của chương iii. Hãy xác định tính đúng, sai của bốn phát biểu sau.
\choiceTF
{\True Kết quả phải được kiểm tra về đơn vị và ý nghĩa vật lí.}
{\True Khi giải bài tổng hợp cần xác định rõ chiều của $B$, dòng điện, lực từ và dòng điện cảm ứng.}
{Có thể dùng trực tiếp số đo góc theo độ trong mọi biểu thức máy tính mà không kiểm tra chế độ.}
{Hiệu suất của mọi thiết bị điện luôn bằng 100 phần trăm.}
\loigiai{
\begin{itemchoice}
\item (Đúng) Kết quả phải được kiểm tra về đơn vị và ý nghĩa vật lí. (Vì): Kết quả phải được kiểm tra về đơn vị và ý nghĩa vật lí. b) Đúng: Khi giải bài tổng hợp cần xác định rõ chiều của $B$, dòng điện, lực từ và dòng điện cảm ứng. c) Sai: Có thể dùng trực tiếp số đo góc theo độ trong mọi biểu thức máy tính mà không kiểm tra chế độ. d) Sai: Hiệu suất của mọi thiết bị điện luôn bằng 100 phần trăm.. Bài toán tổng hợp phải tách các giai đoạn, xác định đúng phương chiều, chọn định luật phù hợp, đổi đơn vị thống nhất và kiểm tra ý nghĩa vật lí của kết quả
\item (Đúng) Khi giải bài tổng hợp cần xác định rõ chiều của $B$, dòng điện, lực từ và dòng điện cảm ứng.
\item (Sai) Có thể dùng trực tiếp số đo góc theo độ trong mọi biểu thức máy tính mà không kiểm tra chế độ.
\item (Sai) Hiệu suất của mọi thiết bị điện luôn bằng 100 phần trăm.
\end{itemchoice}
}
\end{ex}

% ===== Câu 255 =====
% ID: L12C3B20-59-TL
% Mức: VD
\dangbt{Tổng hợp máy phát điện xoay chiều và đồ thị điện từ}
\begin{bt}
% ID: L12C3B20-59-TL
% Mức: VD
Xây dựng một tình huống thực tiễn thuộc dạng tổng hợp máy phát điện xoay chiều và đồ thị điện từ và chỉ ra định luật cần dùng.
\loigiai{
Máy phát điện biến cơ năng thành điện năng nhờ cảm ứng điện từ; nếu Phi = Phi0 cos(omega t + phi) thì e = E0 sin(omega t + phi), E0 = NBS omega. Có thể chọn thiết bị hoặc hiện tượng thực tế tương ứng, nêu dữ kiện và chỉ rõ định luật dùng để giải thích hoặc tính toán.
}
\end{bt}

% ===== Câu 256 =====
% ID: L12C3B20-59-TLN
% Mức: VD
\begin{ex}
% ID: L12C3B20-59-TLN
% Mức: VD
Trong bốn phát biểu sau về tổng hợp máy phát điện xoay chiều và đồ thị điện từ, có bao nhiêu phát biểu đúng? Chỉ ghi một số nguyên. (1) Tần số suất điện động bằng tần số quay của khung trong mô hình một cặp cực. (2) Suất điện động cảm ứng lệch pha pi/2 so với từ thông. (3) Từ thông và suất điện động luôn cùng pha. (4) Suất điện động cực đại không phụ thuộc số vòng dây.
\shortans{2}
\loigiai{
Đối chiếu từng phát biểu với kiến thức: Máy phát điện biến cơ năng thành điện năng nhờ cảm ứng điện từ; nếu Phi = Phi0 cos(omega t + phi) thì e = E0 sin(omega t + phi), E0 = NBS omega. Có 2 phát biểu đúng.
}
\end{ex}

% ===== Câu 257 =====
% ID: L12C3B20-59-TN
% Mức: VD
\dangbt{Tổng hợp đường sức từ và xác định chiều từ trường}
\begin{ex}
% ID: L12C3B20-59-TN
% Mức: VD
Khi phân tích tổng hợp đường sức từ và xác định chiều từ trường, phát biểu nào đúng?
\choice
{\True Bên ngoài nam châm, đường sức từ có chiều đi ra từ cực Bắc và đi vào cực Nam.}
{Các đường sức từ có thể cắt nhau tại một điểm.}
{Đường sức từ luôn là những đường thẳng song song.}
{Đổi chiều dòng điện không làm đổi chiều từ trường.}
\loigiai{
Bên ngoài nam châm, đường sức từ có chiều đi ra từ cực Bắc và đi vào cực Nam. Do đó chọn A. Đường sức từ không cắt nhau; chiều từ trường được xác định bằng kim nam châm hoặc quy tắc nắm tay phải. Bên ngoài nam châm, đường sức đi từ cực Bắc đến cực Nam.
}
\end{ex}

% ===== Câu 258 =====
% ID: L12C3B20-60-DS
% Mức: VD
\dangbt{Bài toán vận dụng cao kết hợp nhiều nội dung của Chương III}
\begin{ex}
% ID: L12C3B20-60-DS
% Mức: VD
Xét bài toán thuộc dạng bài toán vận dụng cao kết hợp nhiều nội dung của chương iii. Hãy xác định tính đúng, sai của bốn phát biểu sau.
\choiceTF
{Có thể cộng đại số mọi vectơ mà không xét phương chiều.}
{\True Các đại lượng phải được đổi về cùng hệ đơn vị trước khi thay số.}
{\True Định luật Lenz quyết định chiều, còn định luật Faraday xác định độ lớn suất điện động cảm ứng.}
{Trong bài tổng hợp có thể bỏ qua điều kiện áp dụng của từng định luật.}
\loigiai{
\begin{itemchoice}
\item (Sai) Có thể cộng đại số mọi vectơ mà không xét phương chiều. (Vì): Có thể cộng đại số mọi vectơ mà không xét phương chiều. b) Đúng: Các đại lượng phải được đổi về cùng hệ đơn vị trước khi thay số. c) Đúng: Định luật Lenz quyết định chiều, còn định luật Faraday xác định độ lớn suất điện động cảm ứng. d) Sai: Trong bài tổng hợp có thể bỏ qua điều kiện áp dụng của từng định luật.. Bài toán tổng hợp phải tách các giai đoạn, xác định đúng phương chiều, chọn định luật phù hợp, đổi đơn vị thống nhất và kiểm tra ý nghĩa vật lí của kết quả
\item (Đúng) Các đại lượng phải được đổi về cùng hệ đơn vị trước khi thay số.
\item (Đúng) Định luật Lenz quyết định chiều, còn định luật Faraday xác định độ lớn suất điện động cảm ứng.
\item (Sai) Trong bài tổng hợp có thể bỏ qua điều kiện áp dụng của từng định luật.
\end{itemchoice}
}
\end{ex}

% ===== Câu 259 =====
% ID: L12C3B20-60-TL
% Mức: VDC
\dangbt{Tổng hợp máy phát điện xoay chiều và đồ thị điện từ}
\begin{bt}
% ID: L12C3B20-60-TL
% Mức: VDC
So sánh một nhận định đúng và một nhận định sai trong dạng tổng hợp máy phát điện xoay chiều và đồ thị điện từ; giải thích căn cứ Vật lí.
\loigiai{
Nhận định đúng: Máy phát điện xoay chiều hoạt động dựa trên hiện tượng cảm ứng điện từ. Nhận định sai: Máy phát điện biến điện năng thành cơ năng. Căn cứ: Máy phát điện biến cơ năng thành điện năng nhờ cảm ứng điện từ; nếu Phi = Phi0 cos(omega t + phi) thì e = E0 sin(omega t + phi), E0 = NBS omega.
}
\end{bt}

% ===== Câu 260 =====
% ID: L12C3B20-60-TLN
% Mức: VDC
\begin{ex}
% ID: L12C3B20-60-TLN
% Mức: VDC
Trong bốn phát biểu sau về tổng hợp máy phát điện xoay chiều và đồ thị điện từ, có bao nhiêu phát biểu đúng? Chỉ ghi một số nguyên. (1) Máy phát điện xoay chiều hoạt động dựa trên hiện tượng cảm ứng điện từ. (2) Tần số suất điện động bằng tần số quay của khung trong mô hình một cặp cực. (3) Tăng tốc độ quay làm giảm biên độ suất điện động. (4) Suất điện động cực đại không phụ thuộc số vòng dây.
\shortans{2}
\loigiai{
Đối chiếu từng phát biểu với kiến thức: Máy phát điện biến cơ năng thành điện năng nhờ cảm ứng điện từ; nếu Phi = Phi0 cos(omega t + phi) thì e = E0 sin(omega t + phi), E0 = NBS omega. Có 2 phát biểu đúng.
}
\end{ex}

% ===== Câu 261 =====
% ID: L12C3B20-60-TN
% Mức: VD
\dangbt{Tổng hợp đường sức từ và xác định chiều từ trường}
\begin{ex}
% ID: L12C3B20-60-TN
% Mức: VD
Cơ sở Vật lí nào sau đây cần được sử dụng trong dạng tổng hợp đường sức từ và xác định chiều từ trường?
\choice
{Đường sức từ luôn là những đường thẳng song song.}
{\True Đường sức từ của dòng điện thẳng là các đường tròn đồng tâm nằm trong mặt phẳng vuông góc với dây dẫn.}
{Đổi chiều dòng điện không làm đổi chiều từ trường.}
{Từ trường đều có các đường sức cong và mật độ thay đổi.}
\loigiai{
Đường sức từ của dòng điện thẳng là các đường tròn đồng tâm nằm trong mặt phẳng vuông góc với dây dẫn. Do đó chọn B. Đường sức từ không cắt nhau; chiều từ trường được xác định bằng kim nam châm hoặc quy tắc nắm tay phải. Bên ngoài nam châm, đường sức đi từ cực Bắc đến cực Nam.
}
\end{ex}

% ===== Câu 262 =====
% ID: L12C3B20-61-TL
% Mức: TH
\dangbt{Tổng hợp máy biến áp, truyền tải điện năng và ứng dụng cảm ứng điện từ}
\begin{bt}
% ID: L12C3B20-61-TL
% Mức: TH
Trình bày kiến thức cốt lõi và các bước giải bài thuộc dạng tổng hợp máy biến áp, truyền tải điện năng và ứng dụng cảm ứng điện từ.
\loigiai{
Máy biến áp lí tưởng có U2/U1 = N2/N1 và U1I1 = U2I2. Truyền tải cùng công suất ở điện áp cao làm giảm hao phí I^2R. Các bước: tóm tắt dữ kiện; xác định phương, chiều và điều kiện; chọn hệ thức; đổi đơn vị; tính toán; kiểm tra kết quả.
}
\end{bt}

% ===== Câu 263 =====
% ID: L12C3B20-61-TLN
% Mức: TH
\begin{ex}
% ID: L12C3B20-61-TLN
% Mức: TH
Trong bốn phát biểu sau về tổng hợp máy biến áp, truyền tải điện năng và ứng dụng cảm ứng điện từ, có bao nhiêu phát biểu đúng? Chỉ ghi một số nguyên. (1) Máy biến áp hoạt động với dòng điện xoay chiều dựa trên cảm ứng điện từ. (2) Máy biến áp làm việc bình thường với nguồn một chiều không đổi. (3) Máy biến áp lí tưởng thỏa U2/U1 = N2/N1. (4) Máy tăng áp luôn có số vòng cuộn thứ cấp nhỏ hơn cuộn sơ cấp.
\shortans{2}
\loigiai{
Đối chiếu từng phát biểu với kiến thức: Máy biến áp lí tưởng có U2/U1 = N2/N1 và U1I1 = U2I2. Truyền tải cùng công suất ở điện áp cao làm giảm hao phí I^2R. Có 2 phát biểu đúng.
}
\end{ex}

% ===== Câu 264 =====
% ID: L12C3B20-61-TN
% Mức: NB
\dangbt{Tổng hợp đường sức từ và xác định chiều từ trường}
\begin{ex}
% ID: L12C3B20-61-TN
% Mức: NB
Phát biểu nào sau đây đúng về tổng hợp đường sức từ và xác định chiều từ trường?
\choice
{\True Bên ngoài nam châm, đường sức từ có chiều đi ra từ cực Bắc và đi vào cực Nam.}
{Các đường sức từ có thể cắt nhau tại một điểm.}
{Đường sức từ luôn là những đường thẳng song song.}
{Đổi chiều dòng điện không làm đổi chiều từ trường.}
\loigiai{
Bên ngoài nam châm, đường sức từ có chiều đi ra từ cực Bắc và đi vào cực Nam. Do đó chọn A. Đường sức từ không cắt nhau; chiều từ trường được xác định bằng kim nam châm hoặc quy tắc nắm tay phải. Bên ngoài nam châm, đường sức đi từ cực Bắc đến cực Nam.
}
\end{ex}

% ===== Câu 265 =====
% ID: L12C3B20-62-TL
% Mức: TH
\dangbt{Tổng hợp máy biến áp, truyền tải điện năng và ứng dụng cảm ứng điện từ}
\begin{bt}
% ID: L12C3B20-62-TL
% Mức: TH
Phân tích các đại lượng, phương chiều và điều kiện áp dụng trong dạng tổng hợp máy biến áp, truyền tải điện năng và ứng dụng cảm ứng điện từ.
\loigiai{
Máy biến áp lí tưởng có U2/U1 = N2/N1 và U1I1 = U2I2. Truyền tải cùng công suất ở điện áp cao làm giảm hao phí I^2R. Cần xác định rõ đại lượng vô hướng hay vectơ, quy ước chiều và điều kiện để công thức được áp dụng.
}
\end{bt}

% ===== Câu 266 =====
% ID: L12C3B20-62-TLN
% Mức: TH
\begin{ex}
% ID: L12C3B20-62-TLN
% Mức: TH
Trong bốn phát biểu sau về tổng hợp máy biến áp, truyền tải điện năng và ứng dụng cảm ứng điện từ, có bao nhiêu phát biểu đúng? Chỉ ghi một số nguyên. (1) Tăng điện áp truyền tải với cùng công suất làm giảm hao phí trên đường dây. (2) Dòng điện Foucault được ứng dụng trong bếp từ và phanh điện từ. (3) Hao phí truyền tải tỉ lệ nghịch với bình phương điện trở dây. (4) Lõi máy biến áp ghép từ một khối thép đặc để tăng dòng điện Foucault.
\shortans{2}
\loigiai{
Đối chiếu từng phát biểu với kiến thức: Máy biến áp lí tưởng có U2/U1 = N2/N1 và U1I1 = U2I2. Truyền tải cùng công suất ở điện áp cao làm giảm hao phí I^2R. Có 2 phát biểu đúng.
}
\end{ex}

% ===== Câu 267 =====
% ID: L12C3B20-62-TN
% Mức: NB
\dangbt{Tổng hợp đường sức từ và xác định chiều từ trường}
\begin{ex}
% ID: L12C3B20-62-TN
% Mức: NB
Trong Chương III, kết luận nào đúng khi giải bài thuộc dạng tổng hợp đường sức từ và xác định chiều từ trường?
\choice
{Đường sức từ luôn là những đường thẳng song song.}
{\True Đường sức từ của dòng điện thẳng là các đường tròn đồng tâm nằm trong mặt phẳng vuông góc với dây dẫn.}
{Đổi chiều dòng điện không làm đổi chiều từ trường.}
{Từ trường đều có các đường sức cong và mật độ thay đổi.}
\loigiai{
Đường sức từ của dòng điện thẳng là các đường tròn đồng tâm nằm trong mặt phẳng vuông góc với dây dẫn. Do đó chọn B. Đường sức từ không cắt nhau; chiều từ trường được xác định bằng kim nam châm hoặc quy tắc nắm tay phải. Bên ngoài nam châm, đường sức đi từ cực Bắc đến cực Nam.
}
\end{ex}

% ===== Câu 268 =====
% ID: L12C3B20-63-TL
% Mức: VD
\dangbt{Tổng hợp máy biến áp, truyền tải điện năng và ứng dụng cảm ứng điện từ}
\begin{bt}
% ID: L12C3B20-63-TL
% Mức: VD
Nêu một sai lầm thường gặp khi giải tổng hợp máy biến áp, truyền tải điện năng và ứng dụng cảm ứng điện từ và cách khắc phục.
\loigiai{
Máy biến áp lí tưởng có U2/U1 = N2/N1 và U1I1 = U2I2. Truyền tải cùng công suất ở điện áp cao làm giảm hao phí I^2R. Sai lầm thường gặp là bỏ qua phương chiều, góc hoặc đơn vị. Khắc phục bằng hình vẽ, quy ước dấu và đổi đơn vị trước khi tính.
}
\end{bt}

% ===== Câu 269 =====
% ID: L12C3B20-63-TLN
% Mức: VD
\begin{ex}
% ID: L12C3B20-63-TLN
% Mức: VD
Trong bốn phát biểu sau về tổng hợp máy biến áp, truyền tải điện năng và ứng dụng cảm ứng điện từ, có bao nhiêu phát biểu đúng? Chỉ ghi một số nguyên. (1) Máy biến áp hoạt động với dòng điện xoay chiều dựa trên cảm ứng điện từ. (2) Tăng điện áp truyền tải với cùng công suất làm giảm hao phí trên đường dây. (3) Máy tăng áp luôn có số vòng cuộn thứ cấp nhỏ hơn cuộn sơ cấp. (4) Lõi máy biến áp ghép từ một khối thép đặc để tăng dòng điện Foucault.
\shortans{2}
\loigiai{
Đối chiếu từng phát biểu với kiến thức: Máy biến áp lí tưởng có U2/U1 = N2/N1 và U1I1 = U2I2. Truyền tải cùng công suất ở điện áp cao làm giảm hao phí I^2R. Có 2 phát biểu đúng.
}
\end{ex}

% ===== Câu 270 =====
% ID: L12C3B20-63-TN
% Mức: NB
\dangbt{Tổng hợp đường sức từ và xác định chiều từ trường}
\begin{ex}
% ID: L12C3B20-63-TN
% Mức: NB
Chọn nhận định phù hợp với kiến thức Vật lí về tổng hợp đường sức từ và xác định chiều từ trường.
\choice
{Đổi chiều dòng điện không làm đổi chiều từ trường.}
{Từ trường đều có các đường sức cong và mật độ thay đổi.}
{\True Chiều từ trường của dòng điện được xác định bằng quy tắc nắm tay phải.}
{Các đường sức từ có thể cắt nhau tại một điểm.}
\loigiai{
Chiều từ trường của dòng điện được xác định bằng quy tắc nắm tay phải. Do đó chọn C. Đường sức từ không cắt nhau; chiều từ trường được xác định bằng kim nam châm hoặc quy tắc nắm tay phải. Bên ngoài nam châm, đường sức đi từ cực Bắc đến cực Nam.
}
\end{ex}

% ===== Câu 271 =====
% ID: L12C3B20-64-TL
% Mức: VD
\dangbt{Tổng hợp máy biến áp, truyền tải điện năng và ứng dụng cảm ứng điện từ}
\begin{bt}
% ID: L12C3B20-64-TL
% Mức: VD
Đề xuất quy trình kiểm tra kết quả của một bài toán tổng hợp máy biến áp, truyền tải điện năng và ứng dụng cảm ứng điện từ.
\loigiai{
Máy biến áp lí tưởng có U2/U1 = N2/N1 và U1I1 = U2I2. Truyền tải cùng công suất ở điện áp cao làm giảm hao phí I^2R. Kiểm tra thứ nguyên, dấu, giới hạn đặc biệt và sự phù hợp với ý nghĩa vật lí.
}
\end{bt}

% ===== Câu 272 =====
% ID: L12C3B20-64-TLN
% Mức: VD
\begin{ex}
% ID: L12C3B20-64-TLN
% Mức: VD
Trong bốn phát biểu sau về tổng hợp máy biến áp, truyền tải điện năng và ứng dụng cảm ứng điện từ, có bao nhiêu phát biểu đúng? Chỉ ghi một số nguyên. (1) Máy biến áp hoạt động với dòng điện xoay chiều dựa trên cảm ứng điện từ. (2) Máy biến áp làm việc bình thường với nguồn một chiều không đổi. (3) Máy biến áp lí tưởng thỏa U2/U1 = N2/N1. (4) Máy tăng áp luôn có số vòng cuộn thứ cấp nhỏ hơn cuộn sơ cấp.
\shortans{2}
\loigiai{
Đối chiếu từng phát biểu với kiến thức: Máy biến áp lí tưởng có U2/U1 = N2/N1 và U1I1 = U2I2. Truyền tải cùng công suất ở điện áp cao làm giảm hao phí I^2R. Có 2 phát biểu đúng.
}
\end{ex}

% ===== Câu 273 =====
% ID: L12C3B20-64-TN
% Mức: TH
\dangbt{Tổng hợp đường sức từ và xác định chiều từ trường}
\begin{ex}
% ID: L12C3B20-64-TN
% Mức: TH
Khi phân tích tổng hợp đường sức từ và xác định chiều từ trường, phát biểu nào đúng?
\choice
{Từ trường đều có các đường sức cong và mật độ thay đổi.}
{Các đường sức từ có thể cắt nhau tại một điểm.}
{Đường sức từ luôn là những đường thẳng song song.}
{\True Tại một điểm, hướng của vectơ cảm ứng từ trùng với hướng mà cực Bắc của kim nam châm chỉ.}
\loigiai{
Tại một điểm, hướng của vectơ cảm ứng từ trùng với hướng mà cực Bắc của kim nam châm chỉ. Do đó chọn D. Đường sức từ không cắt nhau; chiều từ trường được xác định bằng kim nam châm hoặc quy tắc nắm tay phải. Bên ngoài nam châm, đường sức đi từ cực Bắc đến cực Nam.
}
\end{ex}

% ===== Câu 274 =====
% ID: L12C3B20-65-TL
% Mức: VD
\dangbt{Tổng hợp máy biến áp, truyền tải điện năng và ứng dụng cảm ứng điện từ}
\begin{bt}
% ID: L12C3B20-65-TL
% Mức: VD
Xây dựng một tình huống thực tiễn thuộc dạng tổng hợp máy biến áp, truyền tải điện năng và ứng dụng cảm ứng điện từ và chỉ ra định luật cần dùng.
\loigiai{
Máy biến áp lí tưởng có U2/U1 = N2/N1 và U1I1 = U2I2. Truyền tải cùng công suất ở điện áp cao làm giảm hao phí I^2R. Có thể chọn thiết bị hoặc hiện tượng thực tế tương ứng, nêu dữ kiện và chỉ rõ định luật dùng để giải thích hoặc tính toán.
}
\end{bt}

% ===== Câu 275 =====
% ID: L12C3B20-65-TLN
% Mức: VD
\begin{ex}
% ID: L12C3B20-65-TLN
% Mức: VD
Trong bốn phát biểu sau về tổng hợp máy biến áp, truyền tải điện năng và ứng dụng cảm ứng điện từ, có bao nhiêu phát biểu đúng? Chỉ ghi một số nguyên. (1) Tăng điện áp truyền tải với cùng công suất làm giảm hao phí trên đường dây. (2) Dòng điện Foucault được ứng dụng trong bếp từ và phanh điện từ. (3) Hao phí truyền tải tỉ lệ nghịch với bình phương điện trở dây. (4) Lõi máy biến áp ghép từ một khối thép đặc để tăng dòng điện Foucault.
\shortans{2}
\loigiai{
Đối chiếu từng phát biểu với kiến thức: Máy biến áp lí tưởng có U2/U1 = N2/N1 và U1I1 = U2I2. Truyền tải cùng công suất ở điện áp cao làm giảm hao phí I^2R. Có 2 phát biểu đúng.
}
\end{ex}

% ===== Câu 276 =====
% ID: L12C3B20-65-TN
% Mức: TH
\dangbt{Tổng hợp đường sức từ và xác định chiều từ trường}
\begin{ex}
% ID: L12C3B20-65-TN
% Mức: TH
Cơ sở Vật lí nào sau đây cần được sử dụng trong dạng tổng hợp đường sức từ và xác định chiều từ trường?
\choice
{\True Bên ngoài nam châm, đường sức từ có chiều đi ra từ cực Bắc và đi vào cực Nam.}
{Các đường sức từ có thể cắt nhau tại một điểm.}
{Đường sức từ luôn là những đường thẳng song song.}
{Đổi chiều dòng điện không làm đổi chiều từ trường.}
\loigiai{
Bên ngoài nam châm, đường sức từ có chiều đi ra từ cực Bắc và đi vào cực Nam. Do đó chọn A. Đường sức từ không cắt nhau; chiều từ trường được xác định bằng kim nam châm hoặc quy tắc nắm tay phải. Bên ngoài nam châm, đường sức đi từ cực Bắc đến cực Nam.
}
\end{ex}

% ===== Câu 277 =====
% ID: L12C3B20-66-TL
% Mức: VDC
\dangbt{Tổng hợp máy biến áp, truyền tải điện năng và ứng dụng cảm ứng điện từ}
\begin{bt}
% ID: L12C3B20-66-TL
% Mức: VDC
So sánh một nhận định đúng và một nhận định sai trong dạng tổng hợp máy biến áp, truyền tải điện năng và ứng dụng cảm ứng điện từ; giải thích căn cứ Vật lí.
\loigiai{
Nhận định đúng: Máy biến áp hoạt động với dòng điện xoay chiều dựa trên cảm ứng điện từ. Nhận định sai: Máy biến áp làm việc bình thường với nguồn một chiều không đổi. Căn cứ: Máy biến áp lí tưởng có U2/U1 = N2/N1 và U1I1 = U2I2. Truyền tải cùng công suất ở điện áp cao làm giảm hao phí I^2R.
}
\end{bt}

% ===== Câu 278 =====
% ID: L12C3B20-66-TLN
% Mức: VDC
\begin{ex}
% ID: L12C3B20-66-TLN
% Mức: VDC
Trong bốn phát biểu sau về tổng hợp máy biến áp, truyền tải điện năng và ứng dụng cảm ứng điện từ, có bao nhiêu phát biểu đúng? Chỉ ghi một số nguyên. (1) Máy biến áp hoạt động với dòng điện xoay chiều dựa trên cảm ứng điện từ. (2) Tăng điện áp truyền tải với cùng công suất làm giảm hao phí trên đường dây. (3) Máy tăng áp luôn có số vòng cuộn thứ cấp nhỏ hơn cuộn sơ cấp. (4) Lõi máy biến áp ghép từ một khối thép đặc để tăng dòng điện Foucault.
\shortans{2}
\loigiai{
Đối chiếu từng phát biểu với kiến thức: Máy biến áp lí tưởng có U2/U1 = N2/N1 và U1I1 = U2I2. Truyền tải cùng công suất ở điện áp cao làm giảm hao phí I^2R. Có 2 phát biểu đúng.
}
\end{ex}

% ===== Câu 279 =====
% ID: L12C3B20-66-TN
% Mức: TH
\dangbt{Tổng hợp đường sức từ và xác định chiều từ trường}
\begin{ex}
% ID: L12C3B20-66-TN
% Mức: TH
Phát biểu nào sau đây đúng về tổng hợp đường sức từ và xác định chiều từ trường?
\choice
{Đường sức từ luôn là những đường thẳng song song.}
{\True Đường sức từ của dòng điện thẳng là các đường tròn đồng tâm nằm trong mặt phẳng vuông góc với dây dẫn.}
{Đổi chiều dòng điện không làm đổi chiều từ trường.}
{Từ trường đều có các đường sức cong và mật độ thay đổi.}
\loigiai{
Đường sức từ của dòng điện thẳng là các đường tròn đồng tâm nằm trong mặt phẳng vuông góc với dây dẫn. Do đó chọn B. Đường sức từ không cắt nhau; chiều từ trường được xác định bằng kim nam châm hoặc quy tắc nắm tay phải. Bên ngoài nam châm, đường sức đi từ cực Bắc đến cực Nam.
}
\end{ex}

% ===== Câu 280 =====
% ID: L12C3B20-67-TL
% Mức: TH
\dangbt{Bài toán vận dụng cao kết hợp nhiều nội dung của Chương III}
\begin{bt}
% ID: L12C3B20-67-TL
% Mức: TH
Trình bày kiến thức cốt lõi và các bước giải bài thuộc dạng bài toán vận dụng cao kết hợp nhiều nội dung của chương iii.
\loigiai{
Bài toán tổng hợp phải tách các giai đoạn, xác định đúng phương chiều, chọn định luật phù hợp, đổi đơn vị thống nhất và kiểm tra ý nghĩa vật lí của kết quả. Các bước: tóm tắt dữ kiện; xác định phương, chiều và điều kiện; chọn hệ thức; đổi đơn vị; tính toán; kiểm tra kết quả.
}
\end{bt}

% ===== Câu 281 =====
% ID: L12C3B20-67-TLN
% Mức: TH
\begin{ex}
% ID: L12C3B20-67-TLN
% Mức: TH
Trong bốn phát biểu sau về bài toán vận dụng cao kết hợp nhiều nội dung của chương iii, có bao nhiêu phát biểu đúng? Chỉ ghi một số nguyên. (1) Khi giải bài tổng hợp cần xác định rõ chiều của $B$, dòng điện, lực từ và dòng điện cảm ứng. (2) Có thể cộng đại số mọi vectơ mà không xét phương chiều. (3) Các đại lượng phải được đổi về cùng hệ đơn vị trước khi thay số. (4) Có thể dùng trực tiếp số đo góc theo độ trong mọi biểu thức máy tính mà không kiểm tra chế độ.
\shortans{2}
\loigiai{
Đối chiếu từng phát biểu với kiến thức: Bài toán tổng hợp phải tách các giai đoạn, xác định đúng phương chiều, chọn định luật phù hợp, đổi đơn vị thống nhất và kiểm tra ý nghĩa vật lí của kết quả. Có 2 phát biểu đúng.
}
\end{ex}

% ===== Câu 282 =====
% ID: L12C3B20-67-TN
% Mức: TH
\dangbt{Tổng hợp đường sức từ và xác định chiều từ trường}
\begin{ex}
% ID: L12C3B20-67-TN
% Mức: TH
Trong Chương III, kết luận nào đúng khi giải bài thuộc dạng tổng hợp đường sức từ và xác định chiều từ trường?
\choice
{Đổi chiều dòng điện không làm đổi chiều từ trường.}
{Từ trường đều có các đường sức cong và mật độ thay đổi.}
{\True Chiều từ trường của dòng điện được xác định bằng quy tắc nắm tay phải.}
{Các đường sức từ có thể cắt nhau tại một điểm.}
\loigiai{
Chiều từ trường của dòng điện được xác định bằng quy tắc nắm tay phải. Do đó chọn C. Đường sức từ không cắt nhau; chiều từ trường được xác định bằng kim nam châm hoặc quy tắc nắm tay phải. Bên ngoài nam châm, đường sức đi từ cực Bắc đến cực Nam.
}
\end{ex}

% ===== Câu 283 =====
% ID: L12C3B20-68-TL
% Mức: TH
\dangbt{Bài toán vận dụng cao kết hợp nhiều nội dung của Chương III}
\begin{bt}
% ID: L12C3B20-68-TL
% Mức: TH
Phân tích các đại lượng, phương chiều và điều kiện áp dụng trong dạng bài toán vận dụng cao kết hợp nhiều nội dung của chương iii.
\loigiai{
Bài toán tổng hợp phải tách các giai đoạn, xác định đúng phương chiều, chọn định luật phù hợp, đổi đơn vị thống nhất và kiểm tra ý nghĩa vật lí của kết quả. Cần xác định rõ đại lượng vô hướng hay vectơ, quy ước chiều và điều kiện để công thức được áp dụng.
}
\end{bt}

% ===== Câu 284 =====
% ID: L12C3B20-68-TLN
% Mức: TH
\begin{ex}
% ID: L12C3B20-68-TLN
% Mức: TH
Trong bốn phát biểu sau về bài toán vận dụng cao kết hợp nhiều nội dung của chương iii, có bao nhiêu phát biểu đúng? Chỉ ghi một số nguyên. (1) Định luật Lenz quyết định chiều, còn định luật Faraday xác định độ lớn suất điện động cảm ứng. (2) Kết quả phải được kiểm tra về đơn vị và ý nghĩa vật lí. (3) Hiệu suất của mọi thiết bị điện luôn bằng 100 phần trăm. (4) Trong bài tổng hợp có thể bỏ qua điều kiện áp dụng của từng định luật.
\shortans{2}
\loigiai{
Đối chiếu từng phát biểu với kiến thức: Bài toán tổng hợp phải tách các giai đoạn, xác định đúng phương chiều, chọn định luật phù hợp, đổi đơn vị thống nhất và kiểm tra ý nghĩa vật lí của kết quả. Có 2 phát biểu đúng.
}
\end{ex}

% ===== Câu 285 =====
% ID: L12C3B20-68-TN
% Mức: VD
\dangbt{Tổng hợp đường sức từ và xác định chiều từ trường}
\begin{ex}
% ID: L12C3B20-68-TN
% Mức: VD
Chọn nhận định phù hợp với kiến thức Vật lí về tổng hợp đường sức từ và xác định chiều từ trường.
\choice
{Từ trường đều có các đường sức cong và mật độ thay đổi.}
{Các đường sức từ có thể cắt nhau tại một điểm.}
{Đường sức từ luôn là những đường thẳng song song.}
{\True Tại một điểm, hướng của vectơ cảm ứng từ trùng với hướng mà cực Bắc của kim nam châm chỉ.}
\loigiai{
Tại một điểm, hướng của vectơ cảm ứng từ trùng với hướng mà cực Bắc của kim nam châm chỉ. Do đó chọn D. Đường sức từ không cắt nhau; chiều từ trường được xác định bằng kim nam châm hoặc quy tắc nắm tay phải. Bên ngoài nam châm, đường sức đi từ cực Bắc đến cực Nam.
}
\end{ex}

% ===== Câu 286 =====
% ID: L12C3B20-69-TL
% Mức: VD
\dangbt{Bài toán vận dụng cao kết hợp nhiều nội dung của Chương III}
\begin{bt}
% ID: L12C3B20-69-TL
% Mức: VD
Nêu một sai lầm thường gặp khi giải bài toán vận dụng cao kết hợp nhiều nội dung của chương iii và cách khắc phục.
\loigiai{
Bài toán tổng hợp phải tách các giai đoạn, xác định đúng phương chiều, chọn định luật phù hợp, đổi đơn vị thống nhất và kiểm tra ý nghĩa vật lí của kết quả. Sai lầm thường gặp là bỏ qua phương chiều, góc hoặc đơn vị. Khắc phục bằng hình vẽ, quy ước dấu và đổi đơn vị trước khi tính.
}
\end{bt}

% ===== Câu 287 =====
% ID: L12C3B20-69-TLN
% Mức: VD
\begin{ex}
% ID: L12C3B20-69-TLN
% Mức: VD
Trong bốn phát biểu sau về bài toán vận dụng cao kết hợp nhiều nội dung của chương iii, có bao nhiêu phát biểu đúng? Chỉ ghi một số nguyên. (1) Khi giải bài tổng hợp cần xác định rõ chiều của $B$, dòng điện, lực từ và dòng điện cảm ứng. (2) Định luật Lenz quyết định chiều, còn định luật Faraday xác định độ lớn suất điện động cảm ứng. (3) Có thể dùng trực tiếp số đo góc theo độ trong mọi biểu thức máy tính mà không kiểm tra chế độ. (4) Trong bài tổng hợp có thể bỏ qua điều kiện áp dụng của từng định luật.
\shortans{2}
\loigiai{
Đối chiếu từng phát biểu với kiến thức: Bài toán tổng hợp phải tách các giai đoạn, xác định đúng phương chiều, chọn định luật phù hợp, đổi đơn vị thống nhất và kiểm tra ý nghĩa vật lí của kết quả. Có 2 phát biểu đúng.
}
\end{ex}

% ===== Câu 288 =====
% ID: L12C3B20-69-TN
% Mức: VD
\dangbt{Tổng hợp đường sức từ và xác định chiều từ trường}
\begin{ex}
% ID: L12C3B20-69-TN
% Mức: VD
Khi phân tích tổng hợp đường sức từ và xác định chiều từ trường, phát biểu nào đúng?
\choice
{\True Bên ngoài nam châm, đường sức từ có chiều đi ra từ cực Bắc và đi vào cực Nam.}
{Các đường sức từ có thể cắt nhau tại một điểm.}
{Đường sức từ luôn là những đường thẳng song song.}
{Đổi chiều dòng điện không làm đổi chiều từ trường.}
\loigiai{
Bên ngoài nam châm, đường sức từ có chiều đi ra từ cực Bắc và đi vào cực Nam. Do đó chọn A. Đường sức từ không cắt nhau; chiều từ trường được xác định bằng kim nam châm hoặc quy tắc nắm tay phải. Bên ngoài nam châm, đường sức đi từ cực Bắc đến cực Nam.
}
\end{ex}

% ===== Câu 289 =====
% ID: L12C3B20-70-TL
% Mức: VD
\dangbt{Bài toán vận dụng cao kết hợp nhiều nội dung của Chương III}
\begin{bt}
% ID: L12C3B20-70-TL
% Mức: VD
Đề xuất quy trình kiểm tra kết quả của một bài toán bài toán vận dụng cao kết hợp nhiều nội dung của chương iii.
\loigiai{
Bài toán tổng hợp phải tách các giai đoạn, xác định đúng phương chiều, chọn định luật phù hợp, đổi đơn vị thống nhất và kiểm tra ý nghĩa vật lí của kết quả. Kiểm tra thứ nguyên, dấu, giới hạn đặc biệt và sự phù hợp với ý nghĩa vật lí.
}
\end{bt}

% ===== Câu 290 =====
% ID: L12C3B20-70-TLN
% Mức: VD
\begin{ex}
% ID: L12C3B20-70-TLN
% Mức: VD
Trong bốn phát biểu sau về bài toán vận dụng cao kết hợp nhiều nội dung của chương iii, có bao nhiêu phát biểu đúng? Chỉ ghi một số nguyên. (1) Khi giải bài tổng hợp cần xác định rõ chiều của $B$, dòng điện, lực từ và dòng điện cảm ứng. (2) Có thể cộng đại số mọi vectơ mà không xét phương chiều. (3) Các đại lượng phải được đổi về cùng hệ đơn vị trước khi thay số. (4) Có thể dùng trực tiếp số đo góc theo độ trong mọi biểu thức máy tính mà không kiểm tra chế độ.
\shortans{2}
\loigiai{
Đối chiếu từng phát biểu với kiến thức: Bài toán tổng hợp phải tách các giai đoạn, xác định đúng phương chiều, chọn định luật phù hợp, đổi đơn vị thống nhất và kiểm tra ý nghĩa vật lí của kết quả. Có 2 phát biểu đúng.
}
\end{ex}

% ===== Câu 291 =====
% ID: L12C3B20-70-TN
% Mức: VD
\dangbt{Tổng hợp đường sức từ và xác định chiều từ trường}
\begin{ex}
% ID: L12C3B20-70-TN
% Mức: VD
Cơ sở Vật lí nào sau đây cần được sử dụng trong dạng tổng hợp đường sức từ và xác định chiều từ trường?
\choice
{Đường sức từ luôn là những đường thẳng song song.}
{\True Đường sức từ của dòng điện thẳng là các đường tròn đồng tâm nằm trong mặt phẳng vuông góc với dây dẫn.}
{Đổi chiều dòng điện không làm đổi chiều từ trường.}
{Từ trường đều có các đường sức cong và mật độ thay đổi.}
\loigiai{
Đường sức từ của dòng điện thẳng là các đường tròn đồng tâm nằm trong mặt phẳng vuông góc với dây dẫn. Do đó chọn B. Đường sức từ không cắt nhau; chiều từ trường được xác định bằng kim nam châm hoặc quy tắc nắm tay phải. Bên ngoài nam châm, đường sức đi từ cực Bắc đến cực Nam.
}
\end{ex}

% ===== Câu 292 =====
% ID: L12C3B20-71-TL
% Mức: VD
\dangbt{Bài toán vận dụng cao kết hợp nhiều nội dung của Chương III}
\begin{bt}
% ID: L12C3B20-71-TL
% Mức: VD
Xây dựng một tình huống thực tiễn thuộc dạng bài toán vận dụng cao kết hợp nhiều nội dung của chương iii và chỉ ra định luật cần dùng.
\loigiai{
Bài toán tổng hợp phải tách các giai đoạn, xác định đúng phương chiều, chọn định luật phù hợp, đổi đơn vị thống nhất và kiểm tra ý nghĩa vật lí của kết quả. Có thể chọn thiết bị hoặc hiện tượng thực tế tương ứng, nêu dữ kiện và chỉ rõ định luật dùng để giải thích hoặc tính toán.
}
\end{bt}

% ===== Câu 293 =====
% ID: L12C3B20-71-TLN
% Mức: VD
\begin{ex}
% ID: L12C3B20-71-TLN
% Mức: VD
Trong bốn phát biểu sau về bài toán vận dụng cao kết hợp nhiều nội dung của chương iii, có bao nhiêu phát biểu đúng? Chỉ ghi một số nguyên. (1) Định luật Lenz quyết định chiều, còn định luật Faraday xác định độ lớn suất điện động cảm ứng. (2) Kết quả phải được kiểm tra về đơn vị và ý nghĩa vật lí. (3) Hiệu suất của mọi thiết bị điện luôn bằng 100 phần trăm. (4) Trong bài tổng hợp có thể bỏ qua điều kiện áp dụng của từng định luật.
\shortans{2}
\loigiai{
Đối chiếu từng phát biểu với kiến thức: Bài toán tổng hợp phải tách các giai đoạn, xác định đúng phương chiều, chọn định luật phù hợp, đổi đơn vị thống nhất và kiểm tra ý nghĩa vật lí của kết quả. Có 2 phát biểu đúng.
}
\end{ex}

% ===== Câu 294 =====
% ID: L12C3B20-71-TN
% Mức: NB
\dangbt{Tổng hợp lực từ tác dụng lên dây dẫn mang dòng điện}
\begin{ex}
% ID: L12C3B20-71-TN
% Mức: NB
Phát biểu nào sau đây đúng về tổng hợp lực từ tác dụng lên dây dẫn mang dòng điện?
\choice
{\True Độ lớn lực từ tác dụng lên đoạn dây thẳng là $F$ = BIl sin(alpha).}
{Lực từ luôn cùng chiều với dòng điện.}
{Khi dây dẫn vuông góc với từ trường thì lực từ bằng 0.}
{Độ lớn lực từ không phụ thuộc chiều dài đoạn dây.}
\loigiai{
Độ lớn lực từ tác dụng lên đoạn dây thẳng là $F$ = BIl sin(alpha). Do đó chọn A. Với đoạn dây thẳng dài l mang dòng $I$ trong từ trường đều $B$, $F$ = BIl sin(alpha); chiều lực xác định bằng quy tắc bàn tay trái.
}
\end{ex}

% ===== Câu 295 =====
% ID: L12C3B20-72-TL
% Mức: VDC
\dangbt{Bài toán vận dụng cao kết hợp nhiều nội dung của Chương III}
\begin{bt}
% ID: L12C3B20-72-TL
% Mức: VDC
So sánh một nhận định đúng và một nhận định sai trong dạng bài toán vận dụng cao kết hợp nhiều nội dung của chương iii; giải thích căn cứ Vật lí.
\loigiai{
Nhận định đúng: Khi giải bài tổng hợp cần xác định rõ chiều của $B$, dòng điện, lực từ và dòng điện cảm ứng. Nhận định sai: Có thể cộng đại số mọi vectơ mà không xét phương chiều. Căn cứ: Bài toán tổng hợp phải tách các giai đoạn, xác định đúng phương chiều, chọn định luật phù hợp, đổi đơn vị thống nhất và kiểm tra ý nghĩa vật lí của kết quả.
}
\end{bt}

% ===== Câu 296 =====
% ID: L12C3B20-72-TLN
% Mức: VDC
\begin{ex}
% ID: L12C3B20-72-TLN
% Mức: VDC
Trong bốn phát biểu sau về bài toán vận dụng cao kết hợp nhiều nội dung của chương iii, có bao nhiêu phát biểu đúng? Chỉ ghi một số nguyên. (1) Khi giải bài tổng hợp cần xác định rõ chiều của $B$, dòng điện, lực từ và dòng điện cảm ứng. (2) Định luật Lenz quyết định chiều, còn định luật Faraday xác định độ lớn suất điện động cảm ứng. (3) Có thể dùng trực tiếp số đo góc theo độ trong mọi biểu thức máy tính mà không kiểm tra chế độ. (4) Trong bài tổng hợp có thể bỏ qua điều kiện áp dụng của từng định luật.
\shortans{2}
\loigiai{
Đối chiếu từng phát biểu với kiến thức: Bài toán tổng hợp phải tách các giai đoạn, xác định đúng phương chiều, chọn định luật phù hợp, đổi đơn vị thống nhất và kiểm tra ý nghĩa vật lí của kết quả. Có 2 phát biểu đúng.
}
\end{ex}

% ===== Câu 297 =====
% ID: L12C3B20-72-TN
% Mức: NB
\dangbt{Tổng hợp lực từ tác dụng lên dây dẫn mang dòng điện}
\begin{ex}
% ID: L12C3B20-72-TN
% Mức: NB
Trong Chương III, kết luận nào đúng khi giải bài thuộc dạng tổng hợp lực từ tác dụng lên dây dẫn mang dòng điện?
\choice
{Khi dây dẫn vuông góc với từ trường thì lực từ bằng 0.}
{\True Lực từ vuông góc với cả dòng điện và vectơ cảm ứng từ.}
{Độ lớn lực từ không phụ thuộc chiều dài đoạn dây.}
{Đổi chiều cảm ứng từ nhưng giữ dòng điện thì chiều lực từ không đổi.}
\loigiai{
Lực từ vuông góc với cả dòng điện và vectơ cảm ứng từ. Do đó chọn B. Với đoạn dây thẳng dài l mang dòng $I$ trong từ trường đều $B$, $F$ = BIl sin(alpha); chiều lực xác định bằng quy tắc bàn tay trái.
}
\end{ex}

% ===== Câu 298 =====
% ID: L12C3B20-73-TN
% Mức: NB
\begin{ex}
% ID: L12C3B20-73-TN
% Mức: NB
Chọn nhận định phù hợp với kiến thức Vật lí về tổng hợp lực từ tác dụng lên dây dẫn mang dòng điện.
\choice
{Độ lớn lực từ không phụ thuộc chiều dài đoạn dây.}
{Đổi chiều cảm ứng từ nhưng giữ dòng điện thì chiều lực từ không đổi.}
{\True Khi dây dẫn song song với đường sức từ thì lực từ bằng 0.}
{Lực từ luôn cùng chiều với dòng điện.}
\loigiai{
Khi dây dẫn song song với đường sức từ thì lực từ bằng 0. Do đó chọn C. Với đoạn dây thẳng dài l mang dòng $I$ trong từ trường đều $B$, $F$ = BIl sin(alpha); chiều lực xác định bằng quy tắc bàn tay trái.
}
\end{ex}

% ===== Câu 299 =====
% ID: L12C3B20-74-TN
% Mức: TH
\begin{ex}
% ID: L12C3B20-74-TN
% Mức: TH
Khi phân tích tổng hợp lực từ tác dụng lên dây dẫn mang dòng điện, phát biểu nào đúng?
\choice
{Đổi chiều cảm ứng từ nhưng giữ dòng điện thì chiều lực từ không đổi.}
{Lực từ luôn cùng chiều với dòng điện.}
{Khi dây dẫn vuông góc với từ trường thì lực từ bằng 0.}
{\True Đổi chiều dòng điện thì chiều lực từ đổi ngược.}
\loigiai{
Đổi chiều dòng điện thì chiều lực từ đổi ngược. Do đó chọn D. Với đoạn dây thẳng dài l mang dòng $I$ trong từ trường đều $B$, $F$ = BIl sin(alpha); chiều lực xác định bằng quy tắc bàn tay trái.
}
\end{ex}

% ===== Câu 300 =====
% ID: L12C3B20-75-TN
% Mức: TH
\begin{ex}
% ID: L12C3B20-75-TN
% Mức: TH
Cơ sở Vật lí nào sau đây cần được sử dụng trong dạng tổng hợp lực từ tác dụng lên dây dẫn mang dòng điện?
\choice
{\True Độ lớn lực từ tác dụng lên đoạn dây thẳng là $F$ = BIl sin(alpha).}
{Lực từ luôn cùng chiều với dòng điện.}
{Khi dây dẫn vuông góc với từ trường thì lực từ bằng 0.}
{Độ lớn lực từ không phụ thuộc chiều dài đoạn dây.}
\loigiai{
Độ lớn lực từ tác dụng lên đoạn dây thẳng là $F$ = BIl sin(alpha). Do đó chọn A. Với đoạn dây thẳng dài l mang dòng $I$ trong từ trường đều $B$, $F$ = BIl sin(alpha); chiều lực xác định bằng quy tắc bàn tay trái.
}
\end{ex}

% ===== Câu 301 =====
% ID: L12C3B20-76-TN
% Mức: TH
\begin{ex}
% ID: L12C3B20-76-TN
% Mức: TH
Phát biểu nào sau đây đúng về tổng hợp lực từ tác dụng lên dây dẫn mang dòng điện?
\choice
{Khi dây dẫn vuông góc với từ trường thì lực từ bằng 0.}
{\True Lực từ vuông góc với cả dòng điện và vectơ cảm ứng từ.}
{Độ lớn lực từ không phụ thuộc chiều dài đoạn dây.}
{Đổi chiều cảm ứng từ nhưng giữ dòng điện thì chiều lực từ không đổi.}
\loigiai{
Lực từ vuông góc với cả dòng điện và vectơ cảm ứng từ. Do đó chọn B. Với đoạn dây thẳng dài l mang dòng $I$ trong từ trường đều $B$, $F$ = BIl sin(alpha); chiều lực xác định bằng quy tắc bàn tay trái.
}
\end{ex}

% ===== Câu 302 =====
% ID: L12C3B20-77-TN
% Mức: TH
\begin{ex}
% ID: L12C3B20-77-TN
% Mức: TH
Trong Chương III, kết luận nào đúng khi giải bài thuộc dạng tổng hợp lực từ tác dụng lên dây dẫn mang dòng điện?
\choice
{Độ lớn lực từ không phụ thuộc chiều dài đoạn dây.}
{Đổi chiều cảm ứng từ nhưng giữ dòng điện thì chiều lực từ không đổi.}
{\True Khi dây dẫn song song với đường sức từ thì lực từ bằng 0.}
{Lực từ luôn cùng chiều với dòng điện.}
\loigiai{
Khi dây dẫn song song với đường sức từ thì lực từ bằng 0. Do đó chọn C. Với đoạn dây thẳng dài l mang dòng $I$ trong từ trường đều $B$, $F$ = BIl sin(alpha); chiều lực xác định bằng quy tắc bàn tay trái.
}
\end{ex}

% ===== Câu 303 =====
% ID: L12C3B20-78-TN
% Mức: VD
\begin{ex}
% ID: L12C3B20-78-TN
% Mức: VD
Chọn nhận định phù hợp với kiến thức Vật lí về tổng hợp lực từ tác dụng lên dây dẫn mang dòng điện.
\choice
{Đổi chiều cảm ứng từ nhưng giữ dòng điện thì chiều lực từ không đổi.}
{Lực từ luôn cùng chiều với dòng điện.}
{Khi dây dẫn vuông góc với từ trường thì lực từ bằng 0.}
{\True Đổi chiều dòng điện thì chiều lực từ đổi ngược.}
\loigiai{
Đổi chiều dòng điện thì chiều lực từ đổi ngược. Do đó chọn D. Với đoạn dây thẳng dài l mang dòng $I$ trong từ trường đều $B$, $F$ = BIl sin(alpha); chiều lực xác định bằng quy tắc bàn tay trái.
}
\end{ex}

% ===== Câu 304 =====
% ID: L12C3B20-79-TN
% Mức: VD
\begin{ex}
% ID: L12C3B20-79-TN
% Mức: VD
Khi phân tích tổng hợp lực từ tác dụng lên dây dẫn mang dòng điện, phát biểu nào đúng?
\choice
{\True Độ lớn lực từ tác dụng lên đoạn dây thẳng là $F$ = BIl sin(alpha).}
{Lực từ luôn cùng chiều với dòng điện.}
{Khi dây dẫn vuông góc với từ trường thì lực từ bằng 0.}
{Độ lớn lực từ không phụ thuộc chiều dài đoạn dây.}
\loigiai{
Độ lớn lực từ tác dụng lên đoạn dây thẳng là $F$ = BIl sin(alpha). Do đó chọn A. Với đoạn dây thẳng dài l mang dòng $I$ trong từ trường đều $B$, $F$ = BIl sin(alpha); chiều lực xác định bằng quy tắc bàn tay trái.
}
\end{ex}

% ===== Câu 305 =====
% ID: L12C3B20-80-TN
% Mức: VD
\begin{ex}
% ID: L12C3B20-80-TN
% Mức: VD
Cơ sở Vật lí nào sau đây cần được sử dụng trong dạng tổng hợp lực từ tác dụng lên dây dẫn mang dòng điện?
\choice
{Khi dây dẫn vuông góc với từ trường thì lực từ bằng 0.}
{\True Lực từ vuông góc với cả dòng điện và vectơ cảm ứng từ.}
{Độ lớn lực từ không phụ thuộc chiều dài đoạn dây.}
{Đổi chiều cảm ứng từ nhưng giữ dòng điện thì chiều lực từ không đổi.}
\loigiai{
Lực từ vuông góc với cả dòng điện và vectơ cảm ứng từ. Do đó chọn B. Với đoạn dây thẳng dài l mang dòng $I$ trong từ trường đều $B$, $F$ = BIl sin(alpha); chiều lực xác định bằng quy tắc bàn tay trái.
}
\end{ex}

% ===== Câu 306 =====
% ID: L12C3B20-81-TN
% Mức: NB
\dangbt{Tổng hợp từ thông, định luật Lenz và suất điện động cảm ứng}
\begin{ex}
% ID: L12C3B20-81-TN
% Mức: NB
Phát biểu nào sau đây đúng về tổng hợp từ thông, định luật lenz và suất điện động cảm ứng?
\choice
{\True Từ thông qua khung phẳng là Phi = BS cos(alpha), với alpha là góc giữa $B$ và pháp tuyến.}
{Từ thông là đại lượng vectơ.}
{Dòng điện cảm ứng luôn làm tăng từ thông ban đầu.}
{Suất điện động cảm ứng chỉ phụ thuộc điện trở của mạch.}
\loigiai{
Từ thông qua khung phẳng là Phi = BS cos(alpha), với alpha là góc giữa $B$ và pháp tuyến. Do đó chọn A. Phi = BS cos(alpha). Suất điện động cảm ứng có độ lớn |e| = |Delta Phi/Delta t|; chiều dòng điện cảm ứng tuân theo định luật Lenz.
}
\end{ex}

% ===== Câu 307 =====
% ID: L12C3B20-82-TN
% Mức: NB
\begin{ex}
% ID: L12C3B20-82-TN
% Mức: NB
Trong Chương III, kết luận nào đúng khi giải bài thuộc dạng tổng hợp từ thông, định luật lenz và suất điện động cảm ứng?
\choice
{Dòng điện cảm ứng luôn làm tăng từ thông ban đầu.}
{\True Dòng điện cảm ứng có chiều chống lại sự biến thiên từ thông sinh ra nó.}
{Suất điện động cảm ứng chỉ phụ thuộc điện trở của mạch.}
{Khung dây đứng yên trong từ trường đều không đổi luôn có dòng điện cảm ứng.}
\loigiai{
Dòng điện cảm ứng có chiều chống lại sự biến thiên từ thông sinh ra nó. Do đó chọn B. Phi = BS cos(alpha). Suất điện động cảm ứng có độ lớn |e| = |Delta Phi/Delta t|; chiều dòng điện cảm ứng tuân theo định luật Lenz.
}
\end{ex}

% ===== Câu 308 =====
% ID: L12C3B20-83-TN
% Mức: NB
\begin{ex}
% ID: L12C3B20-83-TN
% Mức: NB
Chọn nhận định phù hợp với kiến thức Vật lí về tổng hợp từ thông, định luật lenz và suất điện động cảm ứng.
\choice
{Suất điện động cảm ứng chỉ phụ thuộc điện trở của mạch.}
{Khung dây đứng yên trong từ trường đều không đổi luôn có dòng điện cảm ứng.}
{\True Độ lớn suất điện động cảm ứng bằng độ lớn tốc độ biến thiên từ thông.}
{Từ thông là đại lượng vectơ.}
\loigiai{
Độ lớn suất điện động cảm ứng bằng độ lớn tốc độ biến thiên từ thông. Do đó chọn C. Phi = BS cos(alpha). Suất điện động cảm ứng có độ lớn |e| = |Delta Phi/Delta t|; chiều dòng điện cảm ứng tuân theo định luật Lenz.
}
\end{ex}

% ===== Câu 309 =====
% ID: L12C3B20-84-TN
% Mức: TH
\begin{ex}
% ID: L12C3B20-84-TN
% Mức: TH
Khi phân tích tổng hợp từ thông, định luật lenz và suất điện động cảm ứng, phát biểu nào đúng?
\choice
{Khung dây đứng yên trong từ trường đều không đổi luôn có dòng điện cảm ứng.}
{Từ thông là đại lượng vectơ.}
{Dòng điện cảm ứng luôn làm tăng từ thông ban đầu.}
{\True Từ thông không đổi thì không xuất hiện suất điện động cảm ứng.}
\loigiai{
Từ thông không đổi thì không xuất hiện suất điện động cảm ứng. Do đó chọn D. Phi = BS cos(alpha). Suất điện động cảm ứng có độ lớn |e| = |Delta Phi/Delta t|; chiều dòng điện cảm ứng tuân theo định luật Lenz.
}
\end{ex}

% ===== Câu 310 =====
% ID: L12C3B20-85-TN
% Mức: TH
\begin{ex}
% ID: L12C3B20-85-TN
% Mức: TH
Cơ sở Vật lí nào sau đây cần được sử dụng trong dạng tổng hợp từ thông, định luật lenz và suất điện động cảm ứng?
\choice
{\True Từ thông qua khung phẳng là Phi = BS cos(alpha), với alpha là góc giữa $B$ và pháp tuyến.}
{Từ thông là đại lượng vectơ.}
{Dòng điện cảm ứng luôn làm tăng từ thông ban đầu.}
{Suất điện động cảm ứng chỉ phụ thuộc điện trở của mạch.}
\loigiai{
Từ thông qua khung phẳng là Phi = BS cos(alpha), với alpha là góc giữa $B$ và pháp tuyến. Do đó chọn A. Phi = BS cos(alpha). Suất điện động cảm ứng có độ lớn |e| = |Delta Phi/Delta t|; chiều dòng điện cảm ứng tuân theo định luật Lenz.
}
\end{ex}

% ===== Câu 311 =====
% ID: L12C3B20-86-TN
% Mức: TH
\begin{ex}
% ID: L12C3B20-86-TN
% Mức: TH
Phát biểu nào sau đây đúng về tổng hợp từ thông, định luật lenz và suất điện động cảm ứng?
\choice
{Dòng điện cảm ứng luôn làm tăng từ thông ban đầu.}
{\True Dòng điện cảm ứng có chiều chống lại sự biến thiên từ thông sinh ra nó.}
{Suất điện động cảm ứng chỉ phụ thuộc điện trở của mạch.}
{Khung dây đứng yên trong từ trường đều không đổi luôn có dòng điện cảm ứng.}
\loigiai{
Dòng điện cảm ứng có chiều chống lại sự biến thiên từ thông sinh ra nó. Do đó chọn B. Phi = BS cos(alpha). Suất điện động cảm ứng có độ lớn |e| = |Delta Phi/Delta t|; chiều dòng điện cảm ứng tuân theo định luật Lenz.
}
\end{ex}

% ===== Câu 312 =====
% ID: L12C3B20-87-TN
% Mức: TH
\begin{ex}
% ID: L12C3B20-87-TN
% Mức: TH
Trong Chương III, kết luận nào đúng khi giải bài thuộc dạng tổng hợp từ thông, định luật lenz và suất điện động cảm ứng?
\choice
{Suất điện động cảm ứng chỉ phụ thuộc điện trở của mạch.}
{Khung dây đứng yên trong từ trường đều không đổi luôn có dòng điện cảm ứng.}
{\True Độ lớn suất điện động cảm ứng bằng độ lớn tốc độ biến thiên từ thông.}
{Từ thông là đại lượng vectơ.}
\loigiai{
Độ lớn suất điện động cảm ứng bằng độ lớn tốc độ biến thiên từ thông. Do đó chọn C. Phi = BS cos(alpha). Suất điện động cảm ứng có độ lớn |e| = |Delta Phi/Delta t|; chiều dòng điện cảm ứng tuân theo định luật Lenz.
}
\end{ex}

% ===== Câu 313 =====
% ID: L12C3B20-88-TN
% Mức: VD
\begin{ex}
% ID: L12C3B20-88-TN
% Mức: VD
Chọn nhận định phù hợp với kiến thức Vật lí về tổng hợp từ thông, định luật lenz và suất điện động cảm ứng.
\choice
{Khung dây đứng yên trong từ trường đều không đổi luôn có dòng điện cảm ứng.}
{Từ thông là đại lượng vectơ.}
{Dòng điện cảm ứng luôn làm tăng từ thông ban đầu.}
{\True Từ thông không đổi thì không xuất hiện suất điện động cảm ứng.}
\loigiai{
Từ thông không đổi thì không xuất hiện suất điện động cảm ứng. Do đó chọn D. Phi = BS cos(alpha). Suất điện động cảm ứng có độ lớn |e| = |Delta Phi/Delta t|; chiều dòng điện cảm ứng tuân theo định luật Lenz.
}
\end{ex}

% ===== Câu 314 =====
% ID: L12C3B20-89-TN
% Mức: VD
\begin{ex}
% ID: L12C3B20-89-TN
% Mức: VD
Khi phân tích tổng hợp từ thông, định luật lenz và suất điện động cảm ứng, phát biểu nào đúng?
\choice
{\True Từ thông qua khung phẳng là Phi = BS cos(alpha), với alpha là góc giữa $B$ và pháp tuyến.}
{Từ thông là đại lượng vectơ.}
{Dòng điện cảm ứng luôn làm tăng từ thông ban đầu.}
{Suất điện động cảm ứng chỉ phụ thuộc điện trở của mạch.}
\loigiai{
Từ thông qua khung phẳng là Phi = BS cos(alpha), với alpha là góc giữa $B$ và pháp tuyến. Do đó chọn A. Phi = BS cos(alpha). Suất điện động cảm ứng có độ lớn |e| = |Delta Phi/Delta t|; chiều dòng điện cảm ứng tuân theo định luật Lenz.
}
\end{ex}

% ===== Câu 315 =====
% ID: L12C3B20-90-TN
% Mức: VD
\begin{ex}
% ID: L12C3B20-90-TN
% Mức: VD
Cơ sở Vật lí nào sau đây cần được sử dụng trong dạng tổng hợp từ thông, định luật lenz và suất điện động cảm ứng?
\choice
{Dòng điện cảm ứng luôn làm tăng từ thông ban đầu.}
{\True Dòng điện cảm ứng có chiều chống lại sự biến thiên từ thông sinh ra nó.}
{Suất điện động cảm ứng chỉ phụ thuộc điện trở của mạch.}
{Khung dây đứng yên trong từ trường đều không đổi luôn có dòng điện cảm ứng.}
\loigiai{
Dòng điện cảm ứng có chiều chống lại sự biến thiên từ thông sinh ra nó. Do đó chọn B. Phi = BS cos(alpha). Suất điện động cảm ứng có độ lớn |e| = |Delta Phi/Delta t|; chiều dòng điện cảm ứng tuân theo định luật Lenz.
}
\end{ex}

% ===== Câu 316 =====
% ID: L12C3B20-91-TN
% Mức: NB
\dangbt{Tổng hợp máy phát điện xoay chiều và đồ thị điện từ}
\begin{ex}
% ID: L12C3B20-91-TN
% Mức: NB
Phát biểu nào sau đây đúng về tổng hợp máy phát điện xoay chiều và đồ thị điện từ?
\choice
{\True Máy phát điện xoay chiều hoạt động dựa trên hiện tượng cảm ứng điện từ.}
{Máy phát điện biến điện năng thành cơ năng.}
{Tăng tốc độ quay làm giảm biên độ suất điện động.}
{Từ thông và suất điện động luôn cùng pha.}
\loigiai{
Máy phát điện xoay chiều hoạt động dựa trên hiện tượng cảm ứng điện từ. Do đó chọn A. Máy phát điện biến cơ năng thành điện năng nhờ cảm ứng điện từ; nếu Phi = Phi0 cos(omega t + phi) thì e = E0 sin(omega t + phi), E0 = NBS omega.
}
\end{ex}

% ===== Câu 317 =====
% ID: L12C3B20-92-TN
% Mức: NB
\begin{ex}
% ID: L12C3B20-92-TN
% Mức: NB
Trong Chương III, kết luận nào đúng khi giải bài thuộc dạng tổng hợp máy phát điện xoay chiều và đồ thị điện từ?
\choice
{Tăng tốc độ quay làm giảm biên độ suất điện động.}
{\True Khung $N$ vòng quay đều có suất điện động cực đại E0 = NBS omega.}
{Từ thông và suất điện động luôn cùng pha.}
{Suất điện động cực đại không phụ thuộc số vòng dây.}
\loigiai{
Khung $N$ vòng quay đều có suất điện động cực đại E0 = NBS omega. Do đó chọn B. Máy phát điện biến cơ năng thành điện năng nhờ cảm ứng điện từ; nếu Phi = Phi0 cos(omega t + phi) thì e = E0 sin(omega t + phi), E0 = NBS omega.
}
\end{ex}

% ===== Câu 318 =====
% ID: L12C3B20-93-TN
% Mức: NB
\begin{ex}
% ID: L12C3B20-93-TN
% Mức: NB
Chọn nhận định phù hợp với kiến thức Vật lí về tổng hợp máy phát điện xoay chiều và đồ thị điện từ.
\choice
{Từ thông và suất điện động luôn cùng pha.}
{Suất điện động cực đại không phụ thuộc số vòng dây.}
{\True Tần số suất điện động bằng tần số quay của khung trong mô hình một cặp cực.}
{Máy phát điện biến điện năng thành cơ năng.}
\loigiai{
Tần số suất điện động bằng tần số quay của khung trong mô hình một cặp cực. Do đó chọn C. Máy phát điện biến cơ năng thành điện năng nhờ cảm ứng điện từ; nếu Phi = Phi0 cos(omega t + phi) thì e = E0 sin(omega t + phi), E0 = NBS omega.
}
\end{ex}

% ===== Câu 319 =====
% ID: L12C3B20-94-TN
% Mức: TH
\begin{ex}
% ID: L12C3B20-94-TN
% Mức: TH
Khi phân tích tổng hợp máy phát điện xoay chiều và đồ thị điện từ, phát biểu nào đúng?
\choice
{Suất điện động cực đại không phụ thuộc số vòng dây.}
{Máy phát điện biến điện năng thành cơ năng.}
{Tăng tốc độ quay làm giảm biên độ suất điện động.}
{\True Suất điện động cảm ứng lệch pha pi/2 so với từ thông.}
\loigiai{
Suất điện động cảm ứng lệch pha pi/2 so với từ thông. Do đó chọn D. Máy phát điện biến cơ năng thành điện năng nhờ cảm ứng điện từ; nếu Phi = Phi0 cos(omega t + phi) thì e = E0 sin(omega t + phi), E0 = NBS omega.
}
\end{ex}

% ===== Câu 320 =====
% ID: L12C3B20-95-TN
% Mức: TH
\begin{ex}
% ID: L12C3B20-95-TN
% Mức: TH
Cơ sở Vật lí nào sau đây cần được sử dụng trong dạng tổng hợp máy phát điện xoay chiều và đồ thị điện từ?
\choice
{\True Máy phát điện xoay chiều hoạt động dựa trên hiện tượng cảm ứng điện từ.}
{Máy phát điện biến điện năng thành cơ năng.}
{Tăng tốc độ quay làm giảm biên độ suất điện động.}
{Từ thông và suất điện động luôn cùng pha.}
\loigiai{
Máy phát điện xoay chiều hoạt động dựa trên hiện tượng cảm ứng điện từ. Do đó chọn A. Máy phát điện biến cơ năng thành điện năng nhờ cảm ứng điện từ; nếu Phi = Phi0 cos(omega t + phi) thì e = E0 sin(omega t + phi), E0 = NBS omega.
}
\end{ex}

% ===== Câu 321 =====
% ID: L12C3B20-96-TN
% Mức: TH
\begin{ex}
% ID: L12C3B20-96-TN
% Mức: TH
Phát biểu nào sau đây đúng về tổng hợp máy phát điện xoay chiều và đồ thị điện từ?
\choice
{Tăng tốc độ quay làm giảm biên độ suất điện động.}
{\True Khung $N$ vòng quay đều có suất điện động cực đại E0 = NBS omega.}
{Từ thông và suất điện động luôn cùng pha.}
{Suất điện động cực đại không phụ thuộc số vòng dây.}
\loigiai{
Khung $N$ vòng quay đều có suất điện động cực đại E0 = NBS omega. Do đó chọn B. Máy phát điện biến cơ năng thành điện năng nhờ cảm ứng điện từ; nếu Phi = Phi0 cos(omega t + phi) thì e = E0 sin(omega t + phi), E0 = NBS omega.
}
\end{ex}

% ===== Câu 322 =====
% ID: L12C3B20-97-TN
% Mức: TH
\begin{ex}
% ID: L12C3B20-97-TN
% Mức: TH
Trong Chương III, kết luận nào đúng khi giải bài thuộc dạng tổng hợp máy phát điện xoay chiều và đồ thị điện từ?
\choice
{Từ thông và suất điện động luôn cùng pha.}
{Suất điện động cực đại không phụ thuộc số vòng dây.}
{\True Tần số suất điện động bằng tần số quay của khung trong mô hình một cặp cực.}
{Máy phát điện biến điện năng thành cơ năng.}
\loigiai{
Tần số suất điện động bằng tần số quay của khung trong mô hình một cặp cực. Do đó chọn C. Máy phát điện biến cơ năng thành điện năng nhờ cảm ứng điện từ; nếu Phi = Phi0 cos(omega t + phi) thì e = E0 sin(omega t + phi), E0 = NBS omega.
}
\end{ex}

% ===== Câu 323 =====
% ID: L12C3B20-98-TN
% Mức: VD
\begin{ex}
% ID: L12C3B20-98-TN
% Mức: VD
Chọn nhận định phù hợp với kiến thức Vật lí về tổng hợp máy phát điện xoay chiều và đồ thị điện từ.
\choice
{Suất điện động cực đại không phụ thuộc số vòng dây.}
{Máy phát điện biến điện năng thành cơ năng.}
{Tăng tốc độ quay làm giảm biên độ suất điện động.}
{\True Suất điện động cảm ứng lệch pha pi/2 so với từ thông.}
\loigiai{
Suất điện động cảm ứng lệch pha pi/2 so với từ thông. Do đó chọn D. Máy phát điện biến cơ năng thành điện năng nhờ cảm ứng điện từ; nếu Phi = Phi0 cos(omega t + phi) thì e = E0 sin(omega t + phi), E0 = NBS omega.
}
\end{ex}

% ===== Câu 324 =====
% ID: L12C3B20-99-TN
% Mức: VD
\begin{ex}
% ID: L12C3B20-99-TN
% Mức: VD
Khi phân tích tổng hợp máy phát điện xoay chiều và đồ thị điện từ, phát biểu nào đúng?
\choice
{\True Máy phát điện xoay chiều hoạt động dựa trên hiện tượng cảm ứng điện từ.}
{Máy phát điện biến điện năng thành cơ năng.}
{Tăng tốc độ quay làm giảm biên độ suất điện động.}
{Từ thông và suất điện động luôn cùng pha.}
\loigiai{
Máy phát điện xoay chiều hoạt động dựa trên hiện tượng cảm ứng điện từ. Do đó chọn A. Máy phát điện biến cơ năng thành điện năng nhờ cảm ứng điện từ; nếu Phi = Phi0 cos(omega t + phi) thì e = E0 sin(omega t + phi), E0 = NBS omega.
}
\end{ex}
